\documentclass[a4paper, 12pt, italian]{report}

\usepackage[italian]{babel}


%%%%%%%% javadoc doclet http://doclet.github.io   %%%%%%%%%%%%%%%%%%%%%%

\usepackage{blindtext}
\usepackage{ifthen}
\usepackage{ifpdf}
\usepackage[headings]{fullpage}
\ifpdf \usepackage[pdftex, pdfpagemode={UseOutlines},bookmarks,colorlinks,linkcolor={blue},plainpages=false,pdfpagelabels,citecolor={red},breaklinks=true]{hyperref}
  \usepackage[pdftex]{graphicx}
  \pdfcompresslevel=9
  \DeclareGraphicsRule{*}{mps}{*}{}
\else
  \usepackage[dvips]{graphicx}
\fi

\newcommand{\entityintro}[3]{%
  \hbox to \hsize{%
    \vbox{%
      \hbox to .2in{}%
    }%
    {\bf  #1}%
    \dotfill\pageref{#2}%
  }
  \makebox[\hsize]{%
    \parbox{.4in}{}%
    \parbox[l]{5in}{%
      \vspace{1mm}%
      #3%
      \vspace{1mm}%
    }%
  }%
}
\newcommand{\refdefined}[1]{
\expandafter\ifx\csname r@#1\endcsname\relax
\relax\else
{$($in \ref{#1}, page \pageref{#1}$)$}\fi}
\date{null}
\chardef\textbackslash=`\\


\usepackage{fullpage} 
\usepackage{hyperref}
\hypersetup{
	colorlinks=true,allcolors=blue,pdfpagemode=UseOutlines,pdfpagelayout=TwoPageRight
}
\usepackage{hypcap}

%%%%%%%%%%%%%%%%%%%%%%%%%%%%%%%%%%%%%%%%%%%%%%%%%%%%

\usepackage{ctable}
\usepackage{url}
\usepackage[utf8]{inputenc}
%\usepackage{graphicx}
\usepackage{amsmath}
\usepackage{color}
\usepackage{listings}

%tabella
\usepackage{array}
\usepackage{xcolor,colortbl}
\usepackage{booktabs}
\usepackage{longtable}

\usepackage[raggedright]{titlesec}
\usepackage{blindtext}

\titleformat{\chapter}[hang]{\bfseries\huge}{\thechapter.}{2pc}{}
\titlelabel{\thetitle.\quad}   % For consistency in all headin

\definecolor{Gray}{gray}{0.8}
\newcolumntype{a}{>{\columncolor{Gray}}p{.25\textwidth}}

\begin{document}

% === Per gli snippet =====================================
\definecolor{mygreen}{rgb}{0,0.6,0}
\definecolor{mygray}{rgb}{0.5,0.5,0.5}
\definecolor{mymauve}{rgb}{0.58,0,0.82}
\lstset{
backgroundcolor=\color{white}, % choose the background color; you must add \usepackage{color} or \usepackage{xcolor}
basicstyle=\footnotesize\ttfamily, % the size of the fonts that are used for the code
breakatwhitespace=false, % sets if automatic breaks should only happen at whitespace
breaklines=true, % sets automatic line breaking
captionpos=n, % sets the caption-position to bottom
commentstyle=\color{mygreen}, % comment style
deletekeywords={}, % if you want to delete keywords from the given language
escapeinside={\%}{)}, % if you want to add LaTeX within your code
extendedchars=false, % lets you use non-ASCII characters; for 8-bits encodings only, does not work with UTF-8
frame=single, % adds a frame around the code
keepspaces=true, % keeps spaces in text, useful for keeping indentation of code (possibly needs columns=flexible)
keywordstyle=\color{blue}, % keyword style
language=Java, % the language of the code
morekeywords={}, % if you want to add more keywords to the set
numbers=left, % where to put the line-numbers; possible values are (none, left, right)
numbersep=5pt, % how far the line-numbers are from the code
numberstyle=\tiny\color{mygray}, % the style that is used for the line-numbers
rulecolor=\color{black}, % if not set, the frame-color may be changed on line-breaks within not-black text (e.g. comments (green here))
showspaces=false, % show spaces everywhere adding particular underscores; it overrides 'showstringspaces'
showstringspaces=false, % underline spaces within strings only
showtabs=false, % show tabs within strings adding particular underscores
stepnumber=1, % the step between two line-numbers. If it's 1, each line will be numbered
stringstyle=\color{mymauve}, % string literal style
tabsize=2, % sets default tabsize to 2 spaces
title=\lstname % show the filename of files included with \lstinputlisting; also try caption instead of title
literate=
  {á}{{\'a}}1 {é}{{\'e}}1 {í}{{\'i}}1 {ó}{{\'o}}1 {ú}{{\'u}}1
  {Á}{{\'A}}1 {É}{{\'E}}1 {Í}{{\'I}}1 {Ó}{{\'O}}1 {Ú}{{\'U}}1
  {à}{{\`a}}1 {è}{{\`e}}1 {ì}{{\`i}}1 {ò}{{\`o}}1 {ù}{{\`u}}1
  {À}{{\`A}}1 {È}{{\'E}}1 {Ì}{{\`I}}1 {Ò}{{\`O}}1 {Ù}{{\`U}}1
  {ä}{{\"a}}1 {ë}{{\"e}}1 {ï}{{\"i}}1 {ö}{{\"o}}1 {ü}{{\"u}}1
  {Ä}{{\"A}}1 {Ë}{{\"E}}1 {Ï}{{\"I}}1 {Ö}{{\"O}}1 {Ü}{{\"U}}1
  {â}{{\^a}}1 {ê}{{\^e}}1 {î}{{\^i}}1 {ô}{{\^o}}1 {û}{{\^u}}1
  {Â}{{\^A}}1 {Ê}{{\^E}}1 {Î}{{\^I}}1 {Ô}{{\^O}}1 {Û}{{\^U}}1
  {œ}{{\oe}}1 {Œ}{{\OE}}1 {æ}{{\ae}}1 {Æ}{{\AE}}1 {ß}{{\ss}}1
  {ç}{{\c c}}1 {Ç}{{\c C}}1 {ø}{{\o}}1 {å}{{\r a}}1 {Å}{{\r A}}1
  {€}{{\EUR}}1 {£}{{\pounds}}1
}
% === Fine ===================================================

\begin{titlepage}
\newcommand{\HRule}{\rule{\linewidth}{0.5mm}} 
\center 
\textsc{\LARGE Università degli studi di Salerno}\\[1cm] 
\includegraphics[width=3.5cm]{img/logo.jpg} \\[1cm]
\textsc{\large Progetto di Ingegneria del Software II}\\[0.5cm]
\textsc{\Large Repominer Evolution}\\[0.5cm] 
 \HRule \\[0.4cm]
{ \large \bfseries Object Design Document}\\[0.4cm] 
\HRule \\[1.5cm]

\begin{minipage}{0.4\textwidth}
\begin{flushleft} \large
\emph{Autori:}\\
Matteo \textsc{Merola}\\
Carlo \textsc{Branca}\\
Simone \textsc{Scalabrino}\\
Giovanni \textsc{Grano}\\
\end{flushleft}
\end{minipage}
~
\begin{minipage}{0.4\textwidth}
\begin{flushright} \large
\emph{Supervisore:} \\
Prof. Andrea \textsc{De Lucia}
\end{flushright}
\end{minipage}\\[2.5cm]

{\Large Object Design Document}\\
Versione 1.0\\[1cm]

{\large \today} % Date, change the \today to a set date if you want to be precise

\vfill

\end{titlepage}	
       \chapter*{Revision History}

\ctable[
    caption = Tabella delle revisioni del documento,
    label   = width,
    width   = 145mm,
    pos     = ht
 ]{ccc>{\raggedright}X}{}
 { 
\FL
\multicolumn{1}{c}{Data} &\multicolumn{1}{c}{Versione}&\multicolumn{1}{c}{Descrizione}&\multicolumn{1}{c}{Autori}\ML
 14/07/2014 &
 1.0 &
 Versione iniziale&
 Simone Scalabrino, Giovanni Grano, Carlo Branca, Matteo Merola 
 \LL
}
	\setcounter{tocdepth}{1}	
	\tableofcontents
	\chapter{Introduzione}
\section{Scopo del documento}
Lo scopo di questo documento è di fornire informazioni sulla gestione del progetto. In particolare, il periodo di riferimento è quello che va dalla data di rilascio del primo incremento (13/06/2014) alla data di completamento del secondo incremento previsto in fase di pianificazione (1/07/2014).

\section{Acronimi e abbreviazioni}
\begin{table}[ht]
\centering
\begin{tabular}{|c|c|}
 \hline
 \rowcolor{Gray}\textbf{Acronimo}			& \textbf{Descrizione}				\\
 \hline
 SPMP							& Software Project Management Plan		\\
 \hline
 ODD							& Object Design Document			\\
 \hline
 TP							& Test Plan					\\
 \hline
 TCS							& Test Case Specification			\\
 \hline
 TL							& Test Log					\\
 \hline
 TIR							& Test Incident Report				\\
 \hline
 TSR							& Test Summary Report				\\
 \hline
 BCC							& Basic Code Change model			\\
 \hline
 CB							& Carlo Branca					\\
 \hline
 GG							& Giovanni Grano				\\
 \hline
 MM							& Matteo Merola					\\
 \hline
 SS							& Simone Scalabrino				\\
 \hline
\end{tabular}
\end{table}
	\chapter{Linee guida per l'implementazione}

\section{Nomi di file}
Il software Java utilizza i seguenti suffissi per i file:
\begin{itemize}
  \item Per i sorgenti Java è: .java;
  \item Per i file Bytecode è .class.
\end{itemize}

\section{Organizzazione dei file}
Un file consiste di sezioni che dovrebbero essere separate da linee bianche e un commento opzionale che identifica ogni sezione. File più lunghi di 2000 linee sono ingombranti e devono essere evitati.

\subsection{File sorgenti}
Ogni file sorgente Java contiene una singola classe pubblica o un’interfaccia. Quando ci sono classi e interfacce private associate con la classe pubblica, è possibile inserirle nello stesso file sorgente della classe pubblica. La classe pubblica deve essere la prima classe o interfaccia nel file.

I file sorgenti Java hanno la seguente struttura:
\begin{itemize}
  \item \textsc{Commenti di inizio:} tutti i file sorgenti devono iniziare con un commento in stile C che elenca il nome della classe, descrizione, autore e informazioni sulla versione, informazioni di copyright.
  \begin{lstlisting}
    /**
    * Nome della classe
    * 
    * Descrizione
    *
    * Autore
    *
    * Informazione di versione
    *
    * 2014 - Copyright by University of Salerno
    */
  \end{lstlisting}
  
  \item \textsc{Istruzioni di package e import:} la prima linea non commento di molti file sorgenti Java è l'istruzione \textit{package} che può essere seguita da istruzioni \textit{import}. Ad esempio:\\
  \begin{lstlisting}
    package sie.miner.parser;
    
    import java.util.Map;
  \end{lstlisting}

  \item \textsc{Dichiarazioni di classe e di interfaccia:} l'ordine in cui le dichiarazioni di una classe o interfaccia devono apparire è il seguente: 
\begin{itemize}
  \item commento di documentazione della classe/interfaccia (/** ... */);
  \item istruzione \textit{class} o \textit{interface};
  \item commento di implementazione della classe/interfaccia, se necessario: questo commento deve contenere informazioni generali sulla classe o interfaccia che non sono appropriate per il commento di documentazione;
  \item variabili di classe (static): prima le variabili di classe public, poi quelle protected e infine quelle private;
  \item variabili di istanze: prima quelle public, poi quelle protected e infine quelle private;
  \item costruttori;
  \item metodi: questi devono essere raggruppati in base alla loro funzionalità piuttosto che in base a regole di visibilità o accessibilità. Ad esempio, un metodo di classe privato può stare tra due metodi pubblici. L'obiettivo è quello di rendere più semplice la lettura e la comprensione del codice.
\end{itemize}  

\end{itemize}

\section{Indentazione}
Come unità di indentazione devono essere usati quattro spazi ma la costruzione della medesima non è specificata (spazi o tabulazioni sono entrambi accettati). Le tabulazioni devono essere settate ogni otto spazi (non quattro).

\subsection{Lunghezza delle linee}
Evitare linee più lunghe di ottanta caratteri perché esse non vengono ben gestite da molti terminali e strumenti software. Per la documentazione si utilizza una più corta lunghezza di linea, generalmente non più di settanta caratteri.

\subsection{Spostamento di linee}
Quando un'espressione supera la lunghezza della linea, occorre spezzarla secondo i seguenti principi generali:
\begin{itemize}
  \item interrompere la linea dopo una virgola;
  \item interrompere la linea prima di un operatore;
  \item preferire interruzioni di alto livello rispetto ad interruzioni di basso livello;
  \item allineare la nuova linea con l'inizio dell'espressione nella linea precedente;
  \item se le regole precedenti rendono il codice più confuso o il codice è troppo spostato verso il margine destro utilizzare solo otto spazi di indentazione.
\end{itemize}
%\begin{lstlisting}
%  nomeMetodo(espressioneLunga1, espressioneLunga2, espressioneLunga3,
%                       espressioneLunga4, espressioneLunga5);
%  
%  var = nomeMetodo(espressioneLunga1,
%          nomeMetodo2(espressioneLunga2,
%              espressioneLunga3));
%\end{lstlisting}

\section{Commenti}
I programmi Java possono avere due tipi di commenti: commenti d’implementazione e commenti di documentazione. I commenti d’implementazione sono quelli classici del C++, che sono delimitati da /*...*/ e //. I commenti di documentazione (noti anche come doc comments) sono esclusivi del Java, e sono delimitati da /**...*/. I doc comments possono essere estratti in file HTML utilizzando lo strumento Javadoc.
I commenti di implementazione sono dei mezzi per commentare il codice o per commentare una particolare implementazione. I doc comments vengono utilizzati per descrivere la specifica del codice da una prospettiva non implementativa, per essere letti da sviluppatori che non devono necessariamente avere il codice in mano.
I commenti dovrebbero essere usati per dare una panoramica del codice e per fornire informazioni aggiuntive che non sono prontamente disponibili nel codice stesso. I commenti devono contenere solo informazioni rilevanti per leggere e comprendere il programma. Ad esempio, informazioni su come il package corrispondente è costruito o in quale directory risiede non dovrebbero essere incluse in un commento.

La discussione sulle decisioni non banali o non ovvie è adatta, ma bisogna evitare di duplicare le informazioni che sono presenti in maniera chiara nel codice. E’ molto facile che commenti ridondanti diventino obsoleti; in generale, si dovrebbe evitare di inserire commenti suscettibili di diventare obsoleti con l’evoluzione del software.
La frequenza dei commenti talvolta riflette una povera qualità del codice. Quando ci si sente obbligati ad aggiungere un commento, considerare il caso di riscrivere il codice per renderlo più chiaro.
I commenti non dovrebbero essere inclusi in grandi riquadri tracciati con asterischi o altri caratteri, né dovrebbero includere caratteri speciali come backspace.

\subsection{Formattazione commento di implementazione}
I programmi possono avere tre tipi di commenti di implementazione:
\begin{itemize}
  \item \textsc{Commenti di blocco}: sono usati per fornire descrizioni di file, metodi, strutture dati e algoritmi. I commenti di blocco possono essere usati all’inizio di ogni file e prima di ogni metodo. Possono inoltre essere usati in altri punti, come all’interno dei metodi. I commenti di blocco dentro una funzione o un metodo dovrebbero essere indentati allo stesso livello del codice che descrivono. Un commento di blocco dovrebbe essere preceduto da una linea bianca di separazione dal codice.
  \begin{lstlisting}
    /*
     * Questo e' un commento di blocco
     */
  \end{lstlisting}
  
  \item \textsc{Commenti a linea singola}: sono brevi commenti che possono apparire su una singola linea di codice ed indentati al livello del codice che seguono. Se un commento non può essere scritto su una linea singola, deve seguire il formato di commento di blocco. Un commento a linea singola deve essere preceduto da una linea bianca. Quello che segue è un esempio di commento a linea singola nel codice Java.
  \begin{lstlisting}
    if(condizione) {
      /*  Gestisce la condizione */
      ...
    }
  \end{lstlisting}
  
  \item \textsc{Commenti di fine linea}: Il delimitatore di commento // può commentare una linea completa o una parte di essa. Non dovrebbe essere usato su più linee consecutive per commenti testuali; comunque, può essere usato su più linee consecutive per commentare sezioni del codice. Seguono tre esempi dei tre stili.
  \begin{lstlisting}
    if(i>0) {
      //Fa qualcosa
      ...
    } else {
      i--; //Spiega il perche' qui
    }
    
    //if(x==0) {
    //
    //Fa qualcos'altro
    //...
    //}
  \end{lstlisting}
\end{itemize}


	\chapter{Class Interface}
\section{Classi x}
\begin{center}
\begin{longtable}{|a|p{.65\textwidth}|}
\hline 
Nome classe & 
AccountFacade
\tabularnewline
\hline 
Descrizione & 
Questa classe rappresenta un generico account
\tabularnewline
\hline 
Pre-condizioni & 
\textbf{context} AccountFacade::search(pAccount) \textbf{pre}: pAccount!=null
\tabularnewline
\hline 
Post-condizioni & 
\textbf{context} AccountFacade::search(pAccount) \textbf{post}: Account!=null
\tabularnewline
\hline 
Invarianti & 
\textbf{context} AccountFacade \textbf{inv}: pAccount!=null
\tabularnewline
\hline
\end{longtable}
\end{center}

\chapter{Documentazione del codice} {
% 2. TeXDoclet output include
\section*{Class Hierarchy}{
\thispagestyle{empty}
\markboth{Class Hierarchy}{Class Hierarchy}
\addcontentsline{toc}{section}{Class Hierarchy}
\subsection*{Classes}
{\raggedright
\hspace{0.0cm} $\bullet$ java.lang.Object {\tiny \refdefined{java.lang.Object}} \\
\hspace{1.0cm} $\bullet$ AbstractPreferenceInitializer {\tiny } \\
\hspace{2.0cm} $\bullet$ it.unisa.sesa.repominer.preferences.PreferenceInitializer {\tiny \refdefined{it.unisa.sesa.repominer.preferences.PreferenceInitializer}} \\
\hspace{1.0cm} $\bullet$ AbstractUIPlugin {\tiny } \\
\hspace{2.0cm} $\bullet$ it.unisa.sesa.repominer.Activator {\tiny \refdefined{it.unisa.sesa.repominer.Activator}} \\
\hspace{1.0cm} $\bullet$ FieldEditorPreferencePage {\tiny } \\
\hspace{2.0cm} $\bullet$ it.unisa.sesa.repominer.preferences.PreferencePage {\tiny \refdefined{it.unisa.sesa.repominer.preferences.PreferencePage}} \\
\hspace{1.0cm} $\bullet$ it.unisa.sesa.repominer.actions.HistoryMetricsAction {\tiny \refdefined{it.unisa.sesa.repominer.actions.HistoryMetricsAction}} \\
\hspace{1.0cm} $\bullet$ it.unisa.sesa.repominer.db.ChangeDAO {\tiny \refdefined{it.unisa.sesa.repominer.db.ChangeDAO}} \\
\hspace{1.0cm} $\bullet$ it.unisa.sesa.repominer.db.ChangeForCommitDAO {\tiny \refdefined{it.unisa.sesa.repominer.db.ChangeForCommitDAO}} \\
\hspace{1.0cm} $\bullet$ it.unisa.sesa.repominer.db.ConnectionPool {\tiny \refdefined{it.unisa.sesa.repominer.db.ConnectionPool}} \\
\hspace{1.0cm} $\bullet$ it.unisa.sesa.repominer.db.ImportDAO {\tiny \refdefined{it.unisa.sesa.repominer.db.ImportDAO}} \\
\hspace{1.0cm} $\bullet$ it.unisa.sesa.repominer.db.IssueDAO {\tiny \refdefined{it.unisa.sesa.repominer.db.IssueDAO}} \\
\hspace{1.0cm} $\bullet$ it.unisa.sesa.repominer.db.MethodDAO {\tiny \refdefined{it.unisa.sesa.repominer.db.MethodDAO}} \\
\hspace{1.0cm} $\bullet$ it.unisa.sesa.repominer.db.MetricDAO {\tiny \refdefined{it.unisa.sesa.repominer.db.MetricDAO}} \\
\hspace{1.0cm} $\bullet$ it.unisa.sesa.repominer.db.MetricMethodDAO {\tiny \refdefined{it.unisa.sesa.repominer.db.MetricMethodDAO}} \\
\hspace{1.0cm} $\bullet$ it.unisa.sesa.repominer.db.PackageMetricDAO {\tiny \refdefined{it.unisa.sesa.repominer.db.PackageMetricDAO}} \\
\hspace{1.0cm} $\bullet$ it.unisa.sesa.repominer.db.ProjectDAO {\tiny \refdefined{it.unisa.sesa.repominer.db.ProjectDAO}} \\
\hspace{1.0cm} $\bullet$ it.unisa.sesa.repominer.db.ProjectMetricDAO {\tiny \refdefined{it.unisa.sesa.repominer.db.ProjectMetricDAO}} \\
\hspace{1.0cm} $\bullet$ it.unisa.sesa.repominer.db.SourceContainerDAO {\tiny \refdefined{it.unisa.sesa.repominer.db.SourceContainerDAO}} \\
\hspace{1.0cm} $\bullet$ it.unisa.sesa.repominer.db.TypeDAO {\tiny \refdefined{it.unisa.sesa.repominer.db.TypeDAO}} \\
\hspace{1.0cm} $\bullet$ it.unisa.sesa.repominer.db.entities.CatchedException {\tiny \refdefined{it.unisa.sesa.repominer.db.entities.CatchedException}} \\
\hspace{1.0cm} $\bullet$ it.unisa.sesa.repominer.db.entities.Change {\tiny \refdefined{it.unisa.sesa.repominer.db.entities.Change}} \\
\hspace{1.0cm} $\bullet$ it.unisa.sesa.repominer.db.entities.ChangeForCommit {\tiny \refdefined{it.unisa.sesa.repominer.db.entities.ChangeForCommit}} \\
\hspace{1.0cm} $\bullet$ it.unisa.sesa.repominer.db.entities.ClassInvocation {\tiny \refdefined{it.unisa.sesa.repominer.db.entities.ClassInvocation}} \\
\hspace{1.0cm} $\bullet$ it.unisa.sesa.repominer.db.entities.Import {\tiny \refdefined{it.unisa.sesa.repominer.db.entities.Import}} \\
\hspace{1.0cm} $\bullet$ it.unisa.sesa.repominer.db.entities.Imports {\tiny \refdefined{it.unisa.sesa.repominer.db.entities.Imports}} \\
\hspace{1.0cm} $\bullet$ it.unisa.sesa.repominer.db.entities.Issue {\tiny \refdefined{it.unisa.sesa.repominer.db.entities.Issue}} \\
\hspace{1.0cm} $\bullet$ it.unisa.sesa.repominer.db.entities.IssueAttachment {\tiny \refdefined{it.unisa.sesa.repominer.db.entities.IssueAttachment}} \\
\hspace{1.0cm} $\bullet$ it.unisa.sesa.repominer.db.entities.IssueComment {\tiny \refdefined{it.unisa.sesa.repominer.db.entities.IssueComment}} \\
\hspace{1.0cm} $\bullet$ it.unisa.sesa.repominer.db.entities.Method {\tiny \refdefined{it.unisa.sesa.repominer.db.entities.Method}} \\
\hspace{1.0cm} $\bullet$ it.unisa.sesa.repominer.db.entities.MethodComment {\tiny \refdefined{it.unisa.sesa.repominer.db.entities.MethodComment}} \\
\hspace{1.0cm} $\bullet$ it.unisa.sesa.repominer.db.entities.MethodInvocation {\tiny \refdefined{it.unisa.sesa.repominer.db.entities.MethodInvocation}} \\
\hspace{1.0cm} $\bullet$ it.unisa.sesa.repominer.db.entities.MethodsChangeInCommit {\tiny \refdefined{it.unisa.sesa.repominer.db.entities.MethodsChangeInCommit}} \\
\hspace{1.0cm} $\bullet$ it.unisa.sesa.repominer.db.entities.Metric {\tiny \refdefined{it.unisa.sesa.repominer.db.entities.Metric}} \\
\hspace{2.0cm} $\bullet$ it.unisa.sesa.repominer.db.entities.MetricMethod {\tiny \refdefined{it.unisa.sesa.repominer.db.entities.MetricMethod}} \\
\hspace{2.0cm} $\bullet$ it.unisa.sesa.repominer.db.entities.PackageMetric {\tiny \refdefined{it.unisa.sesa.repominer.db.entities.PackageMetric}} \\
\hspace{2.0cm} $\bullet$ it.unisa.sesa.repominer.db.entities.ProjectMetric {\tiny \refdefined{it.unisa.sesa.repominer.db.entities.ProjectMetric}} \\
\hspace{1.0cm} $\bullet$ it.unisa.sesa.repominer.db.entities.Project {\tiny \refdefined{it.unisa.sesa.repominer.db.entities.Project}} \\
\hspace{1.0cm} $\bullet$ it.unisa.sesa.repominer.db.entities.Review {\tiny \refdefined{it.unisa.sesa.repominer.db.entities.Review}} \\
\hspace{1.0cm} $\bullet$ it.unisa.sesa.repominer.db.entities.SourceContainer {\tiny \refdefined{it.unisa.sesa.repominer.db.entities.SourceContainer}} \\
\hspace{1.0cm} $\bullet$ it.unisa.sesa.repominer.db.entities.Type {\tiny \refdefined{it.unisa.sesa.repominer.db.entities.Type}} \\
\hspace{1.0cm} $\bullet$ it.unisa.sesa.repominer.db.entities.TypeMetric {\tiny \refdefined{it.unisa.sesa.repominer.db.entities.TypeMetric}} \\
\hspace{1.0cm} $\bullet$ it.unisa.sesa.repominer.dbscan.ChangePoint {\tiny \refdefined{it.unisa.sesa.repominer.dbscan.ChangePoint}} \\
\hspace{1.0cm} $\bullet$ it.unisa.sesa.repominer.dbscan.Cluster {\tiny \refdefined{it.unisa.sesa.repominer.dbscan.Cluster}} \\
\hspace{1.0cm} $\bullet$ it.unisa.sesa.repominer.dbscan.DBSCAN {\tiny \refdefined{it.unisa.sesa.repominer.dbscan.DBSCAN}} \\
\hspace{1.0cm} $\bullet$ it.unisa.sesa.repominer.metrics.HistoryMetricsCalculator {\tiny \refdefined{it.unisa.sesa.repominer.metrics.HistoryMetricsCalculator}} \\
\hspace{1.0cm} $\bullet$ it.unisa.sesa.repominer.metrics.PackageMetrics {\tiny \refdefined{it.unisa.sesa.repominer.metrics.PackageMetrics}} \\
\hspace{1.0cm} $\bullet$ it.unisa.sesa.repominer.metrics.ProjectMetrics {\tiny \refdefined{it.unisa.sesa.repominer.metrics.ProjectMetrics}} \\
\hspace{1.0cm} $\bullet$ it.unisa.sesa.repominer.preferences.PreferenceConstants {\tiny \refdefined{it.unisa.sesa.repominer.preferences.PreferenceConstants}} \\
\hspace{1.0cm} $\bullet$ it.unisa.sesa.repominer.preferences.Preferences {\tiny \refdefined{it.unisa.sesa.repominer.preferences.Preferences}} \\
\hspace{1.0cm} $\bullet$ it.unisa.sesa.repominer.util.Utils {\tiny \refdefined{it.unisa.sesa.repominer.util.Utils}} \\
\hspace{1.0cm} $\bullet$ it.unisa.sesa.repominer.util.snippets.Test {\tiny \refdefined{it.unisa.sesa.repominer.util.snippets.Test}} \\
}
\subsection*{Exceptions}
\hspace{0.0cm} $\bullet$ java.lang.Object {\tiny \refdefined{java.lang.Object}} \\
\hspace{1.0cm} $\bullet$ java.lang.Throwable {\tiny \refdefined{java.lang.Throwable}} \\
\hspace{2.0cm} $\bullet$ java.lang.Exception {\tiny \refdefined{java.lang.Exception}} \\
\hspace{3.0cm} $\bullet$ it.unisa.sesa.repominer.metrics.exception.NoChangesException {\tiny \refdefined{it.unisa.sesa.repominer.metrics.exception.NoChangesException}} \\
\hspace{3.0cm} $\bullet$ it.unisa.sesa.repominer.preferences.exceptions.IntegerPreferenceException {\tiny \refdefined{it.unisa.sesa.repominer.preferences.exceptions.IntegerPreferenceException}} \\
\hspace{3.0cm} $\bullet$ it.unisa.sesa.repominer.preferences.exceptions.PeriodLengthTooLong {\tiny \refdefined{it.unisa.sesa.repominer.preferences.exceptions.PeriodLengthTooLong}} \\
}
\section{Package it.unisa.sesa.repominer.actions}{
\label{it.unisa.sesa.repominer.actions}\hypertarget{it.unisa.sesa.repominer.actions}{}
\hskip -.05in
\hbox to \hsize{\textit{ Package Contents\hfil Page}}
\vskip .13in
\hbox{{\bf  Classes}}
\entityintro{HistoryMetricsAction}{it.unisa.sesa.repominer.actions.HistoryMetricsAction}{Our sample action implements workbench action delegate.}
\vskip .1in
\vskip .1in
\subsection{\label{it.unisa.sesa.repominer.actions.HistoryMetricsAction}Class HistoryMetricsAction}{
\hypertarget{it.unisa.sesa.repominer.actions.HistoryMetricsAction}{}\vskip .1in 
Our sample action implements workbench action delegate. The action proxy will be created by the workbench and shown in the UI. When the user tries to use the action, this delegate will be created and execution will be delegated to it.\vskip .1in 
\subsubsection{See also}{}

  \begin{list}{-- }{\setlength{\itemsep}{0cm}\setlength{\parsep}{0cm}}
\item{ IWorkbenchWindowActionDelegate} 
  \end{list}
\subsubsection{Declaration}{
\small public class HistoryMetricsAction
\\ {\bf  extends} java.lang.Object
\refdefined{java.lang.Object}}
\subsubsection{Constructor summary}{
\begin{verse}
\hyperlink{it.unisa.sesa.repominer.actions.HistoryMetricsAction()}{{\bf HistoryMetricsAction()}} The constructor.\\
\end{verse}
}
\subsubsection{Method summary}{
\begin{verse}
\hyperlink{it.unisa.sesa.repominer.actions.HistoryMetricsAction.dispose()}{{\bf dispose()}} We can use this method to dispose of any system resources we previously allocated.\\
\hyperlink{it.unisa.sesa.repominer.actions.HistoryMetricsAction.init(IWorkbenchWindow)}{{\bf init(IWorkbenchWindow)}} We will cache window object in order to be able to provide parent shell for the message dialog.\\
\hyperlink{it.unisa.sesa.repominer.actions.HistoryMetricsAction.run(IAction)}{{\bf run(IAction)}} The action has been activated.\\
\hyperlink{it.unisa.sesa.repominer.actions.HistoryMetricsAction.selectionChanged(IAction, ISelection)}{{\bf selectionChanged(IAction, ISelection)}} Selection in the workbench has been changed.\\
\end{verse}
}
\subsubsection{Constructors}{
\vskip -2em
\begin{itemize}
\item{ 
\index{HistoryMetricsAction()}
\hypertarget{it.unisa.sesa.repominer.actions.HistoryMetricsAction()}{{\bf  HistoryMetricsAction}\\}
\texttt{public\ {\bf  HistoryMetricsAction}()
\label{it.unisa.sesa.repominer.actions.HistoryMetricsAction()}}%end signature
\begin{itemize}
\item{
{\bf  Description}

The constructor.
}
\end{itemize}
}%end item
\end{itemize}
}
\subsubsection{Methods}{
\vskip -2em
\begin{itemize}
\item{ 
\index{dispose()}
\hypertarget{it.unisa.sesa.repominer.actions.HistoryMetricsAction.dispose()}{{\bf  dispose}\\}
\texttt{public void\ {\bf  dispose}()
\label{it.unisa.sesa.repominer.actions.HistoryMetricsAction.dispose()}}%end signature
\begin{itemize}
\item{
{\bf  Description}

We can use this method to dispose of any system resources we previously allocated.
}
\item{{\bf  See also}
  \begin{itemize}
\item{ IWorkbenchWindowActionDelegate\#dispose}
  \end{itemize}
}%end item
\end{itemize}
}%end item
\item{ 
\index{init(IWorkbenchWindow)}
\hypertarget{it.unisa.sesa.repominer.actions.HistoryMetricsAction.init(IWorkbenchWindow)}{{\bf  init}\\}
\texttt{public void\ {\bf  init}(\texttt{IWorkbenchWindow} {\bf  window})
\label{it.unisa.sesa.repominer.actions.HistoryMetricsAction.init(IWorkbenchWindow)}}%end signature
\begin{itemize}
\item{
{\bf  Description}

We will cache window object in order to be able to provide parent shell for the message dialog.
}
\item{{\bf  See also}
  \begin{itemize}
\item{ IWorkbenchWindowActionDelegate\#init}
  \end{itemize}
}%end item
\end{itemize}
}%end item
\item{ 
\index{run(IAction)}
\hypertarget{it.unisa.sesa.repominer.actions.HistoryMetricsAction.run(IAction)}{{\bf  run}\\}
\texttt{public void\ {\bf  run}(\texttt{IAction} {\bf  action})
\label{it.unisa.sesa.repominer.actions.HistoryMetricsAction.run(IAction)}}%end signature
\begin{itemize}
\item{
{\bf  Description}

The action has been activated. The argument of the method represents the 'real' action sitting in the workbench UI.
}
\item{{\bf  See also}
  \begin{itemize}
\item{ IWorkbenchWindowActionDelegate\#run}
  \end{itemize}
}%end item
\end{itemize}
}%end item
\item{ 
\index{selectionChanged(IAction, ISelection)}
\hypertarget{it.unisa.sesa.repominer.actions.HistoryMetricsAction.selectionChanged(IAction, ISelection)}{{\bf  selectionChanged}\\}
\texttt{public void\ {\bf  selectionChanged}(\texttt{IAction} {\bf  action},
\texttt{ISelection} {\bf  selection})
\label{it.unisa.sesa.repominer.actions.HistoryMetricsAction.selectionChanged(IAction, ISelection)}}%end signature
\begin{itemize}
\item{
{\bf  Description}

Selection in the workbench has been changed. We can change the state of the 'real' action here if we want, but this can only happen after the delegate has been created.
}
\item{{\bf  See also}
  \begin{itemize}
\item{ IWorkbenchWindowActionDelegate\#selectionChanged}
  \end{itemize}
}%end item
\end{itemize}
}%end item
\end{itemize}
}
}
}
\section{Package it.unisa.sesa.repominer}{
\label{it.unisa.sesa.repominer}\hypertarget{it.unisa.sesa.repominer}{}
\hskip -.05in
\hbox to \hsize{\textit{ Package Contents\hfil Page}}
\vskip .13in
\hbox{{\bf  Classes}}
\entityintro{Activator}{it.unisa.sesa.repominer.Activator}{The activator class controls the plug-in life cycle}
\vskip .1in
\vskip .1in
\subsection{\label{it.unisa.sesa.repominer.Activator}Class Activator}{
\hypertarget{it.unisa.sesa.repominer.Activator}{}\vskip .1in 
The activator class controls the plug-in life cycle\vskip .1in 
\subsubsection{Declaration}{
\small public class Activator
\\ {\bf  extends} AbstractUIPlugin
\refdefined{.AbstractUIPlugin}}
\subsubsection{Field summary}{
\begin{verse}
\hyperlink{it.unisa.sesa.repominer.Activator.PLUGIN_ID}{{\bf PLUGIN\_ID}} \\
\end{verse}
}
\subsubsection{Constructor summary}{
\begin{verse}
\hyperlink{it.unisa.sesa.repominer.Activator()}{{\bf Activator()}} The constructor\\
\end{verse}
}
\subsubsection{Method summary}{
\begin{verse}
\hyperlink{it.unisa.sesa.repominer.Activator.getDefault()}{{\bf getDefault()}} Returns the shared instance\\
\hyperlink{it.unisa.sesa.repominer.Activator.getImageDescriptor(java.lang.String)}{{\bf getImageDescriptor(String)}} Returns an image descriptor for the image file at the given plug-in relative path\\
\hyperlink{it.unisa.sesa.repominer.Activator.start(BundleContext)}{{\bf start(BundleContext)}} \\
\hyperlink{it.unisa.sesa.repominer.Activator.stop(BundleContext)}{{\bf stop(BundleContext)}} \\
\end{verse}
}
\subsubsection{Fields}{
\begin{itemize}
\item{
\index{PLUGIN\_ID}
\label{it.unisa.sesa.repominer.Activator.PLUGIN_ID}\hypertarget{it.unisa.sesa.repominer.Activator.PLUGIN_ID}{public static final java.lang.String {\bf  PLUGIN\_ID}}
}
\end{itemize}
}
\subsubsection{Constructors}{
\vskip -2em
\begin{itemize}
\item{ 
\index{Activator()}
\hypertarget{it.unisa.sesa.repominer.Activator()}{{\bf  Activator}\\}
\texttt{public\ {\bf  Activator}()
\label{it.unisa.sesa.repominer.Activator()}}%end signature
\begin{itemize}
\item{
{\bf  Description}

The constructor
}
\end{itemize}
}%end item
\end{itemize}
}
\subsubsection{Methods}{
\vskip -2em
\begin{itemize}
\item{ 
\index{getDefault()}
\hypertarget{it.unisa.sesa.repominer.Activator.getDefault()}{{\bf  getDefault}\\}
\texttt{public static Activator\ {\bf  getDefault}()
\label{it.unisa.sesa.repominer.Activator.getDefault()}}%end signature
\begin{itemize}
\item{
{\bf  Description}

Returns the shared instance
}
\item{{\bf  Returns} -- 
the shared instance 
}%end item
\end{itemize}
}%end item
\item{ 
\index{getImageDescriptor(String)}
\hypertarget{it.unisa.sesa.repominer.Activator.getImageDescriptor(java.lang.String)}{{\bf  getImageDescriptor}\\}
\texttt{public static ImageDescriptor\ {\bf  getImageDescriptor}(\texttt{java.lang.String} {\bf  path})
\label{it.unisa.sesa.repominer.Activator.getImageDescriptor(java.lang.String)}}%end signature
\begin{itemize}
\item{
{\bf  Description}

Returns an image descriptor for the image file at the given plug-in relative path
}
\item{
{\bf  Parameters}
  \begin{itemize}
   \item{
\texttt{path} -- the path}
  \end{itemize}
}%end item
\item{{\bf  Returns} -- 
the image descriptor 
}%end item
\end{itemize}
}%end item
\item{ 
\index{start(BundleContext)}
\hypertarget{it.unisa.sesa.repominer.Activator.start(BundleContext)}{{\bf  start}\\}
\texttt{public void\ {\bf  start}(\texttt{BundleContext} {\bf  context}) throws java.lang.Exception
\label{it.unisa.sesa.repominer.Activator.start(BundleContext)}}%end signature
}%end item
\item{ 
\index{stop(BundleContext)}
\hypertarget{it.unisa.sesa.repominer.Activator.stop(BundleContext)}{{\bf  stop}\\}
\texttt{public void\ {\bf  stop}(\texttt{BundleContext} {\bf  context}) throws java.lang.Exception
\label{it.unisa.sesa.repominer.Activator.stop(BundleContext)}}%end signature
}%end item
\end{itemize}
}
}
}
\section{Package it.unisa.sesa.repominer.db}{
\label{it.unisa.sesa.repominer.db}\hypertarget{it.unisa.sesa.repominer.db}{}
\hskip -.05in
\hbox to \hsize{\textit{ Package Contents\hfil Page}}
\vskip .13in
\hbox{{\bf  Classes}}
\entityintro{ChangeDAO}{it.unisa.sesa.repominer.db.ChangeDAO}{}
\entityintro{ChangeForCommitDAO}{it.unisa.sesa.repominer.db.ChangeForCommitDAO}{}
\entityintro{ConnectionPool}{it.unisa.sesa.repominer.db.ConnectionPool}{}
\entityintro{ImportDAO}{it.unisa.sesa.repominer.db.ImportDAO}{}
\entityintro{IssueDAO}{it.unisa.sesa.repominer.db.IssueDAO}{}
\entityintro{MethodDAO}{it.unisa.sesa.repominer.db.MethodDAO}{}
\entityintro{MetricDAO}{it.unisa.sesa.repominer.db.MetricDAO}{}
\entityintro{MetricMethodDAO}{it.unisa.sesa.repominer.db.MetricMethodDAO}{}
\entityintro{PackageMetricDAO}{it.unisa.sesa.repominer.db.PackageMetricDAO}{}
\entityintro{ProjectDAO}{it.unisa.sesa.repominer.db.ProjectDAO}{}
\entityintro{ProjectMetricDAO}{it.unisa.sesa.repominer.db.ProjectMetricDAO}{}
\entityintro{SourceContainerDAO}{it.unisa.sesa.repominer.db.SourceContainerDAO}{}
\entityintro{TypeDAO}{it.unisa.sesa.repominer.db.TypeDAO}{}
\vskip .1in
\vskip .1in
\subsection{\label{it.unisa.sesa.repominer.db.ChangeDAO}Class ChangeDAO}{
\hypertarget{it.unisa.sesa.repominer.db.ChangeDAO}{}\vskip .1in 
\subsubsection{Declaration}{
\small public class ChangeDAO
\\ {\bf  extends} java.lang.Object
\refdefined{java.lang.Object}}
\subsubsection{Constructor summary}{
\begin{verse}
\hyperlink{it.unisa.sesa.repominer.db.ChangeDAO()}{{\bf ChangeDAO()}} \\
\end{verse}
}
\subsubsection{Method summary}{
\begin{verse}
\hyperlink{it.unisa.sesa.repominer.db.ChangeDAO.getChangeById(java.lang.Integer)}{{\bf getChangeById(Integer)}} This method return a single commit picked by id\\
\hyperlink{it.unisa.sesa.repominer.db.ChangeDAO.getChangesByDateInterval(it.unisa.sesa.repominer.db.entities.Project, java.util.Date, java.util.Date)}{{\bf getChangesByDateInterval(Project, Date, Date)}} This method return all commit occurred between two given dates in project passed as parameter\\
\hyperlink{it.unisa.sesa.repominer.db.ChangeDAO.getChangesOfProject(it.unisa.sesa.repominer.db.entities.Project)}{{\bf getChangesOfProject(Project)}} This method returns all commits occurred in project passed as parameter\\
\hyperlink{it.unisa.sesa.repominer.db.ChangeDAO.getProjectEndDate(it.unisa.sesa.repominer.db.entities.Project)}{{\bf getProjectEndDate(Project)}} This method return the date of last change occurred in project passed as parameter\\
\hyperlink{it.unisa.sesa.repominer.db.ChangeDAO.getProjectStartDate(it.unisa.sesa.repominer.db.entities.Project)}{{\bf getProjectStartDate(Project)}} This method return the date of first change occurred in project passed as parameter\\
\end{verse}
}
\subsubsection{Constructors}{
\vskip -2em
\begin{itemize}
\item{ 
\index{ChangeDAO()}
\hypertarget{it.unisa.sesa.repominer.db.ChangeDAO()}{{\bf  ChangeDAO}\\}
\texttt{public\ {\bf  ChangeDAO}()
\label{it.unisa.sesa.repominer.db.ChangeDAO()}}%end signature
}%end item
\end{itemize}
}
\subsubsection{Methods}{
\vskip -2em
\begin{itemize}
\item{ 
\index{getChangeById(Integer)}
\hypertarget{it.unisa.sesa.repominer.db.ChangeDAO.getChangeById(java.lang.Integer)}{{\bf  getChangeById}\\}
\texttt{public entities.Change\ {\bf  getChangeById}(\texttt{java.lang.Integer} {\bf  pId})
\label{it.unisa.sesa.repominer.db.ChangeDAO.getChangeById(java.lang.Integer)}}%end signature
\begin{itemize}
\item{
{\bf  Description}

This method return a single commit picked by id
}
\item{
{\bf  Parameters}
  \begin{itemize}
   \item{
\texttt{pId} -- }
  \end{itemize}
}%end item
\item{{\bf  Returns} -- 
A Change object 
}%end item
\end{itemize}
}%end item
\item{ 
\index{getChangesByDateInterval(Project, Date, Date)}
\hypertarget{it.unisa.sesa.repominer.db.ChangeDAO.getChangesByDateInterval(it.unisa.sesa.repominer.db.entities.Project, java.util.Date, java.util.Date)}{{\bf  getChangesByDateInterval}\\}
\texttt{public java.util.List\ {\bf  getChangesByDateInterval}(\texttt{entities.Project} {\bf  pProject},
\texttt{java.util.Date} {\bf  pDate1},
\texttt{java.util.Date} {\bf  pDate2})
\label{it.unisa.sesa.repominer.db.ChangeDAO.getChangesByDateInterval(it.unisa.sesa.repominer.db.entities.Project, java.util.Date, java.util.Date)}}%end signature
\begin{itemize}
\item{
{\bf  Description}

This method return all commit occurred between two given dates in project passed as parameter
}
\item{
{\bf  Parameters}
  \begin{itemize}
   \item{
\texttt{pProject} -- }
   \item{
\texttt{pDate1} -- }
   \item{
\texttt{pDate2} -- }
  \end{itemize}
}%end item
\item{{\bf  Returns} -- 
A list of Change objects 
}%end item
\end{itemize}
}%end item
\item{ 
\index{getChangesOfProject(Project)}
\hypertarget{it.unisa.sesa.repominer.db.ChangeDAO.getChangesOfProject(it.unisa.sesa.repominer.db.entities.Project)}{{\bf  getChangesOfProject}\\}
\texttt{public java.util.List\ {\bf  getChangesOfProject}(\texttt{entities.Project} {\bf  pProject})
\label{it.unisa.sesa.repominer.db.ChangeDAO.getChangesOfProject(it.unisa.sesa.repominer.db.entities.Project)}}%end signature
\begin{itemize}
\item{
{\bf  Description}

This method returns all commits occurred in project passed as parameter
}
\item{
{\bf  Parameters}
  \begin{itemize}
   \item{
\texttt{pProject} -- }
  \end{itemize}
}%end item
\item{{\bf  Returns} -- 
A list of Change objects 
}%end item
\end{itemize}
}%end item
\item{ 
\index{getProjectEndDate(Project)}
\hypertarget{it.unisa.sesa.repominer.db.ChangeDAO.getProjectEndDate(it.unisa.sesa.repominer.db.entities.Project)}{{\bf  getProjectEndDate}\\}
\texttt{public java.util.Date\ {\bf  getProjectEndDate}(\texttt{entities.Project} {\bf  pProject})
\label{it.unisa.sesa.repominer.db.ChangeDAO.getProjectEndDate(it.unisa.sesa.repominer.db.entities.Project)}}%end signature
\begin{itemize}
\item{
{\bf  Description}

This method return the date of last change occurred in project passed as parameter
}
\item{
{\bf  Parameters}
  \begin{itemize}
   \item{
\texttt{pProject} -- }
  \end{itemize}
}%end item
\item{{\bf  Returns} -- 
A Date (last in project); return current date if query get a null change 
}%end item
\end{itemize}
}%end item
\item{ 
\index{getProjectStartDate(Project)}
\hypertarget{it.unisa.sesa.repominer.db.ChangeDAO.getProjectStartDate(it.unisa.sesa.repominer.db.entities.Project)}{{\bf  getProjectStartDate}\\}
\texttt{public java.util.Date\ {\bf  getProjectStartDate}(\texttt{entities.Project} {\bf  pProject})
\label{it.unisa.sesa.repominer.db.ChangeDAO.getProjectStartDate(it.unisa.sesa.repominer.db.entities.Project)}}%end signature
\begin{itemize}
\item{
{\bf  Description}

This method return the date of first change occurred in project passed as parameter
}
\item{
{\bf  Parameters}
  \begin{itemize}
   \item{
\texttt{pProject} -- }
  \end{itemize}
}%end item
\item{{\bf  Returns} -- 
A Date (first in project); return current date if query get a null Change 
}%end item
\end{itemize}
}%end item
\end{itemize}
}
}
\subsection{\label{it.unisa.sesa.repominer.db.ChangeForCommitDAO}Class ChangeForCommitDAO}{
\hypertarget{it.unisa.sesa.repominer.db.ChangeForCommitDAO}{}\vskip .1in 
\subsubsection{Declaration}{
\small public class ChangeForCommitDAO
\\ {\bf  extends} java.lang.Object
\refdefined{java.lang.Object}}
\subsubsection{Constructor summary}{
\begin{verse}
\hyperlink{it.unisa.sesa.repominer.db.ChangeForCommitDAO()}{{\bf ChangeForCommitDAO()}} \\
\end{verse}
}
\subsubsection{Method summary}{
\begin{verse}
\hyperlink{it.unisa.sesa.repominer.db.ChangeForCommitDAO.getChangeForCommitById(java.lang.Integer)}{{\bf getChangeForCommitById(Integer)}} This method returns a single change occurred in a commit picked by id\\
\hyperlink{it.unisa.sesa.repominer.db.ChangeForCommitDAO.getChangeForCommitOfChange(it.unisa.sesa.repominer.db.entities.Change)}{{\bf getChangeForCommitOfChange(Change)}} This method returns all changes occurred for a single commit passed as parameter\\
\hyperlink{it.unisa.sesa.repominer.db.ChangeForCommitDAO.getChangesForCommitByProject(it.unisa.sesa.repominer.db.entities.Project)}{{\bf getChangesForCommitByProject(Project)}} This method returns all changes occurred in all commits which refer to a project passed as parameter\\
\end{verse}
}
\subsubsection{Constructors}{
\vskip -2em
\begin{itemize}
\item{ 
\index{ChangeForCommitDAO()}
\hypertarget{it.unisa.sesa.repominer.db.ChangeForCommitDAO()}{{\bf  ChangeForCommitDAO}\\}
\texttt{public\ {\bf  ChangeForCommitDAO}()
\label{it.unisa.sesa.repominer.db.ChangeForCommitDAO()}}%end signature
}%end item
\end{itemize}
}
\subsubsection{Methods}{
\vskip -2em
\begin{itemize}
\item{ 
\index{getChangeForCommitById(Integer)}
\hypertarget{it.unisa.sesa.repominer.db.ChangeForCommitDAO.getChangeForCommitById(java.lang.Integer)}{{\bf  getChangeForCommitById}\\}
\texttt{public entities.ChangeForCommit\ {\bf  getChangeForCommitById}(\texttt{java.lang.Integer} {\bf  pId})
\label{it.unisa.sesa.repominer.db.ChangeForCommitDAO.getChangeForCommitById(java.lang.Integer)}}%end signature
\begin{itemize}
\item{
{\bf  Description}

This method returns a single change occurred in a commit picked by id
}
\item{
{\bf  Parameters}
  \begin{itemize}
   \item{
\texttt{pId} -- }
  \end{itemize}
}%end item
\item{{\bf  Returns} -- 
A ChangeForCommit object 
}%end item
\end{itemize}
}%end item
\item{ 
\index{getChangeForCommitOfChange(Change)}
\hypertarget{it.unisa.sesa.repominer.db.ChangeForCommitDAO.getChangeForCommitOfChange(it.unisa.sesa.repominer.db.entities.Change)}{{\bf  getChangeForCommitOfChange}\\}
\texttt{public java.util.List\ {\bf  getChangeForCommitOfChange}(\texttt{entities.Change} {\bf  pChange})
\label{it.unisa.sesa.repominer.db.ChangeForCommitDAO.getChangeForCommitOfChange(it.unisa.sesa.repominer.db.entities.Change)}}%end signature
\begin{itemize}
\item{
{\bf  Description}

This method returns all changes occurred for a single commit passed as parameter
}
\item{
{\bf  Parameters}
  \begin{itemize}
   \item{
\texttt{pChange} -- }
  \end{itemize}
}%end item
\item{{\bf  Returns} -- 
A list of ChangeForCommit objects 
}%end item
\end{itemize}
}%end item
\item{ 
\index{getChangesForCommitByProject(Project)}
\hypertarget{it.unisa.sesa.repominer.db.ChangeForCommitDAO.getChangesForCommitByProject(it.unisa.sesa.repominer.db.entities.Project)}{{\bf  getChangesForCommitByProject}\\}
\texttt{public java.util.List\ {\bf  getChangesForCommitByProject}(\texttt{entities.Project} {\bf  pProject})
\label{it.unisa.sesa.repominer.db.ChangeForCommitDAO.getChangesForCommitByProject(it.unisa.sesa.repominer.db.entities.Project)}}%end signature
\begin{itemize}
\item{
{\bf  Description}

This method returns all changes occurred in all commits which refer to a project passed as parameter
}
\item{
{\bf  Parameters}
  \begin{itemize}
   \item{
\texttt{pProject} -- }
  \end{itemize}
}%end item
\item{{\bf  Returns} -- 
A list of ChangeForCommit objects 
}%end item
\end{itemize}
}%end item
\end{itemize}
}
}
\subsection{\label{it.unisa.sesa.repominer.db.ConnectionPool}Class ConnectionPool}{
\hypertarget{it.unisa.sesa.repominer.db.ConnectionPool}{}\vskip .1in 
\subsubsection{Declaration}{
\small public class ConnectionPool
\\ {\bf  extends} java.lang.Object
\refdefined{java.lang.Object}}
\subsubsection{Method summary}{
\begin{verse}
\hyperlink{it.unisa.sesa.repominer.db.ConnectionPool.getConnection()}{{\bf getConnection()}} \\
\hyperlink{it.unisa.sesa.repominer.db.ConnectionPool.getInstance()}{{\bf getInstance()}} \\
\hyperlink{it.unisa.sesa.repominer.db.ConnectionPool.releaseConnection(java.sql.Connection)}{{\bf releaseConnection(Connection)}} \\
\end{verse}
}
\subsubsection{Methods}{
\vskip -2em
\begin{itemize}
\item{ 
\index{getConnection()}
\hypertarget{it.unisa.sesa.repominer.db.ConnectionPool.getConnection()}{{\bf  getConnection}\\}
\texttt{public java.sql.Connection\ {\bf  getConnection}()
\label{it.unisa.sesa.repominer.db.ConnectionPool.getConnection()}}%end signature
}%end item
\item{ 
\index{getInstance()}
\hypertarget{it.unisa.sesa.repominer.db.ConnectionPool.getInstance()}{{\bf  getInstance}\\}
\texttt{public static ConnectionPool\ {\bf  getInstance}()
\label{it.unisa.sesa.repominer.db.ConnectionPool.getInstance()}}%end signature
}%end item
\item{ 
\index{releaseConnection(Connection)}
\hypertarget{it.unisa.sesa.repominer.db.ConnectionPool.releaseConnection(java.sql.Connection)}{{\bf  releaseConnection}\\}
\texttt{public void\ {\bf  releaseConnection}(\texttt{java.sql.Connection} {\bf  connection})
\label{it.unisa.sesa.repominer.db.ConnectionPool.releaseConnection(java.sql.Connection)}}%end signature
}%end item
\end{itemize}
}
}
\subsection{\label{it.unisa.sesa.repominer.db.ImportDAO}Class ImportDAO}{
\hypertarget{it.unisa.sesa.repominer.db.ImportDAO}{}\vskip .1in 
\subsubsection{Declaration}{
\small public class ImportDAO
\\ {\bf  extends} java.lang.Object
\refdefined{java.lang.Object}}
\subsubsection{Constructor summary}{
\begin{verse}
\hyperlink{it.unisa.sesa.repominer.db.ImportDAO()}{{\bf ImportDAO()}} \\
\end{verse}
}
\subsubsection{Method summary}{
\begin{verse}
\hyperlink{it.unisa.sesa.repominer.db.ImportDAO.getImportById(java.lang.Integer)}{{\bf getImportById(Integer)}} This method returns a name of package picked by its id, passed as parameter\\
\end{verse}
}
\subsubsection{Constructors}{
\vskip -2em
\begin{itemize}
\item{ 
\index{ImportDAO()}
\hypertarget{it.unisa.sesa.repominer.db.ImportDAO()}{{\bf  ImportDAO}\\}
\texttt{public\ {\bf  ImportDAO}()
\label{it.unisa.sesa.repominer.db.ImportDAO()}}%end signature
}%end item
\end{itemize}
}
\subsubsection{Methods}{
\vskip -2em
\begin{itemize}
\item{ 
\index{getImportById(Integer)}
\hypertarget{it.unisa.sesa.repominer.db.ImportDAO.getImportById(java.lang.Integer)}{{\bf  getImportById}\\}
\texttt{public entities.Import\ {\bf  getImportById}(\texttt{java.lang.Integer} {\bf  pId})
\label{it.unisa.sesa.repominer.db.ImportDAO.getImportById(java.lang.Integer)}}%end signature
\begin{itemize}
\item{
{\bf  Description}

This method returns a name of package picked by its id, passed as parameter
}
\item{
{\bf  Parameters}
  \begin{itemize}
   \item{
\texttt{pId} -- }
  \end{itemize}
}%end item
\item{{\bf  Returns} -- 
An Import object 
}%end item
\end{itemize}
}%end item
\end{itemize}
}
}
\subsection{\label{it.unisa.sesa.repominer.db.IssueDAO}Class IssueDAO}{
\hypertarget{it.unisa.sesa.repominer.db.IssueDAO}{}\vskip .1in 
\subsubsection{Declaration}{
\small public class IssueDAO
\\ {\bf  extends} java.lang.Object
\refdefined{java.lang.Object}}
\subsubsection{Constructor summary}{
\begin{verse}
\hyperlink{it.unisa.sesa.repominer.db.IssueDAO()}{{\bf IssueDAO()}} \\
\end{verse}
}
\subsubsection{Method summary}{
\begin{verse}
\hyperlink{it.unisa.sesa.repominer.db.IssueDAO.getIssueById(java.lang.Integer)}{{\bf getIssueById(Integer)}} Returns a single issue picked by its id, picked as parameter\\
\hyperlink{it.unisa.sesa.repominer.db.IssueDAO.getIssuesByProject(it.unisa.sesa.repominer.db.entities.Project)}{{\bf getIssuesByProject(Project)}} Returns all issues for a project passed as parameter\\
\end{verse}
}
\subsubsection{Constructors}{
\vskip -2em
\begin{itemize}
\item{ 
\index{IssueDAO()}
\hypertarget{it.unisa.sesa.repominer.db.IssueDAO()}{{\bf  IssueDAO}\\}
\texttt{public\ {\bf  IssueDAO}()
\label{it.unisa.sesa.repominer.db.IssueDAO()}}%end signature
}%end item
\end{itemize}
}
\subsubsection{Methods}{
\vskip -2em
\begin{itemize}
\item{ 
\index{getIssueById(Integer)}
\hypertarget{it.unisa.sesa.repominer.db.IssueDAO.getIssueById(java.lang.Integer)}{{\bf  getIssueById}\\}
\texttt{public entities.Issue\ {\bf  getIssueById}(\texttt{java.lang.Integer} {\bf  pId})
\label{it.unisa.sesa.repominer.db.IssueDAO.getIssueById(java.lang.Integer)}}%end signature
\begin{itemize}
\item{
{\bf  Description}

Returns a single issue picked by its id, picked as parameter
}
\item{
{\bf  Parameters}
  \begin{itemize}
   \item{
\texttt{pId} -- }
  \end{itemize}
}%end item
\item{{\bf  Returns} -- 
An Issue object 
}%end item
\end{itemize}
}%end item
\item{ 
\index{getIssuesByProject(Project)}
\hypertarget{it.unisa.sesa.repominer.db.IssueDAO.getIssuesByProject(it.unisa.sesa.repominer.db.entities.Project)}{{\bf  getIssuesByProject}\\}
\texttt{public java.util.List\ {\bf  getIssuesByProject}(\texttt{entities.Project} {\bf  pProject})
\label{it.unisa.sesa.repominer.db.IssueDAO.getIssuesByProject(it.unisa.sesa.repominer.db.entities.Project)}}%end signature
\begin{itemize}
\item{
{\bf  Description}

Returns all issues for a project passed as parameter
}
\item{
{\bf  Parameters}
  \begin{itemize}
   \item{
\texttt{pProject} -- }
  \end{itemize}
}%end item
\item{{\bf  Returns} -- 
A list of Issue objects 
}%end item
\end{itemize}
}%end item
\end{itemize}
}
}
\subsection{\label{it.unisa.sesa.repominer.db.MethodDAO}Class MethodDAO}{
\hypertarget{it.unisa.sesa.repominer.db.MethodDAO}{}\vskip .1in 
\subsubsection{Declaration}{
\small public class MethodDAO
\\ {\bf  extends} java.lang.Object
\refdefined{java.lang.Object}}
\subsubsection{Constructor summary}{
\begin{verse}
\hyperlink{it.unisa.sesa.repominer.db.MethodDAO()}{{\bf MethodDAO()}} \\
\end{verse}
}
\subsubsection{Method summary}{
\begin{verse}
\hyperlink{it.unisa.sesa.repominer.db.MethodDAO.getMethodsOfType(it.unisa.sesa.repominer.db.entities.Type)}{{\bf getMethodsOfType(Type)}} This method return the list of methods for a class\\
\end{verse}
}
\subsubsection{Constructors}{
\vskip -2em
\begin{itemize}
\item{ 
\index{MethodDAO()}
\hypertarget{it.unisa.sesa.repominer.db.MethodDAO()}{{\bf  MethodDAO}\\}
\texttt{public\ {\bf  MethodDAO}()
\label{it.unisa.sesa.repominer.db.MethodDAO()}}%end signature
}%end item
\end{itemize}
}
\subsubsection{Methods}{
\vskip -2em
\begin{itemize}
\item{ 
\index{getMethodsOfType(Type)}
\hypertarget{it.unisa.sesa.repominer.db.MethodDAO.getMethodsOfType(it.unisa.sesa.repominer.db.entities.Type)}{{\bf  getMethodsOfType}\\}
\texttt{public java.util.List\ {\bf  getMethodsOfType}(\texttt{entities.Type} {\bf  pType})
\label{it.unisa.sesa.repominer.db.MethodDAO.getMethodsOfType(it.unisa.sesa.repominer.db.entities.Type)}}%end signature
\begin{itemize}
\item{
{\bf  Description}

This method return the list of methods for a class
}
\item{
{\bf  Parameters}
  \begin{itemize}
   \item{
\texttt{pType} -- }
  \end{itemize}
}%end item
\item{{\bf  Returns} -- 
A list of Method objects 
}%end item
\end{itemize}
}%end item
\end{itemize}
}
}
\subsection{\label{it.unisa.sesa.repominer.db.MetricDAO}Class MetricDAO}{
\hypertarget{it.unisa.sesa.repominer.db.MetricDAO}{}\vskip .1in 
\subsubsection{Declaration}{
\small public class MetricDAO
\\ {\bf  extends} java.lang.Object
\refdefined{java.lang.Object}}
\subsubsection{Constructor summary}{
\begin{verse}
\hyperlink{it.unisa.sesa.repominer.db.MetricDAO()}{{\bf MetricDAO()}} \\
\end{verse}
}
\subsubsection{Method summary}{
\begin{verse}
\hyperlink{it.unisa.sesa.repominer.db.MetricDAO.saveMetric(it.unisa.sesa.repominer.db.entities.Metric)}{{\bf saveMetric(Metric)}} This method saves in database a metric passed as parameter\\
\end{verse}
}
\subsubsection{Constructors}{
\vskip -2em
\begin{itemize}
\item{ 
\index{MetricDAO()}
\hypertarget{it.unisa.sesa.repominer.db.MetricDAO()}{{\bf  MetricDAO}\\}
\texttt{public\ {\bf  MetricDAO}()
\label{it.unisa.sesa.repominer.db.MetricDAO()}}%end signature
}%end item
\end{itemize}
}
\subsubsection{Methods}{
\vskip -2em
\begin{itemize}
\item{ 
\index{saveMetric(Metric)}
\hypertarget{it.unisa.sesa.repominer.db.MetricDAO.saveMetric(it.unisa.sesa.repominer.db.entities.Metric)}{{\bf  saveMetric}\\}
\texttt{public java.lang.Integer\ {\bf  saveMetric}(\texttt{entities.Metric} {\bf  pMetric})
\label{it.unisa.sesa.repominer.db.MetricDAO.saveMetric(it.unisa.sesa.repominer.db.entities.Metric)}}%end signature
\begin{itemize}
\item{
{\bf  Description}

This method saves in database a metric passed as parameter
}
\item{
{\bf  Parameters}
  \begin{itemize}
   \item{
\texttt{pMetric} -- }
  \end{itemize}
}%end item
\end{itemize}
}%end item
\end{itemize}
}
}
\subsection{\label{it.unisa.sesa.repominer.db.MetricMethodDAO}Class MetricMethodDAO}{
\hypertarget{it.unisa.sesa.repominer.db.MetricMethodDAO}{}\vskip .1in 
\subsubsection{Declaration}{
\small public class MetricMethodDAO
\\ {\bf  extends} java.lang.Object
\refdefined{java.lang.Object}}
\subsubsection{Constructor summary}{
\begin{verse}
\hyperlink{it.unisa.sesa.repominer.db.MetricMethodDAO()}{{\bf MetricMethodDAO()}} \\
\end{verse}
}
\subsubsection{Method summary}{
\begin{verse}
\hyperlink{it.unisa.sesa.repominer.db.MetricMethodDAO.getMetric(it.unisa.sesa.repominer.db.entities.Metric, it.unisa.sesa.repominer.db.entities.Method)}{{\bf getMetric(Metric, Method)}} This method return a metric of a method\\
\hyperlink{it.unisa.sesa.repominer.db.MetricMethodDAO.getMetricsOfMethod(it.unisa.sesa.repominer.db.entities.Method)}{{\bf getMetricsOfMethod(Method)}} This method return a list of metrics method\\
\end{verse}
}
\subsubsection{Constructors}{
\vskip -2em
\begin{itemize}
\item{ 
\index{MetricMethodDAO()}
\hypertarget{it.unisa.sesa.repominer.db.MetricMethodDAO()}{{\bf  MetricMethodDAO}\\}
\texttt{public\ {\bf  MetricMethodDAO}()
\label{it.unisa.sesa.repominer.db.MetricMethodDAO()}}%end signature
}%end item
\end{itemize}
}
\subsubsection{Methods}{
\vskip -2em
\begin{itemize}
\item{ 
\index{getMetric(Metric, Method)}
\hypertarget{it.unisa.sesa.repominer.db.MetricMethodDAO.getMetric(it.unisa.sesa.repominer.db.entities.Metric, it.unisa.sesa.repominer.db.entities.Method)}{{\bf  getMetric}\\}
\texttt{public entities.MetricMethod\ {\bf  getMetric}(\texttt{entities.Metric} {\bf  pMetric},
\texttt{entities.Method} {\bf  pMethod})
\label{it.unisa.sesa.repominer.db.MetricMethodDAO.getMetric(it.unisa.sesa.repominer.db.entities.Metric, it.unisa.sesa.repominer.db.entities.Method)}}%end signature
\begin{itemize}
\item{
{\bf  Description}

This method return a metric of a method
}
\item{
{\bf  Parameters}
  \begin{itemize}
   \item{
\texttt{pMetric} -- }
   \item{
\texttt{pMethod} -- }
  \end{itemize}
}%end item
\item{{\bf  Returns} -- 
A MetricMethod object 
}%end item
\end{itemize}
}%end item
\item{ 
\index{getMetricsOfMethod(Method)}
\hypertarget{it.unisa.sesa.repominer.db.MetricMethodDAO.getMetricsOfMethod(it.unisa.sesa.repominer.db.entities.Method)}{{\bf  getMetricsOfMethod}\\}
\texttt{public java.util.List\ {\bf  getMetricsOfMethod}(\texttt{entities.Method} {\bf  pMethod})
\label{it.unisa.sesa.repominer.db.MetricMethodDAO.getMetricsOfMethod(it.unisa.sesa.repominer.db.entities.Method)}}%end signature
\begin{itemize}
\item{
{\bf  Description}

This method return a list of metrics method
}
\item{
{\bf  Parameters}
  \begin{itemize}
   \item{
\texttt{pMethod} -- }
  \end{itemize}
}%end item
\item{{\bf  Returns} -- 
A list of MetricMethod objects 
}%end item
\end{itemize}
}%end item
\end{itemize}
}
}
\subsection{\label{it.unisa.sesa.repominer.db.PackageMetricDAO}Class PackageMetricDAO}{
\hypertarget{it.unisa.sesa.repominer.db.PackageMetricDAO}{}\vskip .1in 
\subsubsection{Declaration}{
\small public class PackageMetricDAO
\\ {\bf  extends} java.lang.Object
\refdefined{java.lang.Object}}
\subsubsection{Constructor summary}{
\begin{verse}
\hyperlink{it.unisa.sesa.repominer.db.PackageMetricDAO()}{{\bf PackageMetricDAO()}} \\
\end{verse}
}
\subsubsection{Method summary}{
\begin{verse}
\hyperlink{it.unisa.sesa.repominer.db.PackageMetricDAO.getMetric(it.unisa.sesa.repominer.db.entities.Metric, it.unisa.sesa.repominer.db.entities.SourceContainer)}{{\bf getMetric(Metric, SourceContainer)}} This method return a single package metric picked by metric and package id\\
\hyperlink{it.unisa.sesa.repominer.db.PackageMetricDAO.getMetricsOfPackage(it.unisa.sesa.repominer.db.entities.SourceContainer)}{{\bf getMetricsOfPackage(SourceContainer)}} This method returns a list of package metrics for package as parameter\\
\hyperlink{it.unisa.sesa.repominer.db.PackageMetricDAO.saveMetric(it.unisa.sesa.repominer.db.entities.PackageMetric)}{{\bf saveMetric(PackageMetric)}} This method saves or update a package metric into database\\
\end{verse}
}
\subsubsection{Constructors}{
\vskip -2em
\begin{itemize}
\item{ 
\index{PackageMetricDAO()}
\hypertarget{it.unisa.sesa.repominer.db.PackageMetricDAO()}{{\bf  PackageMetricDAO}\\}
\texttt{public\ {\bf  PackageMetricDAO}()
\label{it.unisa.sesa.repominer.db.PackageMetricDAO()}}%end signature
}%end item
\end{itemize}
}
\subsubsection{Methods}{
\vskip -2em
\begin{itemize}
\item{ 
\index{getMetric(Metric, SourceContainer)}
\hypertarget{it.unisa.sesa.repominer.db.PackageMetricDAO.getMetric(it.unisa.sesa.repominer.db.entities.Metric, it.unisa.sesa.repominer.db.entities.SourceContainer)}{{\bf  getMetric}\\}
\texttt{public entities.PackageMetric\ {\bf  getMetric}(\texttt{entities.Metric} {\bf  pMetric},
\texttt{entities.SourceContainer} {\bf  pSourceContainer})
\label{it.unisa.sesa.repominer.db.PackageMetricDAO.getMetric(it.unisa.sesa.repominer.db.entities.Metric, it.unisa.sesa.repominer.db.entities.SourceContainer)}}%end signature
\begin{itemize}
\item{
{\bf  Description}

This method return a single package metric picked by metric and package id
}
\item{
{\bf  Parameters}
  \begin{itemize}
   \item{
\texttt{pMetric} -- }
   \item{
\texttt{pSourceContainer} -- }
  \end{itemize}
}%end item
\item{{\bf  Returns} -- 
A PackageMetric object 
}%end item
\end{itemize}
}%end item
\item{ 
\index{getMetricsOfPackage(SourceContainer)}
\hypertarget{it.unisa.sesa.repominer.db.PackageMetricDAO.getMetricsOfPackage(it.unisa.sesa.repominer.db.entities.SourceContainer)}{{\bf  getMetricsOfPackage}\\}
\texttt{public java.util.List\ {\bf  getMetricsOfPackage}(\texttt{entities.SourceContainer} {\bf  pSourceContainer})
\label{it.unisa.sesa.repominer.db.PackageMetricDAO.getMetricsOfPackage(it.unisa.sesa.repominer.db.entities.SourceContainer)}}%end signature
\begin{itemize}
\item{
{\bf  Description}

This method returns a list of package metrics for package as parameter
}
\item{
{\bf  Parameters}
  \begin{itemize}
   \item{
\texttt{pSourceContainer} -- }
  \end{itemize}
}%end item
\item{{\bf  Returns} -- 
A list of PackageMetric objects 
}%end item
\end{itemize}
}%end item
\item{ 
\index{saveMetric(PackageMetric)}
\hypertarget{it.unisa.sesa.repominer.db.PackageMetricDAO.saveMetric(it.unisa.sesa.repominer.db.entities.PackageMetric)}{{\bf  saveMetric}\\}
\texttt{public void\ {\bf  saveMetric}(\texttt{entities.PackageMetric} {\bf  pPackageMetric})
\label{it.unisa.sesa.repominer.db.PackageMetricDAO.saveMetric(it.unisa.sesa.repominer.db.entities.PackageMetric)}}%end signature
\begin{itemize}
\item{
{\bf  Description}

This method saves or update a package metric into database
}
\item{
{\bf  Parameters}
  \begin{itemize}
   \item{
\texttt{pPackageMetric} -- }
  \end{itemize}
}%end item
\end{itemize}
}%end item
\end{itemize}
}
}
\subsection{\label{it.unisa.sesa.repominer.db.ProjectDAO}Class ProjectDAO}{
\hypertarget{it.unisa.sesa.repominer.db.ProjectDAO}{}\vskip .1in 
\subsubsection{Declaration}{
\small public class ProjectDAO
\\ {\bf  extends} java.lang.Object
\refdefined{java.lang.Object}}
\subsubsection{Constructor summary}{
\begin{verse}
\hyperlink{it.unisa.sesa.repominer.db.ProjectDAO()}{{\bf ProjectDAO()}} \\
\end{verse}
}
\subsubsection{Method summary}{
\begin{verse}
\hyperlink{it.unisa.sesa.repominer.db.ProjectDAO.getProject(java.lang.Integer)}{{\bf getProject(Integer)}} Returns project searched by id\\
\hyperlink{it.unisa.sesa.repominer.db.ProjectDAO.getProject(java.lang.String)}{{\bf getProject(String)}} Returns a project searched by name\\
\end{verse}
}
\subsubsection{Constructors}{
\vskip -2em
\begin{itemize}
\item{ 
\index{ProjectDAO()}
\hypertarget{it.unisa.sesa.repominer.db.ProjectDAO()}{{\bf  ProjectDAO}\\}
\texttt{public\ {\bf  ProjectDAO}()
\label{it.unisa.sesa.repominer.db.ProjectDAO()}}%end signature
}%end item
\end{itemize}
}
\subsubsection{Methods}{
\vskip -2em
\begin{itemize}
\item{ 
\index{getProject(Integer)}
\hypertarget{it.unisa.sesa.repominer.db.ProjectDAO.getProject(java.lang.Integer)}{{\bf  getProject}\\}
\texttt{public entities.Project\ {\bf  getProject}(\texttt{java.lang.Integer} {\bf  pId})
\label{it.unisa.sesa.repominer.db.ProjectDAO.getProject(java.lang.Integer)}}%end signature
\begin{itemize}
\item{
{\bf  Description}

Returns project searched by id
}
\item{
{\bf  Parameters}
  \begin{itemize}
   \item{
\texttt{pId} -- }
  \end{itemize}
}%end item
\item{{\bf  Returns} -- 
A Project class 
}%end item
\end{itemize}
}%end item
\item{ 
\index{getProject(String)}
\hypertarget{it.unisa.sesa.repominer.db.ProjectDAO.getProject(java.lang.String)}{{\bf  getProject}\\}
\texttt{public entities.Project\ {\bf  getProject}(\texttt{java.lang.String} {\bf  pName})
\label{it.unisa.sesa.repominer.db.ProjectDAO.getProject(java.lang.String)}}%end signature
\begin{itemize}
\item{
{\bf  Description}

Returns a project searched by name
}
\item{
{\bf  Parameters}
  \begin{itemize}
   \item{
\texttt{pName} -- }
  \end{itemize}
}%end item
\item{{\bf  Returns} -- 
A Project class 
}%end item
\end{itemize}
}%end item
\end{itemize}
}
}
\subsection{\label{it.unisa.sesa.repominer.db.ProjectMetricDAO}Class ProjectMetricDAO}{
\hypertarget{it.unisa.sesa.repominer.db.ProjectMetricDAO}{}\vskip .1in 
\subsubsection{Declaration}{
\small public class ProjectMetricDAO
\\ {\bf  extends} java.lang.Object
\refdefined{java.lang.Object}}
\subsubsection{Constructor summary}{
\begin{verse}
\hyperlink{it.unisa.sesa.repominer.db.ProjectMetricDAO()}{{\bf ProjectMetricDAO()}} \\
\end{verse}
}
\subsubsection{Method summary}{
\begin{verse}
\hyperlink{it.unisa.sesa.repominer.db.ProjectMetricDAO.getMetric(it.unisa.sesa.repominer.db.entities.Metric, it.unisa.sesa.repominer.db.entities.Project)}{{\bf getMetric(Metric, Project)}} This method returns a single project metric picked by metric and project id\\
\hyperlink{it.unisa.sesa.repominer.db.ProjectMetricDAO.getMetricsOfProject(it.unisa.sesa.repominer.db.entities.Project)}{{\bf getMetricsOfProject(Project)}} This method returns a list of project metrics for a project passed as parameter\\
\hyperlink{it.unisa.sesa.repominer.db.ProjectMetricDAO.saveMetric(it.unisa.sesa.repominer.db.entities.ProjectMetric)}{{\bf saveMetric(ProjectMetric)}} This method save a project metric into database\\
\end{verse}
}
\subsubsection{Constructors}{
\vskip -2em
\begin{itemize}
\item{ 
\index{ProjectMetricDAO()}
\hypertarget{it.unisa.sesa.repominer.db.ProjectMetricDAO()}{{\bf  ProjectMetricDAO}\\}
\texttt{public\ {\bf  ProjectMetricDAO}()
\label{it.unisa.sesa.repominer.db.ProjectMetricDAO()}}%end signature
}%end item
\end{itemize}
}
\subsubsection{Methods}{
\vskip -2em
\begin{itemize}
\item{ 
\index{getMetric(Metric, Project)}
\hypertarget{it.unisa.sesa.repominer.db.ProjectMetricDAO.getMetric(it.unisa.sesa.repominer.db.entities.Metric, it.unisa.sesa.repominer.db.entities.Project)}{{\bf  getMetric}\\}
\texttt{public entities.ProjectMetric\ {\bf  getMetric}(\texttt{entities.Metric} {\bf  pMetric},
\texttt{entities.Project} {\bf  pProject})
\label{it.unisa.sesa.repominer.db.ProjectMetricDAO.getMetric(it.unisa.sesa.repominer.db.entities.Metric, it.unisa.sesa.repominer.db.entities.Project)}}%end signature
\begin{itemize}
\item{
{\bf  Description}

This method returns a single project metric picked by metric and project id
}
\item{
{\bf  Parameters}
  \begin{itemize}
   \item{
\texttt{pMetric} -- }
   \item{
\texttt{pProject} -- }
  \end{itemize}
}%end item
\item{{\bf  Returns} -- 
 
}%end item
\end{itemize}
}%end item
\item{ 
\index{getMetricsOfProject(Project)}
\hypertarget{it.unisa.sesa.repominer.db.ProjectMetricDAO.getMetricsOfProject(it.unisa.sesa.repominer.db.entities.Project)}{{\bf  getMetricsOfProject}\\}
\texttt{public java.util.List\ {\bf  getMetricsOfProject}(\texttt{entities.Project} {\bf  pProject})
\label{it.unisa.sesa.repominer.db.ProjectMetricDAO.getMetricsOfProject(it.unisa.sesa.repominer.db.entities.Project)}}%end signature
\begin{itemize}
\item{
{\bf  Description}

This method returns a list of project metrics for a project passed as parameter
}
\item{
{\bf  Parameters}
  \begin{itemize}
   \item{
\texttt{pProject} -- }
  \end{itemize}
}%end item
\item{{\bf  Returns} -- 
A list of ProjectMetric objects 
}%end item
\end{itemize}
}%end item
\item{ 
\index{saveMetric(ProjectMetric)}
\hypertarget{it.unisa.sesa.repominer.db.ProjectMetricDAO.saveMetric(it.unisa.sesa.repominer.db.entities.ProjectMetric)}{{\bf  saveMetric}\\}
\texttt{public void\ {\bf  saveMetric}(\texttt{entities.ProjectMetric} {\bf  pProjectMetric})
\label{it.unisa.sesa.repominer.db.ProjectMetricDAO.saveMetric(it.unisa.sesa.repominer.db.entities.ProjectMetric)}}%end signature
\begin{itemize}
\item{
{\bf  Description}

This method save a project metric into database
}
\item{
{\bf  Parameters}
  \begin{itemize}
   \item{
\texttt{pProjectMetric} -- }
  \end{itemize}
}%end item
\end{itemize}
}%end item
\end{itemize}
}
}
\subsection{\label{it.unisa.sesa.repominer.db.SourceContainerDAO}Class SourceContainerDAO}{
\hypertarget{it.unisa.sesa.repominer.db.SourceContainerDAO}{}\vskip .1in 
\subsubsection{Declaration}{
\small public class SourceContainerDAO
\\ {\bf  extends} java.lang.Object
\refdefined{java.lang.Object}}
\subsubsection{Constructor summary}{
\begin{verse}
\hyperlink{it.unisa.sesa.repominer.db.SourceContainerDAO()}{{\bf SourceContainerDAO()}} \\
\end{verse}
}
\subsubsection{Method summary}{
\begin{verse}
\hyperlink{it.unisa.sesa.repominer.db.SourceContainerDAO.getPackages(it.unisa.sesa.repominer.db.entities.Project)}{{\bf getPackages(Project)}} Returns all packages contained in project passed as parameter\\
\hyperlink{it.unisa.sesa.repominer.db.SourceContainerDAO.getPackagesById(java.lang.Integer)}{{\bf getPackagesById(Integer)}} Returns a single package picked by id\\
\end{verse}
}
\subsubsection{Constructors}{
\vskip -2em
\begin{itemize}
\item{ 
\index{SourceContainerDAO()}
\hypertarget{it.unisa.sesa.repominer.db.SourceContainerDAO()}{{\bf  SourceContainerDAO}\\}
\texttt{public\ {\bf  SourceContainerDAO}()
\label{it.unisa.sesa.repominer.db.SourceContainerDAO()}}%end signature
}%end item
\end{itemize}
}
\subsubsection{Methods}{
\vskip -2em
\begin{itemize}
\item{ 
\index{getPackages(Project)}
\hypertarget{it.unisa.sesa.repominer.db.SourceContainerDAO.getPackages(it.unisa.sesa.repominer.db.entities.Project)}{{\bf  getPackages}\\}
\texttt{public java.util.List\ {\bf  getPackages}(\texttt{entities.Project} {\bf  pProject})
\label{it.unisa.sesa.repominer.db.SourceContainerDAO.getPackages(it.unisa.sesa.repominer.db.entities.Project)}}%end signature
\begin{itemize}
\item{
{\bf  Description}

Returns all packages contained in project passed as parameter
}
\item{
{\bf  Parameters}
  \begin{itemize}
   \item{
\texttt{pProject} -- }
  \end{itemize}
}%end item
\item{{\bf  Returns} -- 
List of SourceContainer objects 
}%end item
\end{itemize}
}%end item
\item{ 
\index{getPackagesById(Integer)}
\hypertarget{it.unisa.sesa.repominer.db.SourceContainerDAO.getPackagesById(java.lang.Integer)}{{\bf  getPackagesById}\\}
\texttt{public entities.SourceContainer\ {\bf  getPackagesById}(\texttt{java.lang.Integer} {\bf  pId})
\label{it.unisa.sesa.repominer.db.SourceContainerDAO.getPackagesById(java.lang.Integer)}}%end signature
\begin{itemize}
\item{
{\bf  Description}

Returns a single package picked by id
}
\item{
{\bf  Parameters}
  \begin{itemize}
   \item{
\texttt{pId} -- }
  \end{itemize}
}%end item
\item{{\bf  Returns} -- 
A SourceContainer object 
}%end item
\end{itemize}
}%end item
\end{itemize}
}
}
\subsection{\label{it.unisa.sesa.repominer.db.TypeDAO}Class TypeDAO}{
\hypertarget{it.unisa.sesa.repominer.db.TypeDAO}{}\vskip .1in 
\subsubsection{Declaration}{
\small public class TypeDAO
\\ {\bf  extends} java.lang.Object
\refdefined{java.lang.Object}}
\subsubsection{Constructor summary}{
\begin{verse}
\hyperlink{it.unisa.sesa.repominer.db.TypeDAO()}{{\bf TypeDAO()}} \\
\end{verse}
}
\subsubsection{Method summary}{
\begin{verse}
\hyperlink{it.unisa.sesa.repominer.db.TypeDAO.getClassById(java.lang.Integer)}{{\bf getClassById(Integer)}} Returns a single class picked by id\\
\hyperlink{it.unisa.sesa.repominer.db.TypeDAO.getClassesByPackage(it.unisa.sesa.repominer.db.entities.SourceContainer)}{{\bf getClassesByPackage(SourceContainer)}} Returns all classes contained in the package passed as parameter\\
\hyperlink{it.unisa.sesa.repominer.db.TypeDAO.getSystemNumberOfTypes(it.unisa.sesa.repominer.db.entities.Project)}{{\bf getSystemNumberOfTypes(Project)}} Return the overall number of types in project passed as parameter\\
\end{verse}
}
\subsubsection{Constructors}{
\vskip -2em
\begin{itemize}
\item{ 
\index{TypeDAO()}
\hypertarget{it.unisa.sesa.repominer.db.TypeDAO()}{{\bf  TypeDAO}\\}
\texttt{public\ {\bf  TypeDAO}()
\label{it.unisa.sesa.repominer.db.TypeDAO()}}%end signature
}%end item
\end{itemize}
}
\subsubsection{Methods}{
\vskip -2em
\begin{itemize}
\item{ 
\index{getClassById(Integer)}
\hypertarget{it.unisa.sesa.repominer.db.TypeDAO.getClassById(java.lang.Integer)}{{\bf  getClassById}\\}
\texttt{public entities.Type\ {\bf  getClassById}(\texttt{java.lang.Integer} {\bf  pId})
\label{it.unisa.sesa.repominer.db.TypeDAO.getClassById(java.lang.Integer)}}%end signature
\begin{itemize}
\item{
{\bf  Description}

Returns a single class picked by id
}
\item{
{\bf  Parameters}
  \begin{itemize}
   \item{
\texttt{pId} -- }
  \end{itemize}
}%end item
\item{{\bf  Returns} -- 
A Type object 
}%end item
\end{itemize}
}%end item
\item{ 
\index{getClassesByPackage(SourceContainer)}
\hypertarget{it.unisa.sesa.repominer.db.TypeDAO.getClassesByPackage(it.unisa.sesa.repominer.db.entities.SourceContainer)}{{\bf  getClassesByPackage}\\}
\texttt{public java.util.List\ {\bf  getClassesByPackage}(\texttt{entities.SourceContainer} {\bf  pSourceContainer})
\label{it.unisa.sesa.repominer.db.TypeDAO.getClassesByPackage(it.unisa.sesa.repominer.db.entities.SourceContainer)}}%end signature
\begin{itemize}
\item{
{\bf  Description}

Returns all classes contained in the package passed as parameter
}
\item{
{\bf  Parameters}
  \begin{itemize}
   \item{
\texttt{pSourceContainer} -- }
  \end{itemize}
}%end item
\item{{\bf  Returns} -- 
A list of Type objects 
}%end item
\end{itemize}
}%end item
\item{ 
\index{getSystemNumberOfTypes(Project)}
\hypertarget{it.unisa.sesa.repominer.db.TypeDAO.getSystemNumberOfTypes(it.unisa.sesa.repominer.db.entities.Project)}{{\bf  getSystemNumberOfTypes}\\}
\texttt{public int\ {\bf  getSystemNumberOfTypes}(\texttt{entities.Project} {\bf  pProject})
\label{it.unisa.sesa.repominer.db.TypeDAO.getSystemNumberOfTypes(it.unisa.sesa.repominer.db.entities.Project)}}%end signature
\begin{itemize}
\item{
{\bf  Description}

Return the overall number of types in project passed as parameter
}
\item{
{\bf  Parameters}
  \begin{itemize}
   \item{
\texttt{pProject} -- selected project}
  \end{itemize}
}%end item
\item{{\bf  Returns} -- 
number of types in project 
}%end item
\end{itemize}
}%end item
\end{itemize}
}
}
}
\section{Package it.unisa.sesa.repominer.db.entities}{
\label{it.unisa.sesa.repominer.db.entities}\hypertarget{it.unisa.sesa.repominer.db.entities}{}
\hskip -.05in
\hbox to \hsize{\textit{ Package Contents\hfil Page}}
\vskip .13in
\hbox{{\bf  Classes}}
\entityintro{CatchedException}{it.unisa.sesa.repominer.db.entities.CatchedException}{}
\entityintro{Change}{it.unisa.sesa.repominer.db.entities.Change}{}
\entityintro{ChangeForCommit}{it.unisa.sesa.repominer.db.entities.ChangeForCommit}{}
\entityintro{ClassInvocation}{it.unisa.sesa.repominer.db.entities.ClassInvocation}{}
\entityintro{Import}{it.unisa.sesa.repominer.db.entities.Import}{}
\entityintro{Imports}{it.unisa.sesa.repominer.db.entities.Imports}{}
\entityintro{Issue}{it.unisa.sesa.repominer.db.entities.Issue}{}
\entityintro{IssueAttachment}{it.unisa.sesa.repominer.db.entities.IssueAttachment}{}
\entityintro{IssueComment}{it.unisa.sesa.repominer.db.entities.IssueComment}{}
\entityintro{Method}{it.unisa.sesa.repominer.db.entities.Method}{}
\entityintro{MethodComment}{it.unisa.sesa.repominer.db.entities.MethodComment}{}
\entityintro{MethodInvocation}{it.unisa.sesa.repominer.db.entities.MethodInvocation}{}
\entityintro{MethodsChangeInCommit}{it.unisa.sesa.repominer.db.entities.MethodsChangeInCommit}{}
\entityintro{Metric}{it.unisa.sesa.repominer.db.entities.Metric}{}
\entityintro{MetricMethod}{it.unisa.sesa.repominer.db.entities.MetricMethod}{}
\entityintro{PackageMetric}{it.unisa.sesa.repominer.db.entities.PackageMetric}{}
\entityintro{Project}{it.unisa.sesa.repominer.db.entities.Project}{}
\entityintro{ProjectMetric}{it.unisa.sesa.repominer.db.entities.ProjectMetric}{}
\entityintro{Review}{it.unisa.sesa.repominer.db.entities.Review}{}
\entityintro{SourceContainer}{it.unisa.sesa.repominer.db.entities.SourceContainer}{}
\entityintro{Type}{it.unisa.sesa.repominer.db.entities.Type}{}
\entityintro{TypeMetric}{it.unisa.sesa.repominer.db.entities.TypeMetric}{}
\vskip .1in
\vskip .1in
\subsection{\label{it.unisa.sesa.repominer.db.entities.CatchedException}Class CatchedException}{
\hypertarget{it.unisa.sesa.repominer.db.entities.CatchedException}{}\vskip .1in 
\subsubsection{Declaration}{
\small public class CatchedException
\\ {\bf  extends} java.lang.Object
\refdefined{java.lang.Object}}
\subsubsection{Constructor summary}{
\begin{verse}
\hyperlink{it.unisa.sesa.repominer.db.entities.CatchedException()}{{\bf CatchedException()}} \\
\hyperlink{it.unisa.sesa.repominer.db.entities.CatchedException(java.lang.Integer, java.lang.Integer)}{{\bf CatchedException(Integer, Integer)}} \\
\end{verse}
}
\subsubsection{Method summary}{
\begin{verse}
\hyperlink{it.unisa.sesa.repominer.db.entities.CatchedException.equals(java.lang.Object)}{{\bf equals(Object)}} \\
\hyperlink{it.unisa.sesa.repominer.db.entities.CatchedException.getExceptionTypeId()}{{\bf getExceptionTypeId()}} \\
\hyperlink{it.unisa.sesa.repominer.db.entities.CatchedException.getMethodId()}{{\bf getMethodId()}} \\
\hyperlink{it.unisa.sesa.repominer.db.entities.CatchedException.hashCode()}{{\bf hashCode()}} \\
\hyperlink{it.unisa.sesa.repominer.db.entities.CatchedException.setExceptionTypeId(java.lang.Integer)}{{\bf setExceptionTypeId(Integer)}} \\
\hyperlink{it.unisa.sesa.repominer.db.entities.CatchedException.setMethodId(java.lang.Integer)}{{\bf setMethodId(Integer)}} \\
\hyperlink{it.unisa.sesa.repominer.db.entities.CatchedException.toString()}{{\bf toString()}} \\
\end{verse}
}
\subsubsection{Constructors}{
\vskip -2em
\begin{itemize}
\item{ 
\index{CatchedException()}
\hypertarget{it.unisa.sesa.repominer.db.entities.CatchedException()}{{\bf  CatchedException}\\}
\texttt{public\ {\bf  CatchedException}()
\label{it.unisa.sesa.repominer.db.entities.CatchedException()}}%end signature
}%end item
\item{ 
\index{CatchedException(Integer, Integer)}
\hypertarget{it.unisa.sesa.repominer.db.entities.CatchedException(java.lang.Integer, java.lang.Integer)}{{\bf  CatchedException}\\}
\texttt{public\ {\bf  CatchedException}(\texttt{java.lang.Integer} {\bf  methodId},
\texttt{java.lang.Integer} {\bf  exceptionTypeId})
\label{it.unisa.sesa.repominer.db.entities.CatchedException(java.lang.Integer, java.lang.Integer)}}%end signature
}%end item
\end{itemize}
}
\subsubsection{Methods}{
\vskip -2em
\begin{itemize}
\item{ 
\index{equals(Object)}
\hypertarget{it.unisa.sesa.repominer.db.entities.CatchedException.equals(java.lang.Object)}{{\bf  equals}\\}
\texttt{public boolean\ {\bf  equals}(\texttt{java.lang.Object} {\bf  arg0})
\label{it.unisa.sesa.repominer.db.entities.CatchedException.equals(java.lang.Object)}}%end signature
}%end item
\item{ 
\index{getExceptionTypeId()}
\hypertarget{it.unisa.sesa.repominer.db.entities.CatchedException.getExceptionTypeId()}{{\bf  getExceptionTypeId}\\}
\texttt{public java.lang.Integer\ {\bf  getExceptionTypeId}()
\label{it.unisa.sesa.repominer.db.entities.CatchedException.getExceptionTypeId()}}%end signature
}%end item
\item{ 
\index{getMethodId()}
\hypertarget{it.unisa.sesa.repominer.db.entities.CatchedException.getMethodId()}{{\bf  getMethodId}\\}
\texttt{public java.lang.Integer\ {\bf  getMethodId}()
\label{it.unisa.sesa.repominer.db.entities.CatchedException.getMethodId()}}%end signature
}%end item
\item{ 
\index{hashCode()}
\hypertarget{it.unisa.sesa.repominer.db.entities.CatchedException.hashCode()}{{\bf  hashCode}\\}
\texttt{public native int\ {\bf  hashCode}()
\label{it.unisa.sesa.repominer.db.entities.CatchedException.hashCode()}}%end signature
}%end item
\item{ 
\index{setExceptionTypeId(Integer)}
\hypertarget{it.unisa.sesa.repominer.db.entities.CatchedException.setExceptionTypeId(java.lang.Integer)}{{\bf  setExceptionTypeId}\\}
\texttt{public void\ {\bf  setExceptionTypeId}(\texttt{java.lang.Integer} {\bf  exceptionTypeId})
\label{it.unisa.sesa.repominer.db.entities.CatchedException.setExceptionTypeId(java.lang.Integer)}}%end signature
}%end item
\item{ 
\index{setMethodId(Integer)}
\hypertarget{it.unisa.sesa.repominer.db.entities.CatchedException.setMethodId(java.lang.Integer)}{{\bf  setMethodId}\\}
\texttt{public void\ {\bf  setMethodId}(\texttt{java.lang.Integer} {\bf  methodId})
\label{it.unisa.sesa.repominer.db.entities.CatchedException.setMethodId(java.lang.Integer)}}%end signature
}%end item
\item{ 
\index{toString()}
\hypertarget{it.unisa.sesa.repominer.db.entities.CatchedException.toString()}{{\bf  toString}\\}
\texttt{public java.lang.String\ {\bf  toString}()
\label{it.unisa.sesa.repominer.db.entities.CatchedException.toString()}}%end signature
}%end item
\end{itemize}
}
}
\subsection{\label{it.unisa.sesa.repominer.db.entities.Change}Class Change}{
\hypertarget{it.unisa.sesa.repominer.db.entities.Change}{}\vskip .1in 
\subsubsection{Declaration}{
\small public class Change
\\ {\bf  extends} java.lang.Object
\refdefined{java.lang.Object}}
\subsubsection{Constructor summary}{
\begin{verse}
\hyperlink{it.unisa.sesa.repominer.db.entities.Change()}{{\bf Change()}} \\
\hyperlink{it.unisa.sesa.repominer.db.entities.Change(java.lang.Integer, java.lang.String, java.util.Date, java.lang.String, java.lang.String, java.lang.String, java.lang.Integer)}{{\bf Change(Integer, String, Date, String, String, String, Integer)}} \\
\end{verse}
}
\subsubsection{Method summary}{
\begin{verse}
\hyperlink{it.unisa.sesa.repominer.db.entities.Change.equals(java.lang.Object)}{{\bf equals(Object)}} \\
\hyperlink{it.unisa.sesa.repominer.db.entities.Change.getCommitDate()}{{\bf getCommitDate()}} \\
\hyperlink{it.unisa.sesa.repominer.db.entities.Change.getDevId()}{{\bf getDevId()}} \\
\hyperlink{it.unisa.sesa.repominer.db.entities.Change.getDevMail()}{{\bf getDevMail()}} \\
\hyperlink{it.unisa.sesa.repominer.db.entities.Change.getHash()}{{\bf getHash()}} \\
\hyperlink{it.unisa.sesa.repominer.db.entities.Change.getId()}{{\bf getId()}} \\
\hyperlink{it.unisa.sesa.repominer.db.entities.Change.getMessage()}{{\bf getMessage()}} \\
\hyperlink{it.unisa.sesa.repominer.db.entities.Change.getProjectId()}{{\bf getProjectId()}} \\
\hyperlink{it.unisa.sesa.repominer.db.entities.Change.hashCode()}{{\bf hashCode()}} \\
\hyperlink{it.unisa.sesa.repominer.db.entities.Change.setCommitDate(java.util.Date)}{{\bf setCommitDate(Date)}} \\
\hyperlink{it.unisa.sesa.repominer.db.entities.Change.setDevId(java.lang.String)}{{\bf setDevId(String)}} \\
\hyperlink{it.unisa.sesa.repominer.db.entities.Change.setDevMail(java.lang.String)}{{\bf setDevMail(String)}} \\
\hyperlink{it.unisa.sesa.repominer.db.entities.Change.setHash(java.lang.String)}{{\bf setHash(String)}} \\
\hyperlink{it.unisa.sesa.repominer.db.entities.Change.setId(java.lang.Integer)}{{\bf setId(Integer)}} \\
\hyperlink{it.unisa.sesa.repominer.db.entities.Change.setMessage(java.lang.String)}{{\bf setMessage(String)}} \\
\hyperlink{it.unisa.sesa.repominer.db.entities.Change.setProjectId(java.lang.Integer)}{{\bf setProjectId(Integer)}} \\
\hyperlink{it.unisa.sesa.repominer.db.entities.Change.toString()}{{\bf toString()}} \\
\end{verse}
}
\subsubsection{Constructors}{
\vskip -2em
\begin{itemize}
\item{ 
\index{Change()}
\hypertarget{it.unisa.sesa.repominer.db.entities.Change()}{{\bf  Change}\\}
\texttt{public\ {\bf  Change}()
\label{it.unisa.sesa.repominer.db.entities.Change()}}%end signature
}%end item
\item{ 
\index{Change(Integer, String, Date, String, String, String, Integer)}
\hypertarget{it.unisa.sesa.repominer.db.entities.Change(java.lang.Integer, java.lang.String, java.util.Date, java.lang.String, java.lang.String, java.lang.String, java.lang.Integer)}{{\bf  Change}\\}
\texttt{public\ {\bf  Change}(\texttt{java.lang.Integer} {\bf  id},
\texttt{java.lang.String} {\bf  hash},
\texttt{java.util.Date} {\bf  commitDate},
\texttt{java.lang.String} {\bf  devMail},
\texttt{java.lang.String} {\bf  devId},
\texttt{java.lang.String} {\bf  message},
\texttt{java.lang.Integer} {\bf  projectId})
\label{it.unisa.sesa.repominer.db.entities.Change(java.lang.Integer, java.lang.String, java.util.Date, java.lang.String, java.lang.String, java.lang.String, java.lang.Integer)}}%end signature
}%end item
\end{itemize}
}
\subsubsection{Methods}{
\vskip -2em
\begin{itemize}
\item{ 
\index{equals(Object)}
\hypertarget{it.unisa.sesa.repominer.db.entities.Change.equals(java.lang.Object)}{{\bf  equals}\\}
\texttt{public boolean\ {\bf  equals}(\texttt{java.lang.Object} {\bf  arg0})
\label{it.unisa.sesa.repominer.db.entities.Change.equals(java.lang.Object)}}%end signature
}%end item
\item{ 
\index{getCommitDate()}
\hypertarget{it.unisa.sesa.repominer.db.entities.Change.getCommitDate()}{{\bf  getCommitDate}\\}
\texttt{public java.util.Date\ {\bf  getCommitDate}()
\label{it.unisa.sesa.repominer.db.entities.Change.getCommitDate()}}%end signature
}%end item
\item{ 
\index{getDevId()}
\hypertarget{it.unisa.sesa.repominer.db.entities.Change.getDevId()}{{\bf  getDevId}\\}
\texttt{public java.lang.String\ {\bf  getDevId}()
\label{it.unisa.sesa.repominer.db.entities.Change.getDevId()}}%end signature
}%end item
\item{ 
\index{getDevMail()}
\hypertarget{it.unisa.sesa.repominer.db.entities.Change.getDevMail()}{{\bf  getDevMail}\\}
\texttt{public java.lang.String\ {\bf  getDevMail}()
\label{it.unisa.sesa.repominer.db.entities.Change.getDevMail()}}%end signature
}%end item
\item{ 
\index{getHash()}
\hypertarget{it.unisa.sesa.repominer.db.entities.Change.getHash()}{{\bf  getHash}\\}
\texttt{public java.lang.String\ {\bf  getHash}()
\label{it.unisa.sesa.repominer.db.entities.Change.getHash()}}%end signature
}%end item
\item{ 
\index{getId()}
\hypertarget{it.unisa.sesa.repominer.db.entities.Change.getId()}{{\bf  getId}\\}
\texttt{public java.lang.Integer\ {\bf  getId}()
\label{it.unisa.sesa.repominer.db.entities.Change.getId()}}%end signature
}%end item
\item{ 
\index{getMessage()}
\hypertarget{it.unisa.sesa.repominer.db.entities.Change.getMessage()}{{\bf  getMessage}\\}
\texttt{public java.lang.String\ {\bf  getMessage}()
\label{it.unisa.sesa.repominer.db.entities.Change.getMessage()}}%end signature
}%end item
\item{ 
\index{getProjectId()}
\hypertarget{it.unisa.sesa.repominer.db.entities.Change.getProjectId()}{{\bf  getProjectId}\\}
\texttt{public java.lang.Integer\ {\bf  getProjectId}()
\label{it.unisa.sesa.repominer.db.entities.Change.getProjectId()}}%end signature
}%end item
\item{ 
\index{hashCode()}
\hypertarget{it.unisa.sesa.repominer.db.entities.Change.hashCode()}{{\bf  hashCode}\\}
\texttt{public native int\ {\bf  hashCode}()
\label{it.unisa.sesa.repominer.db.entities.Change.hashCode()}}%end signature
}%end item
\item{ 
\index{setCommitDate(Date)}
\hypertarget{it.unisa.sesa.repominer.db.entities.Change.setCommitDate(java.util.Date)}{{\bf  setCommitDate}\\}
\texttt{public void\ {\bf  setCommitDate}(\texttt{java.util.Date} {\bf  commitDate})
\label{it.unisa.sesa.repominer.db.entities.Change.setCommitDate(java.util.Date)}}%end signature
}%end item
\item{ 
\index{setDevId(String)}
\hypertarget{it.unisa.sesa.repominer.db.entities.Change.setDevId(java.lang.String)}{{\bf  setDevId}\\}
\texttt{public void\ {\bf  setDevId}(\texttt{java.lang.String} {\bf  devId})
\label{it.unisa.sesa.repominer.db.entities.Change.setDevId(java.lang.String)}}%end signature
}%end item
\item{ 
\index{setDevMail(String)}
\hypertarget{it.unisa.sesa.repominer.db.entities.Change.setDevMail(java.lang.String)}{{\bf  setDevMail}\\}
\texttt{public void\ {\bf  setDevMail}(\texttt{java.lang.String} {\bf  devMail})
\label{it.unisa.sesa.repominer.db.entities.Change.setDevMail(java.lang.String)}}%end signature
}%end item
\item{ 
\index{setHash(String)}
\hypertarget{it.unisa.sesa.repominer.db.entities.Change.setHash(java.lang.String)}{{\bf  setHash}\\}
\texttt{public void\ {\bf  setHash}(\texttt{java.lang.String} {\bf  hash})
\label{it.unisa.sesa.repominer.db.entities.Change.setHash(java.lang.String)}}%end signature
}%end item
\item{ 
\index{setId(Integer)}
\hypertarget{it.unisa.sesa.repominer.db.entities.Change.setId(java.lang.Integer)}{{\bf  setId}\\}
\texttt{public void\ {\bf  setId}(\texttt{java.lang.Integer} {\bf  id})
\label{it.unisa.sesa.repominer.db.entities.Change.setId(java.lang.Integer)}}%end signature
}%end item
\item{ 
\index{setMessage(String)}
\hypertarget{it.unisa.sesa.repominer.db.entities.Change.setMessage(java.lang.String)}{{\bf  setMessage}\\}
\texttt{public void\ {\bf  setMessage}(\texttt{java.lang.String} {\bf  message})
\label{it.unisa.sesa.repominer.db.entities.Change.setMessage(java.lang.String)}}%end signature
}%end item
\item{ 
\index{setProjectId(Integer)}
\hypertarget{it.unisa.sesa.repominer.db.entities.Change.setProjectId(java.lang.Integer)}{{\bf  setProjectId}\\}
\texttt{public void\ {\bf  setProjectId}(\texttt{java.lang.Integer} {\bf  projectId})
\label{it.unisa.sesa.repominer.db.entities.Change.setProjectId(java.lang.Integer)}}%end signature
}%end item
\item{ 
\index{toString()}
\hypertarget{it.unisa.sesa.repominer.db.entities.Change.toString()}{{\bf  toString}\\}
\texttt{public java.lang.String\ {\bf  toString}()
\label{it.unisa.sesa.repominer.db.entities.Change.toString()}}%end signature
}%end item
\end{itemize}
}
}
\subsection{\label{it.unisa.sesa.repominer.db.entities.ChangeForCommit}Class ChangeForCommit}{
\hypertarget{it.unisa.sesa.repominer.db.entities.ChangeForCommit}{}\vskip .1in 
\subsubsection{Declaration}{
\small public class ChangeForCommit
\\ {\bf  extends} java.lang.Object
\refdefined{java.lang.Object}}
\subsubsection{Constructor summary}{
\begin{verse}
\hyperlink{it.unisa.sesa.repominer.db.entities.ChangeForCommit()}{{\bf ChangeForCommit()}} \\
\hyperlink{it.unisa.sesa.repominer.db.entities.ChangeForCommit(java.lang.Integer, java.lang.Integer, java.lang.String, java.lang.Integer, java.lang.Integer)}{{\bf ChangeForCommit(Integer, Integer, String, Integer, Integer)}} \\
\end{verse}
}
\subsubsection{Method summary}{
\begin{verse}
\hyperlink{it.unisa.sesa.repominer.db.entities.ChangeForCommit.equals(java.lang.Object)}{{\bf equals(Object)}} \\
\hyperlink{it.unisa.sesa.repominer.db.entities.ChangeForCommit.getChangeHashId()}{{\bf getChangeHashId()}} \\
\hyperlink{it.unisa.sesa.repominer.db.entities.ChangeForCommit.getDeletions()}{{\bf getDeletions()}} \\
\hyperlink{it.unisa.sesa.repominer.db.entities.ChangeForCommit.getId()}{{\bf getId()}} \\
\hyperlink{it.unisa.sesa.repominer.db.entities.ChangeForCommit.getInsertions()}{{\bf getInsertions()}} \\
\hyperlink{it.unisa.sesa.repominer.db.entities.ChangeForCommit.getModifiedFile()}{{\bf getModifiedFile()}} \\
\hyperlink{it.unisa.sesa.repominer.db.entities.ChangeForCommit.hashCode()}{{\bf hashCode()}} \\
\hyperlink{it.unisa.sesa.repominer.db.entities.ChangeForCommit.setChangeHashId(java.lang.Integer)}{{\bf setChangeHashId(Integer)}} \\
\hyperlink{it.unisa.sesa.repominer.db.entities.ChangeForCommit.setDeletions(java.lang.Integer)}{{\bf setDeletions(Integer)}} \\
\hyperlink{it.unisa.sesa.repominer.db.entities.ChangeForCommit.setId(java.lang.Integer)}{{\bf setId(Integer)}} \\
\hyperlink{it.unisa.sesa.repominer.db.entities.ChangeForCommit.setInsertions(java.lang.Integer)}{{\bf setInsertions(Integer)}} \\
\hyperlink{it.unisa.sesa.repominer.db.entities.ChangeForCommit.setModifiedFile(java.lang.String)}{{\bf setModifiedFile(String)}} \\
\hyperlink{it.unisa.sesa.repominer.db.entities.ChangeForCommit.toString()}{{\bf toString()}} \\
\end{verse}
}
\subsubsection{Constructors}{
\vskip -2em
\begin{itemize}
\item{ 
\index{ChangeForCommit()}
\hypertarget{it.unisa.sesa.repominer.db.entities.ChangeForCommit()}{{\bf  ChangeForCommit}\\}
\texttt{public\ {\bf  ChangeForCommit}()
\label{it.unisa.sesa.repominer.db.entities.ChangeForCommit()}}%end signature
}%end item
\item{ 
\index{ChangeForCommit(Integer, Integer, String, Integer, Integer)}
\hypertarget{it.unisa.sesa.repominer.db.entities.ChangeForCommit(java.lang.Integer, java.lang.Integer, java.lang.String, java.lang.Integer, java.lang.Integer)}{{\bf  ChangeForCommit}\\}
\texttt{public\ {\bf  ChangeForCommit}(\texttt{java.lang.Integer} {\bf  id},
\texttt{java.lang.Integer} {\bf  changeHashId},
\texttt{java.lang.String} {\bf  modifiedFile},
\texttt{java.lang.Integer} {\bf  insertions},
\texttt{java.lang.Integer} {\bf  deletions})
\label{it.unisa.sesa.repominer.db.entities.ChangeForCommit(java.lang.Integer, java.lang.Integer, java.lang.String, java.lang.Integer, java.lang.Integer)}}%end signature
}%end item
\end{itemize}
}
\subsubsection{Methods}{
\vskip -2em
\begin{itemize}
\item{ 
\index{equals(Object)}
\hypertarget{it.unisa.sesa.repominer.db.entities.ChangeForCommit.equals(java.lang.Object)}{{\bf  equals}\\}
\texttt{public boolean\ {\bf  equals}(\texttt{java.lang.Object} {\bf  arg0})
\label{it.unisa.sesa.repominer.db.entities.ChangeForCommit.equals(java.lang.Object)}}%end signature
}%end item
\item{ 
\index{getChangeHashId()}
\hypertarget{it.unisa.sesa.repominer.db.entities.ChangeForCommit.getChangeHashId()}{{\bf  getChangeHashId}\\}
\texttt{public java.lang.Integer\ {\bf  getChangeHashId}()
\label{it.unisa.sesa.repominer.db.entities.ChangeForCommit.getChangeHashId()}}%end signature
}%end item
\item{ 
\index{getDeletions()}
\hypertarget{it.unisa.sesa.repominer.db.entities.ChangeForCommit.getDeletions()}{{\bf  getDeletions}\\}
\texttt{public java.lang.Integer\ {\bf  getDeletions}()
\label{it.unisa.sesa.repominer.db.entities.ChangeForCommit.getDeletions()}}%end signature
}%end item
\item{ 
\index{getId()}
\hypertarget{it.unisa.sesa.repominer.db.entities.ChangeForCommit.getId()}{{\bf  getId}\\}
\texttt{public java.lang.Integer\ {\bf  getId}()
\label{it.unisa.sesa.repominer.db.entities.ChangeForCommit.getId()}}%end signature
}%end item
\item{ 
\index{getInsertions()}
\hypertarget{it.unisa.sesa.repominer.db.entities.ChangeForCommit.getInsertions()}{{\bf  getInsertions}\\}
\texttt{public java.lang.Integer\ {\bf  getInsertions}()
\label{it.unisa.sesa.repominer.db.entities.ChangeForCommit.getInsertions()}}%end signature
}%end item
\item{ 
\index{getModifiedFile()}
\hypertarget{it.unisa.sesa.repominer.db.entities.ChangeForCommit.getModifiedFile()}{{\bf  getModifiedFile}\\}
\texttt{public java.lang.String\ {\bf  getModifiedFile}()
\label{it.unisa.sesa.repominer.db.entities.ChangeForCommit.getModifiedFile()}}%end signature
}%end item
\item{ 
\index{hashCode()}
\hypertarget{it.unisa.sesa.repominer.db.entities.ChangeForCommit.hashCode()}{{\bf  hashCode}\\}
\texttt{public native int\ {\bf  hashCode}()
\label{it.unisa.sesa.repominer.db.entities.ChangeForCommit.hashCode()}}%end signature
}%end item
\item{ 
\index{setChangeHashId(Integer)}
\hypertarget{it.unisa.sesa.repominer.db.entities.ChangeForCommit.setChangeHashId(java.lang.Integer)}{{\bf  setChangeHashId}\\}
\texttt{public void\ {\bf  setChangeHashId}(\texttt{java.lang.Integer} {\bf  changeHash})
\label{it.unisa.sesa.repominer.db.entities.ChangeForCommit.setChangeHashId(java.lang.Integer)}}%end signature
}%end item
\item{ 
\index{setDeletions(Integer)}
\hypertarget{it.unisa.sesa.repominer.db.entities.ChangeForCommit.setDeletions(java.lang.Integer)}{{\bf  setDeletions}\\}
\texttt{public void\ {\bf  setDeletions}(\texttt{java.lang.Integer} {\bf  deletions})
\label{it.unisa.sesa.repominer.db.entities.ChangeForCommit.setDeletions(java.lang.Integer)}}%end signature
}%end item
\item{ 
\index{setId(Integer)}
\hypertarget{it.unisa.sesa.repominer.db.entities.ChangeForCommit.setId(java.lang.Integer)}{{\bf  setId}\\}
\texttt{public void\ {\bf  setId}(\texttt{java.lang.Integer} {\bf  id})
\label{it.unisa.sesa.repominer.db.entities.ChangeForCommit.setId(java.lang.Integer)}}%end signature
}%end item
\item{ 
\index{setInsertions(Integer)}
\hypertarget{it.unisa.sesa.repominer.db.entities.ChangeForCommit.setInsertions(java.lang.Integer)}{{\bf  setInsertions}\\}
\texttt{public void\ {\bf  setInsertions}(\texttt{java.lang.Integer} {\bf  insertions})
\label{it.unisa.sesa.repominer.db.entities.ChangeForCommit.setInsertions(java.lang.Integer)}}%end signature
}%end item
\item{ 
\index{setModifiedFile(String)}
\hypertarget{it.unisa.sesa.repominer.db.entities.ChangeForCommit.setModifiedFile(java.lang.String)}{{\bf  setModifiedFile}\\}
\texttt{public void\ {\bf  setModifiedFile}(\texttt{java.lang.String} {\bf  modifiedFile})
\label{it.unisa.sesa.repominer.db.entities.ChangeForCommit.setModifiedFile(java.lang.String)}}%end signature
}%end item
\item{ 
\index{toString()}
\hypertarget{it.unisa.sesa.repominer.db.entities.ChangeForCommit.toString()}{{\bf  toString}\\}
\texttt{public java.lang.String\ {\bf  toString}()
\label{it.unisa.sesa.repominer.db.entities.ChangeForCommit.toString()}}%end signature
}%end item
\end{itemize}
}
}
\subsection{\label{it.unisa.sesa.repominer.db.entities.ClassInvocation}Class ClassInvocation}{
\hypertarget{it.unisa.sesa.repominer.db.entities.ClassInvocation}{}\vskip .1in 
\subsubsection{Declaration}{
\small public class ClassInvocation
\\ {\bf  extends} java.lang.Object
\refdefined{java.lang.Object}}
\subsubsection{Constructor summary}{
\begin{verse}
\hyperlink{it.unisa.sesa.repominer.db.entities.ClassInvocation()}{{\bf ClassInvocation()}} \\
\hyperlink{it.unisa.sesa.repominer.db.entities.ClassInvocation(java.lang.Integer, java.lang.Integer)}{{\bf ClassInvocation(Integer, Integer)}} \\
\end{verse}
}
\subsubsection{Method summary}{
\begin{verse}
\hyperlink{it.unisa.sesa.repominer.db.entities.ClassInvocation.equals(java.lang.Object)}{{\bf equals(Object)}} \\
\hyperlink{it.unisa.sesa.repominer.db.entities.ClassInvocation.getInvokedClassId()}{{\bf getInvokedClassId()}} \\
\hyperlink{it.unisa.sesa.repominer.db.entities.ClassInvocation.getInvokerClassId()}{{\bf getInvokerClassId()}} \\
\hyperlink{it.unisa.sesa.repominer.db.entities.ClassInvocation.hashCode()}{{\bf hashCode()}} \\
\hyperlink{it.unisa.sesa.repominer.db.entities.ClassInvocation.setInvokedClassId(java.lang.Integer)}{{\bf setInvokedClassId(Integer)}} \\
\hyperlink{it.unisa.sesa.repominer.db.entities.ClassInvocation.setInvokerClassId(java.lang.Integer)}{{\bf setInvokerClassId(Integer)}} \\
\hyperlink{it.unisa.sesa.repominer.db.entities.ClassInvocation.toString()}{{\bf toString()}} \\
\end{verse}
}
\subsubsection{Constructors}{
\vskip -2em
\begin{itemize}
\item{ 
\index{ClassInvocation()}
\hypertarget{it.unisa.sesa.repominer.db.entities.ClassInvocation()}{{\bf  ClassInvocation}\\}
\texttt{public\ {\bf  ClassInvocation}()
\label{it.unisa.sesa.repominer.db.entities.ClassInvocation()}}%end signature
}%end item
\item{ 
\index{ClassInvocation(Integer, Integer)}
\hypertarget{it.unisa.sesa.repominer.db.entities.ClassInvocation(java.lang.Integer, java.lang.Integer)}{{\bf  ClassInvocation}\\}
\texttt{public\ {\bf  ClassInvocation}(\texttt{java.lang.Integer} {\bf  invokedClassId},
\texttt{java.lang.Integer} {\bf  invokerClassId})
\label{it.unisa.sesa.repominer.db.entities.ClassInvocation(java.lang.Integer, java.lang.Integer)}}%end signature
}%end item
\end{itemize}
}
\subsubsection{Methods}{
\vskip -2em
\begin{itemize}
\item{ 
\index{equals(Object)}
\hypertarget{it.unisa.sesa.repominer.db.entities.ClassInvocation.equals(java.lang.Object)}{{\bf  equals}\\}
\texttt{public boolean\ {\bf  equals}(\texttt{java.lang.Object} {\bf  arg0})
\label{it.unisa.sesa.repominer.db.entities.ClassInvocation.equals(java.lang.Object)}}%end signature
}%end item
\item{ 
\index{getInvokedClassId()}
\hypertarget{it.unisa.sesa.repominer.db.entities.ClassInvocation.getInvokedClassId()}{{\bf  getInvokedClassId}\\}
\texttt{public java.lang.Integer\ {\bf  getInvokedClassId}()
\label{it.unisa.sesa.repominer.db.entities.ClassInvocation.getInvokedClassId()}}%end signature
}%end item
\item{ 
\index{getInvokerClassId()}
\hypertarget{it.unisa.sesa.repominer.db.entities.ClassInvocation.getInvokerClassId()}{{\bf  getInvokerClassId}\\}
\texttt{public java.lang.Integer\ {\bf  getInvokerClassId}()
\label{it.unisa.sesa.repominer.db.entities.ClassInvocation.getInvokerClassId()}}%end signature
}%end item
\item{ 
\index{hashCode()}
\hypertarget{it.unisa.sesa.repominer.db.entities.ClassInvocation.hashCode()}{{\bf  hashCode}\\}
\texttt{public native int\ {\bf  hashCode}()
\label{it.unisa.sesa.repominer.db.entities.ClassInvocation.hashCode()}}%end signature
}%end item
\item{ 
\index{setInvokedClassId(Integer)}
\hypertarget{it.unisa.sesa.repominer.db.entities.ClassInvocation.setInvokedClassId(java.lang.Integer)}{{\bf  setInvokedClassId}\\}
\texttt{public void\ {\bf  setInvokedClassId}(\texttt{java.lang.Integer} {\bf  invokedClassId})
\label{it.unisa.sesa.repominer.db.entities.ClassInvocation.setInvokedClassId(java.lang.Integer)}}%end signature
}%end item
\item{ 
\index{setInvokerClassId(Integer)}
\hypertarget{it.unisa.sesa.repominer.db.entities.ClassInvocation.setInvokerClassId(java.lang.Integer)}{{\bf  setInvokerClassId}\\}
\texttt{public void\ {\bf  setInvokerClassId}(\texttt{java.lang.Integer} {\bf  invokerClassId})
\label{it.unisa.sesa.repominer.db.entities.ClassInvocation.setInvokerClassId(java.lang.Integer)}}%end signature
}%end item
\item{ 
\index{toString()}
\hypertarget{it.unisa.sesa.repominer.db.entities.ClassInvocation.toString()}{{\bf  toString}\\}
\texttt{public java.lang.String\ {\bf  toString}()
\label{it.unisa.sesa.repominer.db.entities.ClassInvocation.toString()}}%end signature
}%end item
\end{itemize}
}
}
\subsection{\label{it.unisa.sesa.repominer.db.entities.Import}Class Import}{
\hypertarget{it.unisa.sesa.repominer.db.entities.Import}{}\vskip .1in 
\subsubsection{Declaration}{
\small public class Import
\\ {\bf  extends} java.lang.Object
\refdefined{java.lang.Object}}
\subsubsection{Constructor summary}{
\begin{verse}
\hyperlink{it.unisa.sesa.repominer.db.entities.Import()}{{\bf Import()}} \\
\hyperlink{it.unisa.sesa.repominer.db.entities.Import(java.lang.Integer, java.lang.String)}{{\bf Import(Integer, String)}} \\
\end{verse}
}
\subsubsection{Method summary}{
\begin{verse}
\hyperlink{it.unisa.sesa.repominer.db.entities.Import.equals(java.lang.Object)}{{\bf equals(Object)}} \\
\hyperlink{it.unisa.sesa.repominer.db.entities.Import.getId()}{{\bf getId()}} \\
\hyperlink{it.unisa.sesa.repominer.db.entities.Import.getName()}{{\bf getName()}} \\
\hyperlink{it.unisa.sesa.repominer.db.entities.Import.hashCode()}{{\bf hashCode()}} \\
\hyperlink{it.unisa.sesa.repominer.db.entities.Import.setId(java.lang.Integer)}{{\bf setId(Integer)}} \\
\hyperlink{it.unisa.sesa.repominer.db.entities.Import.setName(java.lang.String)}{{\bf setName(String)}} \\
\hyperlink{it.unisa.sesa.repominer.db.entities.Import.toString()}{{\bf toString()}} \\
\end{verse}
}
\subsubsection{Constructors}{
\vskip -2em
\begin{itemize}
\item{ 
\index{Import()}
\hypertarget{it.unisa.sesa.repominer.db.entities.Import()}{{\bf  Import}\\}
\texttt{public\ {\bf  Import}()
\label{it.unisa.sesa.repominer.db.entities.Import()}}%end signature
}%end item
\item{ 
\index{Import(Integer, String)}
\hypertarget{it.unisa.sesa.repominer.db.entities.Import(java.lang.Integer, java.lang.String)}{{\bf  Import}\\}
\texttt{public\ {\bf  Import}(\texttt{java.lang.Integer} {\bf  id},
\texttt{java.lang.String} {\bf  name})
\label{it.unisa.sesa.repominer.db.entities.Import(java.lang.Integer, java.lang.String)}}%end signature
}%end item
\end{itemize}
}
\subsubsection{Methods}{
\vskip -2em
\begin{itemize}
\item{ 
\index{equals(Object)}
\hypertarget{it.unisa.sesa.repominer.db.entities.Import.equals(java.lang.Object)}{{\bf  equals}\\}
\texttt{public boolean\ {\bf  equals}(\texttt{java.lang.Object} {\bf  arg0})
\label{it.unisa.sesa.repominer.db.entities.Import.equals(java.lang.Object)}}%end signature
}%end item
\item{ 
\index{getId()}
\hypertarget{it.unisa.sesa.repominer.db.entities.Import.getId()}{{\bf  getId}\\}
\texttt{public java.lang.Integer\ {\bf  getId}()
\label{it.unisa.sesa.repominer.db.entities.Import.getId()}}%end signature
}%end item
\item{ 
\index{getName()}
\hypertarget{it.unisa.sesa.repominer.db.entities.Import.getName()}{{\bf  getName}\\}
\texttt{public java.lang.String\ {\bf  getName}()
\label{it.unisa.sesa.repominer.db.entities.Import.getName()}}%end signature
}%end item
\item{ 
\index{hashCode()}
\hypertarget{it.unisa.sesa.repominer.db.entities.Import.hashCode()}{{\bf  hashCode}\\}
\texttt{public native int\ {\bf  hashCode}()
\label{it.unisa.sesa.repominer.db.entities.Import.hashCode()}}%end signature
}%end item
\item{ 
\index{setId(Integer)}
\hypertarget{it.unisa.sesa.repominer.db.entities.Import.setId(java.lang.Integer)}{{\bf  setId}\\}
\texttt{public void\ {\bf  setId}(\texttt{java.lang.Integer} {\bf  id})
\label{it.unisa.sesa.repominer.db.entities.Import.setId(java.lang.Integer)}}%end signature
}%end item
\item{ 
\index{setName(String)}
\hypertarget{it.unisa.sesa.repominer.db.entities.Import.setName(java.lang.String)}{{\bf  setName}\\}
\texttt{public void\ {\bf  setName}(\texttt{java.lang.String} {\bf  name})
\label{it.unisa.sesa.repominer.db.entities.Import.setName(java.lang.String)}}%end signature
}%end item
\item{ 
\index{toString()}
\hypertarget{it.unisa.sesa.repominer.db.entities.Import.toString()}{{\bf  toString}\\}
\texttt{public java.lang.String\ {\bf  toString}()
\label{it.unisa.sesa.repominer.db.entities.Import.toString()}}%end signature
}%end item
\end{itemize}
}
}
\subsection{\label{it.unisa.sesa.repominer.db.entities.Imports}Class Imports}{
\hypertarget{it.unisa.sesa.repominer.db.entities.Imports}{}\vskip .1in 
\subsubsection{Declaration}{
\small public class Imports
\\ {\bf  extends} java.lang.Object
\refdefined{java.lang.Object}}
\subsubsection{Constructor summary}{
\begin{verse}
\hyperlink{it.unisa.sesa.repominer.db.entities.Imports()}{{\bf Imports()}} \\
\hyperlink{it.unisa.sesa.repominer.db.entities.Imports(java.lang.Integer, java.lang.Integer)}{{\bf Imports(Integer, Integer)}} \\
\end{verse}
}
\subsubsection{Method summary}{
\begin{verse}
\hyperlink{it.unisa.sesa.repominer.db.entities.Imports.equals(java.lang.Object)}{{\bf equals(Object)}} \\
\hyperlink{it.unisa.sesa.repominer.db.entities.Imports.getImportedId()}{{\bf getImportedId()}} \\
\hyperlink{it.unisa.sesa.repominer.db.entities.Imports.getImporterId()}{{\bf getImporterId()}} \\
\hyperlink{it.unisa.sesa.repominer.db.entities.Imports.hashCode()}{{\bf hashCode()}} \\
\hyperlink{it.unisa.sesa.repominer.db.entities.Imports.setImportedId(java.lang.Integer)}{{\bf setImportedId(Integer)}} \\
\hyperlink{it.unisa.sesa.repominer.db.entities.Imports.setImporterId(java.lang.Integer)}{{\bf setImporterId(Integer)}} \\
\hyperlink{it.unisa.sesa.repominer.db.entities.Imports.toString()}{{\bf toString()}} \\
\end{verse}
}
\subsubsection{Constructors}{
\vskip -2em
\begin{itemize}
\item{ 
\index{Imports()}
\hypertarget{it.unisa.sesa.repominer.db.entities.Imports()}{{\bf  Imports}\\}
\texttt{public\ {\bf  Imports}()
\label{it.unisa.sesa.repominer.db.entities.Imports()}}%end signature
}%end item
\item{ 
\index{Imports(Integer, Integer)}
\hypertarget{it.unisa.sesa.repominer.db.entities.Imports(java.lang.Integer, java.lang.Integer)}{{\bf  Imports}\\}
\texttt{public\ {\bf  Imports}(\texttt{java.lang.Integer} {\bf  importerId},
\texttt{java.lang.Integer} {\bf  importedId})
\label{it.unisa.sesa.repominer.db.entities.Imports(java.lang.Integer, java.lang.Integer)}}%end signature
}%end item
\end{itemize}
}
\subsubsection{Methods}{
\vskip -2em
\begin{itemize}
\item{ 
\index{equals(Object)}
\hypertarget{it.unisa.sesa.repominer.db.entities.Imports.equals(java.lang.Object)}{{\bf  equals}\\}
\texttt{public boolean\ {\bf  equals}(\texttt{java.lang.Object} {\bf  arg0})
\label{it.unisa.sesa.repominer.db.entities.Imports.equals(java.lang.Object)}}%end signature
}%end item
\item{ 
\index{getImportedId()}
\hypertarget{it.unisa.sesa.repominer.db.entities.Imports.getImportedId()}{{\bf  getImportedId}\\}
\texttt{public java.lang.Integer\ {\bf  getImportedId}()
\label{it.unisa.sesa.repominer.db.entities.Imports.getImportedId()}}%end signature
}%end item
\item{ 
\index{getImporterId()}
\hypertarget{it.unisa.sesa.repominer.db.entities.Imports.getImporterId()}{{\bf  getImporterId}\\}
\texttt{public java.lang.Integer\ {\bf  getImporterId}()
\label{it.unisa.sesa.repominer.db.entities.Imports.getImporterId()}}%end signature
}%end item
\item{ 
\index{hashCode()}
\hypertarget{it.unisa.sesa.repominer.db.entities.Imports.hashCode()}{{\bf  hashCode}\\}
\texttt{public native int\ {\bf  hashCode}()
\label{it.unisa.sesa.repominer.db.entities.Imports.hashCode()}}%end signature
}%end item
\item{ 
\index{setImportedId(Integer)}
\hypertarget{it.unisa.sesa.repominer.db.entities.Imports.setImportedId(java.lang.Integer)}{{\bf  setImportedId}\\}
\texttt{public void\ {\bf  setImportedId}(\texttt{java.lang.Integer} {\bf  importedId})
\label{it.unisa.sesa.repominer.db.entities.Imports.setImportedId(java.lang.Integer)}}%end signature
}%end item
\item{ 
\index{setImporterId(Integer)}
\hypertarget{it.unisa.sesa.repominer.db.entities.Imports.setImporterId(java.lang.Integer)}{{\bf  setImporterId}\\}
\texttt{public void\ {\bf  setImporterId}(\texttt{java.lang.Integer} {\bf  importerId})
\label{it.unisa.sesa.repominer.db.entities.Imports.setImporterId(java.lang.Integer)}}%end signature
}%end item
\item{ 
\index{toString()}
\hypertarget{it.unisa.sesa.repominer.db.entities.Imports.toString()}{{\bf  toString}\\}
\texttt{public java.lang.String\ {\bf  toString}()
\label{it.unisa.sesa.repominer.db.entities.Imports.toString()}}%end signature
}%end item
\end{itemize}
}
}
\subsection{\label{it.unisa.sesa.repominer.db.entities.Issue}Class Issue}{
\hypertarget{it.unisa.sesa.repominer.db.entities.Issue}{}\vskip .1in 
\subsubsection{Declaration}{
\small public class Issue
\\ {\bf  extends} java.lang.Object
\refdefined{java.lang.Object}}
\subsubsection{Constructor summary}{
\begin{verse}
\hyperlink{it.unisa.sesa.repominer.db.entities.Issue()}{{\bf Issue()}} \\
\hyperlink{it.unisa.sesa.repominer.db.entities.Issue(java.lang.Integer, java.lang.String, java.lang.String, java.lang.String, java.lang.String, java.lang.String, java.lang.String, java.lang.String, java.lang.String, java.lang.String, java.util.Date, java.util.Date, java.util.Date, java.lang.Integer)}{{\bf Issue(Integer, String, String, String, String, String, String, String, String, String, Date, Date, Date, Integer)}} \\
\end{verse}
}
\subsubsection{Method summary}{
\begin{verse}
\hyperlink{it.unisa.sesa.repominer.db.entities.Issue.equals(java.lang.Object)}{{\bf equals(Object)}} \\
\hyperlink{it.unisa.sesa.repominer.db.entities.Issue.getAffectedVersion()}{{\bf getAffectedVersion()}} \\
\hyperlink{it.unisa.sesa.repominer.db.entities.Issue.getAssignee()}{{\bf getAssignee()}} \\
\hyperlink{it.unisa.sesa.repominer.db.entities.Issue.getClosed()}{{\bf getClosed()}} \\
\hyperlink{it.unisa.sesa.repominer.db.entities.Issue.getComponent()}{{\bf getComponent()}} \\
\hyperlink{it.unisa.sesa.repominer.db.entities.Issue.getCreated()}{{\bf getCreated()}} \\
\hyperlink{it.unisa.sesa.repominer.db.entities.Issue.getFixVersion()}{{\bf getFixVersion()}} \\
\hyperlink{it.unisa.sesa.repominer.db.entities.Issue.getId()}{{\bf getId()}} \\
\hyperlink{it.unisa.sesa.repominer.db.entities.Issue.getPriority()}{{\bf getPriority()}} \\
\hyperlink{it.unisa.sesa.repominer.db.entities.Issue.getProjectId()}{{\bf getProjectId()}} \\
\hyperlink{it.unisa.sesa.repominer.db.entities.Issue.getReporter()}{{\bf getReporter()}} \\
\hyperlink{it.unisa.sesa.repominer.db.entities.Issue.getResolution()}{{\bf getResolution()}} \\
\hyperlink{it.unisa.sesa.repominer.db.entities.Issue.getStatus()}{{\bf getStatus()}} \\
\hyperlink{it.unisa.sesa.repominer.db.entities.Issue.getType()}{{\bf getType()}} \\
\hyperlink{it.unisa.sesa.repominer.db.entities.Issue.getUpdated()}{{\bf getUpdated()}} \\
\hyperlink{it.unisa.sesa.repominer.db.entities.Issue.hashCode()}{{\bf hashCode()}} \\
\hyperlink{it.unisa.sesa.repominer.db.entities.Issue.setAffectedVersion(java.lang.String)}{{\bf setAffectedVersion(String)}} \\
\hyperlink{it.unisa.sesa.repominer.db.entities.Issue.setAssignee(java.lang.String)}{{\bf setAssignee(String)}} \\
\hyperlink{it.unisa.sesa.repominer.db.entities.Issue.setClosed(java.util.Date)}{{\bf setClosed(Date)}} \\
\hyperlink{it.unisa.sesa.repominer.db.entities.Issue.setComponent(java.lang.String)}{{\bf setComponent(String)}} \\
\hyperlink{it.unisa.sesa.repominer.db.entities.Issue.setCreated(java.util.Date)}{{\bf setCreated(Date)}} \\
\hyperlink{it.unisa.sesa.repominer.db.entities.Issue.setFixVersion(java.lang.String)}{{\bf setFixVersion(String)}} \\
\hyperlink{it.unisa.sesa.repominer.db.entities.Issue.setId(java.lang.Integer)}{{\bf setId(Integer)}} \\
\hyperlink{it.unisa.sesa.repominer.db.entities.Issue.setPriority(java.lang.String)}{{\bf setPriority(String)}} \\
\hyperlink{it.unisa.sesa.repominer.db.entities.Issue.setProjectId(java.lang.Integer)}{{\bf setProjectId(Integer)}} \\
\hyperlink{it.unisa.sesa.repominer.db.entities.Issue.setReporter(java.lang.String)}{{\bf setReporter(String)}} \\
\hyperlink{it.unisa.sesa.repominer.db.entities.Issue.setResolution(java.lang.String)}{{\bf setResolution(String)}} \\
\hyperlink{it.unisa.sesa.repominer.db.entities.Issue.setStatus(java.lang.String)}{{\bf setStatus(String)}} \\
\hyperlink{it.unisa.sesa.repominer.db.entities.Issue.setType(java.lang.String)}{{\bf setType(String)}} \\
\hyperlink{it.unisa.sesa.repominer.db.entities.Issue.setUpdated(java.util.Date)}{{\bf setUpdated(Date)}} \\
\hyperlink{it.unisa.sesa.repominer.db.entities.Issue.toString()}{{\bf toString()}} \\
\end{verse}
}
\subsubsection{Constructors}{
\vskip -2em
\begin{itemize}
\item{ 
\index{Issue()}
\hypertarget{it.unisa.sesa.repominer.db.entities.Issue()}{{\bf  Issue}\\}
\texttt{public\ {\bf  Issue}()
\label{it.unisa.sesa.repominer.db.entities.Issue()}}%end signature
}%end item
\item{ 
\index{Issue(Integer, String, String, String, String, String, String, String, String, String, Date, Date, Date, Integer)}
\hypertarget{it.unisa.sesa.repominer.db.entities.Issue(java.lang.Integer, java.lang.String, java.lang.String, java.lang.String, java.lang.String, java.lang.String, java.lang.String, java.lang.String, java.lang.String, java.lang.String, java.util.Date, java.util.Date, java.util.Date, java.lang.Integer)}{{\bf  Issue}\\}
\texttt{public\ {\bf  Issue}(\texttt{java.lang.Integer} {\bf  id},
\texttt{java.lang.String} {\bf  type},
\texttt{java.lang.String} {\bf  priority},
\texttt{java.lang.String} {\bf  status},
\texttt{java.lang.String} {\bf  resolution},
\texttt{java.lang.String} {\bf  affectedVersion},
\texttt{java.lang.String} {\bf  component},
\texttt{java.lang.String} {\bf  fixVersion},
\texttt{java.lang.String} {\bf  assignee},
\texttt{java.lang.String} {\bf  reporter},
\texttt{java.util.Date} {\bf  updated},
\texttt{java.util.Date} {\bf  closed},
\texttt{java.util.Date} {\bf  created},
\texttt{java.lang.Integer} {\bf  projectId})
\label{it.unisa.sesa.repominer.db.entities.Issue(java.lang.Integer, java.lang.String, java.lang.String, java.lang.String, java.lang.String, java.lang.String, java.lang.String, java.lang.String, java.lang.String, java.lang.String, java.util.Date, java.util.Date, java.util.Date, java.lang.Integer)}}%end signature
}%end item
\end{itemize}
}
\subsubsection{Methods}{
\vskip -2em
\begin{itemize}
\item{ 
\index{equals(Object)}
\hypertarget{it.unisa.sesa.repominer.db.entities.Issue.equals(java.lang.Object)}{{\bf  equals}\\}
\texttt{public boolean\ {\bf  equals}(\texttt{java.lang.Object} {\bf  arg0})
\label{it.unisa.sesa.repominer.db.entities.Issue.equals(java.lang.Object)}}%end signature
}%end item
\item{ 
\index{getAffectedVersion()}
\hypertarget{it.unisa.sesa.repominer.db.entities.Issue.getAffectedVersion()}{{\bf  getAffectedVersion}\\}
\texttt{public java.lang.String\ {\bf  getAffectedVersion}()
\label{it.unisa.sesa.repominer.db.entities.Issue.getAffectedVersion()}}%end signature
}%end item
\item{ 
\index{getAssignee()}
\hypertarget{it.unisa.sesa.repominer.db.entities.Issue.getAssignee()}{{\bf  getAssignee}\\}
\texttt{public java.lang.String\ {\bf  getAssignee}()
\label{it.unisa.sesa.repominer.db.entities.Issue.getAssignee()}}%end signature
}%end item
\item{ 
\index{getClosed()}
\hypertarget{it.unisa.sesa.repominer.db.entities.Issue.getClosed()}{{\bf  getClosed}\\}
\texttt{public java.util.Date\ {\bf  getClosed}()
\label{it.unisa.sesa.repominer.db.entities.Issue.getClosed()}}%end signature
}%end item
\item{ 
\index{getComponent()}
\hypertarget{it.unisa.sesa.repominer.db.entities.Issue.getComponent()}{{\bf  getComponent}\\}
\texttt{public java.lang.String\ {\bf  getComponent}()
\label{it.unisa.sesa.repominer.db.entities.Issue.getComponent()}}%end signature
}%end item
\item{ 
\index{getCreated()}
\hypertarget{it.unisa.sesa.repominer.db.entities.Issue.getCreated()}{{\bf  getCreated}\\}
\texttt{public java.util.Date\ {\bf  getCreated}()
\label{it.unisa.sesa.repominer.db.entities.Issue.getCreated()}}%end signature
}%end item
\item{ 
\index{getFixVersion()}
\hypertarget{it.unisa.sesa.repominer.db.entities.Issue.getFixVersion()}{{\bf  getFixVersion}\\}
\texttt{public java.lang.String\ {\bf  getFixVersion}()
\label{it.unisa.sesa.repominer.db.entities.Issue.getFixVersion()}}%end signature
}%end item
\item{ 
\index{getId()}
\hypertarget{it.unisa.sesa.repominer.db.entities.Issue.getId()}{{\bf  getId}\\}
\texttt{public java.lang.Integer\ {\bf  getId}()
\label{it.unisa.sesa.repominer.db.entities.Issue.getId()}}%end signature
}%end item
\item{ 
\index{getPriority()}
\hypertarget{it.unisa.sesa.repominer.db.entities.Issue.getPriority()}{{\bf  getPriority}\\}
\texttt{public java.lang.String\ {\bf  getPriority}()
\label{it.unisa.sesa.repominer.db.entities.Issue.getPriority()}}%end signature
}%end item
\item{ 
\index{getProjectId()}
\hypertarget{it.unisa.sesa.repominer.db.entities.Issue.getProjectId()}{{\bf  getProjectId}\\}
\texttt{public java.lang.Integer\ {\bf  getProjectId}()
\label{it.unisa.sesa.repominer.db.entities.Issue.getProjectId()}}%end signature
}%end item
\item{ 
\index{getReporter()}
\hypertarget{it.unisa.sesa.repominer.db.entities.Issue.getReporter()}{{\bf  getReporter}\\}
\texttt{public java.lang.String\ {\bf  getReporter}()
\label{it.unisa.sesa.repominer.db.entities.Issue.getReporter()}}%end signature
}%end item
\item{ 
\index{getResolution()}
\hypertarget{it.unisa.sesa.repominer.db.entities.Issue.getResolution()}{{\bf  getResolution}\\}
\texttt{public java.lang.String\ {\bf  getResolution}()
\label{it.unisa.sesa.repominer.db.entities.Issue.getResolution()}}%end signature
}%end item
\item{ 
\index{getStatus()}
\hypertarget{it.unisa.sesa.repominer.db.entities.Issue.getStatus()}{{\bf  getStatus}\\}
\texttt{public java.lang.String\ {\bf  getStatus}()
\label{it.unisa.sesa.repominer.db.entities.Issue.getStatus()}}%end signature
}%end item
\item{ 
\index{getType()}
\hypertarget{it.unisa.sesa.repominer.db.entities.Issue.getType()}{{\bf  getType}\\}
\texttt{public java.lang.String\ {\bf  getType}()
\label{it.unisa.sesa.repominer.db.entities.Issue.getType()}}%end signature
}%end item
\item{ 
\index{getUpdated()}
\hypertarget{it.unisa.sesa.repominer.db.entities.Issue.getUpdated()}{{\bf  getUpdated}\\}
\texttt{public java.util.Date\ {\bf  getUpdated}()
\label{it.unisa.sesa.repominer.db.entities.Issue.getUpdated()}}%end signature
}%end item
\item{ 
\index{hashCode()}
\hypertarget{it.unisa.sesa.repominer.db.entities.Issue.hashCode()}{{\bf  hashCode}\\}
\texttt{public native int\ {\bf  hashCode}()
\label{it.unisa.sesa.repominer.db.entities.Issue.hashCode()}}%end signature
}%end item
\item{ 
\index{setAffectedVersion(String)}
\hypertarget{it.unisa.sesa.repominer.db.entities.Issue.setAffectedVersion(java.lang.String)}{{\bf  setAffectedVersion}\\}
\texttt{public void\ {\bf  setAffectedVersion}(\texttt{java.lang.String} {\bf  affectedVersion})
\label{it.unisa.sesa.repominer.db.entities.Issue.setAffectedVersion(java.lang.String)}}%end signature
}%end item
\item{ 
\index{setAssignee(String)}
\hypertarget{it.unisa.sesa.repominer.db.entities.Issue.setAssignee(java.lang.String)}{{\bf  setAssignee}\\}
\texttt{public void\ {\bf  setAssignee}(\texttt{java.lang.String} {\bf  assignee})
\label{it.unisa.sesa.repominer.db.entities.Issue.setAssignee(java.lang.String)}}%end signature
}%end item
\item{ 
\index{setClosed(Date)}
\hypertarget{it.unisa.sesa.repominer.db.entities.Issue.setClosed(java.util.Date)}{{\bf  setClosed}\\}
\texttt{public void\ {\bf  setClosed}(\texttt{java.util.Date} {\bf  closed})
\label{it.unisa.sesa.repominer.db.entities.Issue.setClosed(java.util.Date)}}%end signature
}%end item
\item{ 
\index{setComponent(String)}
\hypertarget{it.unisa.sesa.repominer.db.entities.Issue.setComponent(java.lang.String)}{{\bf  setComponent}\\}
\texttt{public void\ {\bf  setComponent}(\texttt{java.lang.String} {\bf  component})
\label{it.unisa.sesa.repominer.db.entities.Issue.setComponent(java.lang.String)}}%end signature
}%end item
\item{ 
\index{setCreated(Date)}
\hypertarget{it.unisa.sesa.repominer.db.entities.Issue.setCreated(java.util.Date)}{{\bf  setCreated}\\}
\texttt{public void\ {\bf  setCreated}(\texttt{java.util.Date} {\bf  created})
\label{it.unisa.sesa.repominer.db.entities.Issue.setCreated(java.util.Date)}}%end signature
}%end item
\item{ 
\index{setFixVersion(String)}
\hypertarget{it.unisa.sesa.repominer.db.entities.Issue.setFixVersion(java.lang.String)}{{\bf  setFixVersion}\\}
\texttt{public void\ {\bf  setFixVersion}(\texttt{java.lang.String} {\bf  fixVersion})
\label{it.unisa.sesa.repominer.db.entities.Issue.setFixVersion(java.lang.String)}}%end signature
}%end item
\item{ 
\index{setId(Integer)}
\hypertarget{it.unisa.sesa.repominer.db.entities.Issue.setId(java.lang.Integer)}{{\bf  setId}\\}
\texttt{public void\ {\bf  setId}(\texttt{java.lang.Integer} {\bf  id})
\label{it.unisa.sesa.repominer.db.entities.Issue.setId(java.lang.Integer)}}%end signature
}%end item
\item{ 
\index{setPriority(String)}
\hypertarget{it.unisa.sesa.repominer.db.entities.Issue.setPriority(java.lang.String)}{{\bf  setPriority}\\}
\texttt{public void\ {\bf  setPriority}(\texttt{java.lang.String} {\bf  priority})
\label{it.unisa.sesa.repominer.db.entities.Issue.setPriority(java.lang.String)}}%end signature
}%end item
\item{ 
\index{setProjectId(Integer)}
\hypertarget{it.unisa.sesa.repominer.db.entities.Issue.setProjectId(java.lang.Integer)}{{\bf  setProjectId}\\}
\texttt{public void\ {\bf  setProjectId}(\texttt{java.lang.Integer} {\bf  projectId})
\label{it.unisa.sesa.repominer.db.entities.Issue.setProjectId(java.lang.Integer)}}%end signature
}%end item
\item{ 
\index{setReporter(String)}
\hypertarget{it.unisa.sesa.repominer.db.entities.Issue.setReporter(java.lang.String)}{{\bf  setReporter}\\}
\texttt{public void\ {\bf  setReporter}(\texttt{java.lang.String} {\bf  reporter})
\label{it.unisa.sesa.repominer.db.entities.Issue.setReporter(java.lang.String)}}%end signature
}%end item
\item{ 
\index{setResolution(String)}
\hypertarget{it.unisa.sesa.repominer.db.entities.Issue.setResolution(java.lang.String)}{{\bf  setResolution}\\}
\texttt{public void\ {\bf  setResolution}(\texttt{java.lang.String} {\bf  resolution})
\label{it.unisa.sesa.repominer.db.entities.Issue.setResolution(java.lang.String)}}%end signature
}%end item
\item{ 
\index{setStatus(String)}
\hypertarget{it.unisa.sesa.repominer.db.entities.Issue.setStatus(java.lang.String)}{{\bf  setStatus}\\}
\texttt{public void\ {\bf  setStatus}(\texttt{java.lang.String} {\bf  status})
\label{it.unisa.sesa.repominer.db.entities.Issue.setStatus(java.lang.String)}}%end signature
}%end item
\item{ 
\index{setType(String)}
\hypertarget{it.unisa.sesa.repominer.db.entities.Issue.setType(java.lang.String)}{{\bf  setType}\\}
\texttt{public void\ {\bf  setType}(\texttt{java.lang.String} {\bf  type})
\label{it.unisa.sesa.repominer.db.entities.Issue.setType(java.lang.String)}}%end signature
}%end item
\item{ 
\index{setUpdated(Date)}
\hypertarget{it.unisa.sesa.repominer.db.entities.Issue.setUpdated(java.util.Date)}{{\bf  setUpdated}\\}
\texttt{public void\ {\bf  setUpdated}(\texttt{java.util.Date} {\bf  updated})
\label{it.unisa.sesa.repominer.db.entities.Issue.setUpdated(java.util.Date)}}%end signature
}%end item
\item{ 
\index{toString()}
\hypertarget{it.unisa.sesa.repominer.db.entities.Issue.toString()}{{\bf  toString}\\}
\texttt{public java.lang.String\ {\bf  toString}()
\label{it.unisa.sesa.repominer.db.entities.Issue.toString()}}%end signature
}%end item
\end{itemize}
}
}
\subsection{\label{it.unisa.sesa.repominer.db.entities.IssueAttachment}Class IssueAttachment}{
\hypertarget{it.unisa.sesa.repominer.db.entities.IssueAttachment}{}\vskip .1in 
\subsubsection{Declaration}{
\small public class IssueAttachment
\\ {\bf  extends} java.lang.Object
\refdefined{java.lang.Object}}
\subsubsection{Constructor summary}{
\begin{verse}
\hyperlink{it.unisa.sesa.repominer.db.entities.IssueAttachment()}{{\bf IssueAttachment()}} \\
\hyperlink{it.unisa.sesa.repominer.db.entities.IssueAttachment(java.lang.Integer, java.sql.Date, java.lang.String, java.lang.Integer)}{{\bf IssueAttachment(Integer, Date, String, Integer)}} \\
\end{verse}
}
\subsubsection{Method summary}{
\begin{verse}
\hyperlink{it.unisa.sesa.repominer.db.entities.IssueAttachment.equals(java.lang.Object)}{{\bf equals(Object)}} \\
\hyperlink{it.unisa.sesa.repominer.db.entities.IssueAttachment.getBelongingIssueId()}{{\bf getBelongingIssueId()}} \\
\hyperlink{it.unisa.sesa.repominer.db.entities.IssueAttachment.getDate()}{{\bf getDate()}} \\
\hyperlink{it.unisa.sesa.repominer.db.entities.IssueAttachment.getId()}{{\bf getId()}} \\
\hyperlink{it.unisa.sesa.repominer.db.entities.IssueAttachment.getName()}{{\bf getName()}} \\
\hyperlink{it.unisa.sesa.repominer.db.entities.IssueAttachment.hashCode()}{{\bf hashCode()}} \\
\hyperlink{it.unisa.sesa.repominer.db.entities.IssueAttachment.setBelongingIssueId(java.lang.Integer)}{{\bf setBelongingIssueId(Integer)}} \\
\hyperlink{it.unisa.sesa.repominer.db.entities.IssueAttachment.setDate(java.sql.Date)}{{\bf setDate(Date)}} \\
\hyperlink{it.unisa.sesa.repominer.db.entities.IssueAttachment.setId(java.lang.Integer)}{{\bf setId(Integer)}} \\
\hyperlink{it.unisa.sesa.repominer.db.entities.IssueAttachment.setName(java.lang.String)}{{\bf setName(String)}} \\
\hyperlink{it.unisa.sesa.repominer.db.entities.IssueAttachment.toString()}{{\bf toString()}} \\
\end{verse}
}
\subsubsection{Constructors}{
\vskip -2em
\begin{itemize}
\item{ 
\index{IssueAttachment()}
\hypertarget{it.unisa.sesa.repominer.db.entities.IssueAttachment()}{{\bf  IssueAttachment}\\}
\texttt{public\ {\bf  IssueAttachment}()
\label{it.unisa.sesa.repominer.db.entities.IssueAttachment()}}%end signature
}%end item
\item{ 
\index{IssueAttachment(Integer, Date, String, Integer)}
\hypertarget{it.unisa.sesa.repominer.db.entities.IssueAttachment(java.lang.Integer, java.sql.Date, java.lang.String, java.lang.Integer)}{{\bf  IssueAttachment}\\}
\texttt{public\ {\bf  IssueAttachment}(\texttt{java.lang.Integer} {\bf  id},
\texttt{java.sql.Date} {\bf  date},
\texttt{java.lang.String} {\bf  name},
\texttt{java.lang.Integer} {\bf  belongingIssueId})
\label{it.unisa.sesa.repominer.db.entities.IssueAttachment(java.lang.Integer, java.sql.Date, java.lang.String, java.lang.Integer)}}%end signature
}%end item
\end{itemize}
}
\subsubsection{Methods}{
\vskip -2em
\begin{itemize}
\item{ 
\index{equals(Object)}
\hypertarget{it.unisa.sesa.repominer.db.entities.IssueAttachment.equals(java.lang.Object)}{{\bf  equals}\\}
\texttt{public boolean\ {\bf  equals}(\texttt{java.lang.Object} {\bf  arg0})
\label{it.unisa.sesa.repominer.db.entities.IssueAttachment.equals(java.lang.Object)}}%end signature
}%end item
\item{ 
\index{getBelongingIssueId()}
\hypertarget{it.unisa.sesa.repominer.db.entities.IssueAttachment.getBelongingIssueId()}{{\bf  getBelongingIssueId}\\}
\texttt{public java.lang.Integer\ {\bf  getBelongingIssueId}()
\label{it.unisa.sesa.repominer.db.entities.IssueAttachment.getBelongingIssueId()}}%end signature
}%end item
\item{ 
\index{getDate()}
\hypertarget{it.unisa.sesa.repominer.db.entities.IssueAttachment.getDate()}{{\bf  getDate}\\}
\texttt{public java.sql.Date\ {\bf  getDate}()
\label{it.unisa.sesa.repominer.db.entities.IssueAttachment.getDate()}}%end signature
}%end item
\item{ 
\index{getId()}
\hypertarget{it.unisa.sesa.repominer.db.entities.IssueAttachment.getId()}{{\bf  getId}\\}
\texttt{public java.lang.Integer\ {\bf  getId}()
\label{it.unisa.sesa.repominer.db.entities.IssueAttachment.getId()}}%end signature
}%end item
\item{ 
\index{getName()}
\hypertarget{it.unisa.sesa.repominer.db.entities.IssueAttachment.getName()}{{\bf  getName}\\}
\texttt{public java.lang.String\ {\bf  getName}()
\label{it.unisa.sesa.repominer.db.entities.IssueAttachment.getName()}}%end signature
}%end item
\item{ 
\index{hashCode()}
\hypertarget{it.unisa.sesa.repominer.db.entities.IssueAttachment.hashCode()}{{\bf  hashCode}\\}
\texttt{public native int\ {\bf  hashCode}()
\label{it.unisa.sesa.repominer.db.entities.IssueAttachment.hashCode()}}%end signature
}%end item
\item{ 
\index{setBelongingIssueId(Integer)}
\hypertarget{it.unisa.sesa.repominer.db.entities.IssueAttachment.setBelongingIssueId(java.lang.Integer)}{{\bf  setBelongingIssueId}\\}
\texttt{public void\ {\bf  setBelongingIssueId}(\texttt{java.lang.Integer} {\bf  belongingIssueId})
\label{it.unisa.sesa.repominer.db.entities.IssueAttachment.setBelongingIssueId(java.lang.Integer)}}%end signature
}%end item
\item{ 
\index{setDate(Date)}
\hypertarget{it.unisa.sesa.repominer.db.entities.IssueAttachment.setDate(java.sql.Date)}{{\bf  setDate}\\}
\texttt{public void\ {\bf  setDate}(\texttt{java.sql.Date} {\bf  date})
\label{it.unisa.sesa.repominer.db.entities.IssueAttachment.setDate(java.sql.Date)}}%end signature
}%end item
\item{ 
\index{setId(Integer)}
\hypertarget{it.unisa.sesa.repominer.db.entities.IssueAttachment.setId(java.lang.Integer)}{{\bf  setId}\\}
\texttt{public void\ {\bf  setId}(\texttt{java.lang.Integer} {\bf  id})
\label{it.unisa.sesa.repominer.db.entities.IssueAttachment.setId(java.lang.Integer)}}%end signature
}%end item
\item{ 
\index{setName(String)}
\hypertarget{it.unisa.sesa.repominer.db.entities.IssueAttachment.setName(java.lang.String)}{{\bf  setName}\\}
\texttt{public void\ {\bf  setName}(\texttt{java.lang.String} {\bf  name})
\label{it.unisa.sesa.repominer.db.entities.IssueAttachment.setName(java.lang.String)}}%end signature
}%end item
\item{ 
\index{toString()}
\hypertarget{it.unisa.sesa.repominer.db.entities.IssueAttachment.toString()}{{\bf  toString}\\}
\texttt{public java.lang.String\ {\bf  toString}()
\label{it.unisa.sesa.repominer.db.entities.IssueAttachment.toString()}}%end signature
}%end item
\end{itemize}
}
}
\subsection{\label{it.unisa.sesa.repominer.db.entities.IssueComment}Class IssueComment}{
\hypertarget{it.unisa.sesa.repominer.db.entities.IssueComment}{}\vskip .1in 
\subsubsection{Declaration}{
\small public class IssueComment
\\ {\bf  extends} java.lang.Object
\refdefined{java.lang.Object}}
\subsubsection{Constructor summary}{
\begin{verse}
\hyperlink{it.unisa.sesa.repominer.db.entities.IssueComment()}{{\bf IssueComment()}} \\
\hyperlink{it.unisa.sesa.repominer.db.entities.IssueComment(java.lang.Integer, java.lang.String, java.lang.String, java.lang.String, java.util.Date, java.lang.Integer)}{{\bf IssueComment(Integer, String, String, String, Date, Integer)}} \\
\end{verse}
}
\subsubsection{Method summary}{
\begin{verse}
\hyperlink{it.unisa.sesa.repominer.db.entities.IssueComment.equals(java.lang.Object)}{{\bf equals(Object)}} \\
\hyperlink{it.unisa.sesa.repominer.db.entities.IssueComment.getBelongingIssueId()}{{\bf getBelongingIssueId()}} \\
\hyperlink{it.unisa.sesa.repominer.db.entities.IssueComment.getDate()}{{\bf getDate()}} \\
\hyperlink{it.unisa.sesa.repominer.db.entities.IssueComment.getDevId()}{{\bf getDevId()}} \\
\hyperlink{it.unisa.sesa.repominer.db.entities.IssueComment.getDevMail()}{{\bf getDevMail()}} \\
\hyperlink{it.unisa.sesa.repominer.db.entities.IssueComment.getId()}{{\bf getId()}} \\
\hyperlink{it.unisa.sesa.repominer.db.entities.IssueComment.getText()}{{\bf getText()}} \\
\hyperlink{it.unisa.sesa.repominer.db.entities.IssueComment.hashCode()}{{\bf hashCode()}} \\
\hyperlink{it.unisa.sesa.repominer.db.entities.IssueComment.setBelongingIssueId(java.lang.Integer)}{{\bf setBelongingIssueId(Integer)}} \\
\hyperlink{it.unisa.sesa.repominer.db.entities.IssueComment.setDate(java.util.Date)}{{\bf setDate(Date)}} \\
\hyperlink{it.unisa.sesa.repominer.db.entities.IssueComment.setDevId(java.lang.String)}{{\bf setDevId(String)}} \\
\hyperlink{it.unisa.sesa.repominer.db.entities.IssueComment.setDevMail(java.lang.String)}{{\bf setDevMail(String)}} \\
\hyperlink{it.unisa.sesa.repominer.db.entities.IssueComment.setId(java.lang.Integer)}{{\bf setId(Integer)}} \\
\hyperlink{it.unisa.sesa.repominer.db.entities.IssueComment.setText(java.lang.String)}{{\bf setText(String)}} \\
\hyperlink{it.unisa.sesa.repominer.db.entities.IssueComment.toString()}{{\bf toString()}} \\
\end{verse}
}
\subsubsection{Constructors}{
\vskip -2em
\begin{itemize}
\item{ 
\index{IssueComment()}
\hypertarget{it.unisa.sesa.repominer.db.entities.IssueComment()}{{\bf  IssueComment}\\}
\texttt{public\ {\bf  IssueComment}()
\label{it.unisa.sesa.repominer.db.entities.IssueComment()}}%end signature
}%end item
\item{ 
\index{IssueComment(Integer, String, String, String, Date, Integer)}
\hypertarget{it.unisa.sesa.repominer.db.entities.IssueComment(java.lang.Integer, java.lang.String, java.lang.String, java.lang.String, java.util.Date, java.lang.Integer)}{{\bf  IssueComment}\\}
\texttt{public\ {\bf  IssueComment}(\texttt{java.lang.Integer} {\bf  id},
\texttt{java.lang.String} {\bf  devId},
\texttt{java.lang.String} {\bf  devMail},
\texttt{java.lang.String} {\bf  text},
\texttt{java.util.Date} {\bf  date},
\texttt{java.lang.Integer} {\bf  belongingIssueId})
\label{it.unisa.sesa.repominer.db.entities.IssueComment(java.lang.Integer, java.lang.String, java.lang.String, java.lang.String, java.util.Date, java.lang.Integer)}}%end signature
}%end item
\end{itemize}
}
\subsubsection{Methods}{
\vskip -2em
\begin{itemize}
\item{ 
\index{equals(Object)}
\hypertarget{it.unisa.sesa.repominer.db.entities.IssueComment.equals(java.lang.Object)}{{\bf  equals}\\}
\texttt{public boolean\ {\bf  equals}(\texttt{java.lang.Object} {\bf  arg0})
\label{it.unisa.sesa.repominer.db.entities.IssueComment.equals(java.lang.Object)}}%end signature
}%end item
\item{ 
\index{getBelongingIssueId()}
\hypertarget{it.unisa.sesa.repominer.db.entities.IssueComment.getBelongingIssueId()}{{\bf  getBelongingIssueId}\\}
\texttt{public java.lang.Integer\ {\bf  getBelongingIssueId}()
\label{it.unisa.sesa.repominer.db.entities.IssueComment.getBelongingIssueId()}}%end signature
}%end item
\item{ 
\index{getDate()}
\hypertarget{it.unisa.sesa.repominer.db.entities.IssueComment.getDate()}{{\bf  getDate}\\}
\texttt{public java.util.Date\ {\bf  getDate}()
\label{it.unisa.sesa.repominer.db.entities.IssueComment.getDate()}}%end signature
}%end item
\item{ 
\index{getDevId()}
\hypertarget{it.unisa.sesa.repominer.db.entities.IssueComment.getDevId()}{{\bf  getDevId}\\}
\texttt{public java.lang.String\ {\bf  getDevId}()
\label{it.unisa.sesa.repominer.db.entities.IssueComment.getDevId()}}%end signature
}%end item
\item{ 
\index{getDevMail()}
\hypertarget{it.unisa.sesa.repominer.db.entities.IssueComment.getDevMail()}{{\bf  getDevMail}\\}
\texttt{public java.lang.String\ {\bf  getDevMail}()
\label{it.unisa.sesa.repominer.db.entities.IssueComment.getDevMail()}}%end signature
}%end item
\item{ 
\index{getId()}
\hypertarget{it.unisa.sesa.repominer.db.entities.IssueComment.getId()}{{\bf  getId}\\}
\texttt{public java.lang.Integer\ {\bf  getId}()
\label{it.unisa.sesa.repominer.db.entities.IssueComment.getId()}}%end signature
}%end item
\item{ 
\index{getText()}
\hypertarget{it.unisa.sesa.repominer.db.entities.IssueComment.getText()}{{\bf  getText}\\}
\texttt{public java.lang.String\ {\bf  getText}()
\label{it.unisa.sesa.repominer.db.entities.IssueComment.getText()}}%end signature
}%end item
\item{ 
\index{hashCode()}
\hypertarget{it.unisa.sesa.repominer.db.entities.IssueComment.hashCode()}{{\bf  hashCode}\\}
\texttt{public native int\ {\bf  hashCode}()
\label{it.unisa.sesa.repominer.db.entities.IssueComment.hashCode()}}%end signature
}%end item
\item{ 
\index{setBelongingIssueId(Integer)}
\hypertarget{it.unisa.sesa.repominer.db.entities.IssueComment.setBelongingIssueId(java.lang.Integer)}{{\bf  setBelongingIssueId}\\}
\texttt{public void\ {\bf  setBelongingIssueId}(\texttt{java.lang.Integer} {\bf  belongingIssueId})
\label{it.unisa.sesa.repominer.db.entities.IssueComment.setBelongingIssueId(java.lang.Integer)}}%end signature
}%end item
\item{ 
\index{setDate(Date)}
\hypertarget{it.unisa.sesa.repominer.db.entities.IssueComment.setDate(java.util.Date)}{{\bf  setDate}\\}
\texttt{public void\ {\bf  setDate}(\texttt{java.util.Date} {\bf  date})
\label{it.unisa.sesa.repominer.db.entities.IssueComment.setDate(java.util.Date)}}%end signature
}%end item
\item{ 
\index{setDevId(String)}
\hypertarget{it.unisa.sesa.repominer.db.entities.IssueComment.setDevId(java.lang.String)}{{\bf  setDevId}\\}
\texttt{public void\ {\bf  setDevId}(\texttt{java.lang.String} {\bf  devId})
\label{it.unisa.sesa.repominer.db.entities.IssueComment.setDevId(java.lang.String)}}%end signature
}%end item
\item{ 
\index{setDevMail(String)}
\hypertarget{it.unisa.sesa.repominer.db.entities.IssueComment.setDevMail(java.lang.String)}{{\bf  setDevMail}\\}
\texttt{public void\ {\bf  setDevMail}(\texttt{java.lang.String} {\bf  devMail})
\label{it.unisa.sesa.repominer.db.entities.IssueComment.setDevMail(java.lang.String)}}%end signature
}%end item
\item{ 
\index{setId(Integer)}
\hypertarget{it.unisa.sesa.repominer.db.entities.IssueComment.setId(java.lang.Integer)}{{\bf  setId}\\}
\texttt{public void\ {\bf  setId}(\texttt{java.lang.Integer} {\bf  id})
\label{it.unisa.sesa.repominer.db.entities.IssueComment.setId(java.lang.Integer)}}%end signature
}%end item
\item{ 
\index{setText(String)}
\hypertarget{it.unisa.sesa.repominer.db.entities.IssueComment.setText(java.lang.String)}{{\bf  setText}\\}
\texttt{public void\ {\bf  setText}(\texttt{java.lang.String} {\bf  text})
\label{it.unisa.sesa.repominer.db.entities.IssueComment.setText(java.lang.String)}}%end signature
}%end item
\item{ 
\index{toString()}
\hypertarget{it.unisa.sesa.repominer.db.entities.IssueComment.toString()}{{\bf  toString}\\}
\texttt{public java.lang.String\ {\bf  toString}()
\label{it.unisa.sesa.repominer.db.entities.IssueComment.toString()}}%end signature
}%end item
\end{itemize}
}
}
\subsection{\label{it.unisa.sesa.repominer.db.entities.Method}Class Method}{
\hypertarget{it.unisa.sesa.repominer.db.entities.Method}{}\vskip .1in 
\subsubsection{Declaration}{
\small public class Method
\\ {\bf  extends} java.lang.Object
\refdefined{java.lang.Object}}
\subsubsection{Constructor summary}{
\begin{verse}
\hyperlink{it.unisa.sesa.repominer.db.entities.Method()}{{\bf Method()}} \\
\hyperlink{it.unisa.sesa.repominer.db.entities.Method(java.lang.Integer, java.lang.Integer, java.lang.String, java.lang.Integer, java.lang.Integer)}{{\bf Method(Integer, Integer, String, Integer, Integer)}} \\
\end{verse}
}
\subsubsection{Method summary}{
\begin{verse}
\hyperlink{it.unisa.sesa.repominer.db.entities.Method.equals(java.lang.Object)}{{\bf equals(Object)}} \\
\hyperlink{it.unisa.sesa.repominer.db.entities.Method.getBelongingTypeId()}{{\bf getBelongingTypeId()}} \\
\hyperlink{it.unisa.sesa.repominer.db.entities.Method.getId()}{{\bf getId()}} \\
\hyperlink{it.unisa.sesa.repominer.db.entities.Method.getIsConstructor()}{{\bf getIsConstructor()}} \\
\hyperlink{it.unisa.sesa.repominer.db.entities.Method.getLineNumber()}{{\bf getLineNumber()}} \\
\hyperlink{it.unisa.sesa.repominer.db.entities.Method.getReturnTypeId()}{{\bf getReturnTypeId()}} \\
\hyperlink{it.unisa.sesa.repominer.db.entities.Method.hashCode()}{{\bf hashCode()}} \\
\hyperlink{it.unisa.sesa.repominer.db.entities.Method.setBelongingTypeId(java.lang.Integer)}{{\bf setBelongingTypeId(Integer)}} \\
\hyperlink{it.unisa.sesa.repominer.db.entities.Method.setId(java.lang.Integer)}{{\bf setId(Integer)}} \\
\hyperlink{it.unisa.sesa.repominer.db.entities.Method.setIsConstructor(java.lang.String)}{{\bf setIsConstructor(String)}} \\
\hyperlink{it.unisa.sesa.repominer.db.entities.Method.setLineNumber(java.lang.Integer)}{{\bf setLineNumber(Integer)}} \\
\hyperlink{it.unisa.sesa.repominer.db.entities.Method.setReturnTypeId(java.lang.Integer)}{{\bf setReturnTypeId(Integer)}} \\
\hyperlink{it.unisa.sesa.repominer.db.entities.Method.toString()}{{\bf toString()}} \\
\end{verse}
}
\subsubsection{Constructors}{
\vskip -2em
\begin{itemize}
\item{ 
\index{Method()}
\hypertarget{it.unisa.sesa.repominer.db.entities.Method()}{{\bf  Method}\\}
\texttt{public\ {\bf  Method}()
\label{it.unisa.sesa.repominer.db.entities.Method()}}%end signature
}%end item
\item{ 
\index{Method(Integer, Integer, String, Integer, Integer)}
\hypertarget{it.unisa.sesa.repominer.db.entities.Method(java.lang.Integer, java.lang.Integer, java.lang.String, java.lang.Integer, java.lang.Integer)}{{\bf  Method}\\}
\texttt{public\ {\bf  Method}(\texttt{java.lang.Integer} {\bf  id},
\texttt{java.lang.Integer} {\bf  lineNumber},
\texttt{java.lang.String} {\bf  isConstructor},
\texttt{java.lang.Integer} {\bf  belongingTypeId},
\texttt{java.lang.Integer} {\bf  returnTypeId})
\label{it.unisa.sesa.repominer.db.entities.Method(java.lang.Integer, java.lang.Integer, java.lang.String, java.lang.Integer, java.lang.Integer)}}%end signature
}%end item
\end{itemize}
}
\subsubsection{Methods}{
\vskip -2em
\begin{itemize}
\item{ 
\index{equals(Object)}
\hypertarget{it.unisa.sesa.repominer.db.entities.Method.equals(java.lang.Object)}{{\bf  equals}\\}
\texttt{public boolean\ {\bf  equals}(\texttt{java.lang.Object} {\bf  arg0})
\label{it.unisa.sesa.repominer.db.entities.Method.equals(java.lang.Object)}}%end signature
}%end item
\item{ 
\index{getBelongingTypeId()}
\hypertarget{it.unisa.sesa.repominer.db.entities.Method.getBelongingTypeId()}{{\bf  getBelongingTypeId}\\}
\texttt{public java.lang.Integer\ {\bf  getBelongingTypeId}()
\label{it.unisa.sesa.repominer.db.entities.Method.getBelongingTypeId()}}%end signature
}%end item
\item{ 
\index{getId()}
\hypertarget{it.unisa.sesa.repominer.db.entities.Method.getId()}{{\bf  getId}\\}
\texttt{public java.lang.Integer\ {\bf  getId}()
\label{it.unisa.sesa.repominer.db.entities.Method.getId()}}%end signature
}%end item
\item{ 
\index{getIsConstructor()}
\hypertarget{it.unisa.sesa.repominer.db.entities.Method.getIsConstructor()}{{\bf  getIsConstructor}\\}
\texttt{public java.lang.String\ {\bf  getIsConstructor}()
\label{it.unisa.sesa.repominer.db.entities.Method.getIsConstructor()}}%end signature
}%end item
\item{ 
\index{getLineNumber()}
\hypertarget{it.unisa.sesa.repominer.db.entities.Method.getLineNumber()}{{\bf  getLineNumber}\\}
\texttt{public java.lang.Integer\ {\bf  getLineNumber}()
\label{it.unisa.sesa.repominer.db.entities.Method.getLineNumber()}}%end signature
}%end item
\item{ 
\index{getReturnTypeId()}
\hypertarget{it.unisa.sesa.repominer.db.entities.Method.getReturnTypeId()}{{\bf  getReturnTypeId}\\}
\texttt{public java.lang.Integer\ {\bf  getReturnTypeId}()
\label{it.unisa.sesa.repominer.db.entities.Method.getReturnTypeId()}}%end signature
}%end item
\item{ 
\index{hashCode()}
\hypertarget{it.unisa.sesa.repominer.db.entities.Method.hashCode()}{{\bf  hashCode}\\}
\texttt{public native int\ {\bf  hashCode}()
\label{it.unisa.sesa.repominer.db.entities.Method.hashCode()}}%end signature
}%end item
\item{ 
\index{setBelongingTypeId(Integer)}
\hypertarget{it.unisa.sesa.repominer.db.entities.Method.setBelongingTypeId(java.lang.Integer)}{{\bf  setBelongingTypeId}\\}
\texttt{public void\ {\bf  setBelongingTypeId}(\texttt{java.lang.Integer} {\bf  belongingTypeId})
\label{it.unisa.sesa.repominer.db.entities.Method.setBelongingTypeId(java.lang.Integer)}}%end signature
}%end item
\item{ 
\index{setId(Integer)}
\hypertarget{it.unisa.sesa.repominer.db.entities.Method.setId(java.lang.Integer)}{{\bf  setId}\\}
\texttt{public void\ {\bf  setId}(\texttt{java.lang.Integer} {\bf  id})
\label{it.unisa.sesa.repominer.db.entities.Method.setId(java.lang.Integer)}}%end signature
}%end item
\item{ 
\index{setIsConstructor(String)}
\hypertarget{it.unisa.sesa.repominer.db.entities.Method.setIsConstructor(java.lang.String)}{{\bf  setIsConstructor}\\}
\texttt{public void\ {\bf  setIsConstructor}(\texttt{java.lang.String} {\bf  isConstructor})
\label{it.unisa.sesa.repominer.db.entities.Method.setIsConstructor(java.lang.String)}}%end signature
}%end item
\item{ 
\index{setLineNumber(Integer)}
\hypertarget{it.unisa.sesa.repominer.db.entities.Method.setLineNumber(java.lang.Integer)}{{\bf  setLineNumber}\\}
\texttt{public void\ {\bf  setLineNumber}(\texttt{java.lang.Integer} {\bf  lineNumber})
\label{it.unisa.sesa.repominer.db.entities.Method.setLineNumber(java.lang.Integer)}}%end signature
}%end item
\item{ 
\index{setReturnTypeId(Integer)}
\hypertarget{it.unisa.sesa.repominer.db.entities.Method.setReturnTypeId(java.lang.Integer)}{{\bf  setReturnTypeId}\\}
\texttt{public void\ {\bf  setReturnTypeId}(\texttt{java.lang.Integer} {\bf  returnTypeId})
\label{it.unisa.sesa.repominer.db.entities.Method.setReturnTypeId(java.lang.Integer)}}%end signature
}%end item
\item{ 
\index{toString()}
\hypertarget{it.unisa.sesa.repominer.db.entities.Method.toString()}{{\bf  toString}\\}
\texttt{public java.lang.String\ {\bf  toString}()
\label{it.unisa.sesa.repominer.db.entities.Method.toString()}}%end signature
}%end item
\end{itemize}
}
}
\subsection{\label{it.unisa.sesa.repominer.db.entities.MethodComment}Class MethodComment}{
\hypertarget{it.unisa.sesa.repominer.db.entities.MethodComment}{}\vskip .1in 
\subsubsection{Declaration}{
\small public class MethodComment
\\ {\bf  extends} java.lang.Object
\refdefined{java.lang.Object}}
\subsubsection{Constructor summary}{
\begin{verse}
\hyperlink{it.unisa.sesa.repominer.db.entities.MethodComment()}{{\bf MethodComment()}} \\
\hyperlink{it.unisa.sesa.repominer.db.entities.MethodComment(java.lang.Integer, java.lang.Integer)}{{\bf MethodComment(Integer, Integer)}} \\
\end{verse}
}
\subsubsection{Method summary}{
\begin{verse}
\hyperlink{it.unisa.sesa.repominer.db.entities.MethodComment.equals(java.lang.Object)}{{\bf equals(Object)}} \\
\hyperlink{it.unisa.sesa.repominer.db.entities.MethodComment.getCommentId()}{{\bf getCommentId()}} \\
\hyperlink{it.unisa.sesa.repominer.db.entities.MethodComment.getMethodId()}{{\bf getMethodId()}} \\
\hyperlink{it.unisa.sesa.repominer.db.entities.MethodComment.hashCode()}{{\bf hashCode()}} \\
\hyperlink{it.unisa.sesa.repominer.db.entities.MethodComment.setCommentId(java.lang.Integer)}{{\bf setCommentId(Integer)}} \\
\hyperlink{it.unisa.sesa.repominer.db.entities.MethodComment.setMethodId(java.lang.Integer)}{{\bf setMethodId(Integer)}} \\
\hyperlink{it.unisa.sesa.repominer.db.entities.MethodComment.toString()}{{\bf toString()}} \\
\end{verse}
}
\subsubsection{Constructors}{
\vskip -2em
\begin{itemize}
\item{ 
\index{MethodComment()}
\hypertarget{it.unisa.sesa.repominer.db.entities.MethodComment()}{{\bf  MethodComment}\\}
\texttt{public\ {\bf  MethodComment}()
\label{it.unisa.sesa.repominer.db.entities.MethodComment()}}%end signature
}%end item
\item{ 
\index{MethodComment(Integer, Integer)}
\hypertarget{it.unisa.sesa.repominer.db.entities.MethodComment(java.lang.Integer, java.lang.Integer)}{{\bf  MethodComment}\\}
\texttt{public\ {\bf  MethodComment}(\texttt{java.lang.Integer} {\bf  methodId},
\texttt{java.lang.Integer} {\bf  commentId})
\label{it.unisa.sesa.repominer.db.entities.MethodComment(java.lang.Integer, java.lang.Integer)}}%end signature
}%end item
\end{itemize}
}
\subsubsection{Methods}{
\vskip -2em
\begin{itemize}
\item{ 
\index{equals(Object)}
\hypertarget{it.unisa.sesa.repominer.db.entities.MethodComment.equals(java.lang.Object)}{{\bf  equals}\\}
\texttt{public boolean\ {\bf  equals}(\texttt{java.lang.Object} {\bf  arg0})
\label{it.unisa.sesa.repominer.db.entities.MethodComment.equals(java.lang.Object)}}%end signature
}%end item
\item{ 
\index{getCommentId()}
\hypertarget{it.unisa.sesa.repominer.db.entities.MethodComment.getCommentId()}{{\bf  getCommentId}\\}
\texttt{public java.lang.Integer\ {\bf  getCommentId}()
\label{it.unisa.sesa.repominer.db.entities.MethodComment.getCommentId()}}%end signature
}%end item
\item{ 
\index{getMethodId()}
\hypertarget{it.unisa.sesa.repominer.db.entities.MethodComment.getMethodId()}{{\bf  getMethodId}\\}
\texttt{public java.lang.Integer\ {\bf  getMethodId}()
\label{it.unisa.sesa.repominer.db.entities.MethodComment.getMethodId()}}%end signature
}%end item
\item{ 
\index{hashCode()}
\hypertarget{it.unisa.sesa.repominer.db.entities.MethodComment.hashCode()}{{\bf  hashCode}\\}
\texttt{public native int\ {\bf  hashCode}()
\label{it.unisa.sesa.repominer.db.entities.MethodComment.hashCode()}}%end signature
}%end item
\item{ 
\index{setCommentId(Integer)}
\hypertarget{it.unisa.sesa.repominer.db.entities.MethodComment.setCommentId(java.lang.Integer)}{{\bf  setCommentId}\\}
\texttt{public void\ {\bf  setCommentId}(\texttt{java.lang.Integer} {\bf  commentId})
\label{it.unisa.sesa.repominer.db.entities.MethodComment.setCommentId(java.lang.Integer)}}%end signature
}%end item
\item{ 
\index{setMethodId(Integer)}
\hypertarget{it.unisa.sesa.repominer.db.entities.MethodComment.setMethodId(java.lang.Integer)}{{\bf  setMethodId}\\}
\texttt{public void\ {\bf  setMethodId}(\texttt{java.lang.Integer} {\bf  methodId})
\label{it.unisa.sesa.repominer.db.entities.MethodComment.setMethodId(java.lang.Integer)}}%end signature
}%end item
\item{ 
\index{toString()}
\hypertarget{it.unisa.sesa.repominer.db.entities.MethodComment.toString()}{{\bf  toString}\\}
\texttt{public java.lang.String\ {\bf  toString}()
\label{it.unisa.sesa.repominer.db.entities.MethodComment.toString()}}%end signature
}%end item
\end{itemize}
}
}
\subsection{\label{it.unisa.sesa.repominer.db.entities.MethodInvocation}Class MethodInvocation}{
\hypertarget{it.unisa.sesa.repominer.db.entities.MethodInvocation}{}\vskip .1in 
\subsubsection{Declaration}{
\small public class MethodInvocation
\\ {\bf  extends} java.lang.Object
\refdefined{java.lang.Object}}
\subsubsection{Constructor summary}{
\begin{verse}
\hyperlink{it.unisa.sesa.repominer.db.entities.MethodInvocation()}{{\bf MethodInvocation()}} \\
\hyperlink{it.unisa.sesa.repominer.db.entities.MethodInvocation(java.lang.Integer, java.lang.Integer)}{{\bf MethodInvocation(Integer, Integer)}} \\
\end{verse}
}
\subsubsection{Method summary}{
\begin{verse}
\hyperlink{it.unisa.sesa.repominer.db.entities.MethodInvocation.equals(java.lang.Object)}{{\bf equals(Object)}} \\
\hyperlink{it.unisa.sesa.repominer.db.entities.MethodInvocation.getInvokedMethodId()}{{\bf getInvokedMethodId()}} \\
\hyperlink{it.unisa.sesa.repominer.db.entities.MethodInvocation.getInvokerMethodId()}{{\bf getInvokerMethodId()}} \\
\hyperlink{it.unisa.sesa.repominer.db.entities.MethodInvocation.hashCode()}{{\bf hashCode()}} \\
\hyperlink{it.unisa.sesa.repominer.db.entities.MethodInvocation.setInvokedMethodId(java.lang.Integer)}{{\bf setInvokedMethodId(Integer)}} \\
\hyperlink{it.unisa.sesa.repominer.db.entities.MethodInvocation.setInvokerMethodId(java.lang.Integer)}{{\bf setInvokerMethodId(Integer)}} \\
\hyperlink{it.unisa.sesa.repominer.db.entities.MethodInvocation.toString()}{{\bf toString()}} \\
\end{verse}
}
\subsubsection{Constructors}{
\vskip -2em
\begin{itemize}
\item{ 
\index{MethodInvocation()}
\hypertarget{it.unisa.sesa.repominer.db.entities.MethodInvocation()}{{\bf  MethodInvocation}\\}
\texttt{public\ {\bf  MethodInvocation}()
\label{it.unisa.sesa.repominer.db.entities.MethodInvocation()}}%end signature
}%end item
\item{ 
\index{MethodInvocation(Integer, Integer)}
\hypertarget{it.unisa.sesa.repominer.db.entities.MethodInvocation(java.lang.Integer, java.lang.Integer)}{{\bf  MethodInvocation}\\}
\texttt{public\ {\bf  MethodInvocation}(\texttt{java.lang.Integer} {\bf  invokerMethodId},
\texttt{java.lang.Integer} {\bf  invokedMethodId})
\label{it.unisa.sesa.repominer.db.entities.MethodInvocation(java.lang.Integer, java.lang.Integer)}}%end signature
}%end item
\end{itemize}
}
\subsubsection{Methods}{
\vskip -2em
\begin{itemize}
\item{ 
\index{equals(Object)}
\hypertarget{it.unisa.sesa.repominer.db.entities.MethodInvocation.equals(java.lang.Object)}{{\bf  equals}\\}
\texttt{public boolean\ {\bf  equals}(\texttt{java.lang.Object} {\bf  arg0})
\label{it.unisa.sesa.repominer.db.entities.MethodInvocation.equals(java.lang.Object)}}%end signature
}%end item
\item{ 
\index{getInvokedMethodId()}
\hypertarget{it.unisa.sesa.repominer.db.entities.MethodInvocation.getInvokedMethodId()}{{\bf  getInvokedMethodId}\\}
\texttt{public java.lang.Integer\ {\bf  getInvokedMethodId}()
\label{it.unisa.sesa.repominer.db.entities.MethodInvocation.getInvokedMethodId()}}%end signature
}%end item
\item{ 
\index{getInvokerMethodId()}
\hypertarget{it.unisa.sesa.repominer.db.entities.MethodInvocation.getInvokerMethodId()}{{\bf  getInvokerMethodId}\\}
\texttt{public java.lang.Integer\ {\bf  getInvokerMethodId}()
\label{it.unisa.sesa.repominer.db.entities.MethodInvocation.getInvokerMethodId()}}%end signature
}%end item
\item{ 
\index{hashCode()}
\hypertarget{it.unisa.sesa.repominer.db.entities.MethodInvocation.hashCode()}{{\bf  hashCode}\\}
\texttt{public native int\ {\bf  hashCode}()
\label{it.unisa.sesa.repominer.db.entities.MethodInvocation.hashCode()}}%end signature
}%end item
\item{ 
\index{setInvokedMethodId(Integer)}
\hypertarget{it.unisa.sesa.repominer.db.entities.MethodInvocation.setInvokedMethodId(java.lang.Integer)}{{\bf  setInvokedMethodId}\\}
\texttt{public void\ {\bf  setInvokedMethodId}(\texttt{java.lang.Integer} {\bf  invokedMethodId})
\label{it.unisa.sesa.repominer.db.entities.MethodInvocation.setInvokedMethodId(java.lang.Integer)}}%end signature
}%end item
\item{ 
\index{setInvokerMethodId(Integer)}
\hypertarget{it.unisa.sesa.repominer.db.entities.MethodInvocation.setInvokerMethodId(java.lang.Integer)}{{\bf  setInvokerMethodId}\\}
\texttt{public void\ {\bf  setInvokerMethodId}(\texttt{java.lang.Integer} {\bf  invokerMethodId})
\label{it.unisa.sesa.repominer.db.entities.MethodInvocation.setInvokerMethodId(java.lang.Integer)}}%end signature
}%end item
\item{ 
\index{toString()}
\hypertarget{it.unisa.sesa.repominer.db.entities.MethodInvocation.toString()}{{\bf  toString}\\}
\texttt{public java.lang.String\ {\bf  toString}()
\label{it.unisa.sesa.repominer.db.entities.MethodInvocation.toString()}}%end signature
}%end item
\end{itemize}
}
}
\subsection{\label{it.unisa.sesa.repominer.db.entities.MethodsChangeInCommit}Class MethodsChangeInCommit}{
\hypertarget{it.unisa.sesa.repominer.db.entities.MethodsChangeInCommit}{}\vskip .1in 
\subsubsection{Declaration}{
\small public class MethodsChangeInCommit
\\ {\bf  extends} java.lang.Object
\refdefined{java.lang.Object}}
\subsubsection{Constructor summary}{
\begin{verse}
\hyperlink{it.unisa.sesa.repominer.db.entities.MethodsChangeInCommit()}{{\bf MethodsChangeInCommit()}} \\
\hyperlink{it.unisa.sesa.repominer.db.entities.MethodsChangeInCommit(java.lang.Integer, java.lang.String, java.lang.Integer)}{{\bf MethodsChangeInCommit(Integer, String, Integer)}} \\
\end{verse}
}
\subsubsection{Method summary}{
\begin{verse}
\hyperlink{it.unisa.sesa.repominer.db.entities.MethodsChangeInCommit.equals(java.lang.Object)}{{\bf equals(Object)}} \\
\hyperlink{it.unisa.sesa.repominer.db.entities.MethodsChangeInCommit.hashCode()}{{\bf hashCode()}} \\
\hyperlink{it.unisa.sesa.repominer.db.entities.MethodsChangeInCommit.toString()}{{\bf toString()}} \\
\end{verse}
}
\subsubsection{Constructors}{
\vskip -2em
\begin{itemize}
\item{ 
\index{MethodsChangeInCommit()}
\hypertarget{it.unisa.sesa.repominer.db.entities.MethodsChangeInCommit()}{{\bf  MethodsChangeInCommit}\\}
\texttt{public\ {\bf  MethodsChangeInCommit}()
\label{it.unisa.sesa.repominer.db.entities.MethodsChangeInCommit()}}%end signature
}%end item
\item{ 
\index{MethodsChangeInCommit(Integer, String, Integer)}
\hypertarget{it.unisa.sesa.repominer.db.entities.MethodsChangeInCommit(java.lang.Integer, java.lang.String, java.lang.Integer)}{{\bf  MethodsChangeInCommit}\\}
\texttt{public\ {\bf  MethodsChangeInCommit}(\texttt{java.lang.Integer} {\bf  id},
\texttt{java.lang.String} {\bf  modifiedMethod},
\texttt{java.lang.Integer} {\bf  proprietaryFileId})
\label{it.unisa.sesa.repominer.db.entities.MethodsChangeInCommit(java.lang.Integer, java.lang.String, java.lang.Integer)}}%end signature
}%end item
\end{itemize}
}
\subsubsection{Methods}{
\vskip -2em
\begin{itemize}
\item{ 
\index{equals(Object)}
\hypertarget{it.unisa.sesa.repominer.db.entities.MethodsChangeInCommit.equals(java.lang.Object)}{{\bf  equals}\\}
\texttt{public boolean\ {\bf  equals}(\texttt{java.lang.Object} {\bf  arg0})
\label{it.unisa.sesa.repominer.db.entities.MethodsChangeInCommit.equals(java.lang.Object)}}%end signature
}%end item
\item{ 
\index{hashCode()}
\hypertarget{it.unisa.sesa.repominer.db.entities.MethodsChangeInCommit.hashCode()}{{\bf  hashCode}\\}
\texttt{public native int\ {\bf  hashCode}()
\label{it.unisa.sesa.repominer.db.entities.MethodsChangeInCommit.hashCode()}}%end signature
}%end item
\item{ 
\index{toString()}
\hypertarget{it.unisa.sesa.repominer.db.entities.MethodsChangeInCommit.toString()}{{\bf  toString}\\}
\texttt{public java.lang.String\ {\bf  toString}()
\label{it.unisa.sesa.repominer.db.entities.MethodsChangeInCommit.toString()}}%end signature
}%end item
\end{itemize}
}
}
\subsection{\label{it.unisa.sesa.repominer.db.entities.Metric}Class Metric}{
\hypertarget{it.unisa.sesa.repominer.db.entities.Metric}{}\vskip .1in 
\subsubsection{Declaration}{
\small public class Metric
\\ {\bf  extends} java.lang.Object
\refdefined{java.lang.Object}}
\subsubsection{All known subclasses}{ProjectMetric\small{\refdefined{it.unisa.sesa.repominer.db.entities.ProjectMetric}}, PackageMetric\small{\refdefined{it.unisa.sesa.repominer.db.entities.PackageMetric}}, MetricMethod\small{\refdefined{it.unisa.sesa.repominer.db.entities.MetricMethod}}}
\subsubsection{Field summary}{
\begin{verse}
\hyperlink{it.unisa.sesa.repominer.db.entities.Metric.BCCM_DESCRIPTION}{{\bf BCCM\_DESCRIPTION}} \\
\hyperlink{it.unisa.sesa.repominer.db.entities.Metric.BCCM_NAME}{{\bf BCCM\_NAME}} \\
\hyperlink{it.unisa.sesa.repominer.db.entities.Metric.CHANGE_SET_SIZE_DESCRIPTION}{{\bf CHANGE\_SET\_SIZE\_DESCRIPTION}} \\
\hyperlink{it.unisa.sesa.repominer.db.entities.Metric.CHANGE_SET_SIZE_NAME}{{\bf CHANGE\_SET\_SIZE\_NAME}} \\
\hyperlink{it.unisa.sesa.repominer.db.entities.Metric.ECCM_DESCRIPTION}{{\bf ECCM\_DESCRIPTION}} \\
\hyperlink{it.unisa.sesa.repominer.db.entities.Metric.ECCM_NAME}{{\bf ECCM\_NAME}} \\
\hyperlink{it.unisa.sesa.repominer.db.entities.Metric.ECCM_STATIC_DESCRIPTION}{{\bf ECCM\_STATIC\_DESCRIPTION}} \\
\hyperlink{it.unisa.sesa.repominer.db.entities.Metric.ECCM_STATIC_NAME}{{\bf ECCM\_STATIC\_NAME}} \\
\hyperlink{it.unisa.sesa.repominer.db.entities.Metric.MAX_LINES_DESCRIPTION}{{\bf MAX\_LINES\_DESCRIPTION}} \\
\hyperlink{it.unisa.sesa.repominer.db.entities.Metric.MAX_LINES_NAME}{{\bf MAX\_LINES\_NAME}} \\
\hyperlink{it.unisa.sesa.repominer.db.entities.Metric.MEAN_LINES_DESCRIPTION}{{\bf MEAN\_LINES\_DESCRIPTION}} \\
\hyperlink{it.unisa.sesa.repominer.db.entities.Metric.MEAN_LINES_NAME}{{\bf MEAN\_LINES\_NAME}} \\
\hyperlink{it.unisa.sesa.repominer.db.entities.Metric.NAUTH_DESCRIPTION}{{\bf NAUTH\_DESCRIPTION}} \\
\hyperlink{it.unisa.sesa.repominer.db.entities.Metric.NAUTH_NAME}{{\bf NAUTH\_NAME}} \\
\hyperlink{it.unisa.sesa.repominer.db.entities.Metric.NCHANGE_DESCRIPTION}{{\bf NCHANGE\_DESCRIPTION}} \\
\hyperlink{it.unisa.sesa.repominer.db.entities.Metric.NCHANGE_NAME}{{\bf NCHANGE\_NAME}} \\
\hyperlink{it.unisa.sesa.repominer.db.entities.Metric.NFIX_DESCRIPTION}{{\bf NFIX\_DESCRIPTION}} \\
\hyperlink{it.unisa.sesa.repominer.db.entities.Metric.NFIX_NAME}{{\bf NFIX\_NAME}} \\
\hyperlink{it.unisa.sesa.repominer.db.entities.Metric.NREF_DESCRIPTION}{{\bf NREF\_DESCRIPTION}} \\
\hyperlink{it.unisa.sesa.repominer.db.entities.Metric.NREF_NAME}{{\bf NREF\_NAME}} \\
\hyperlink{it.unisa.sesa.repominer.db.entities.Metric.NUM_REVISION_DESCRIPTION}{{\bf NUM\_REVISION\_DESCRIPTION}} \\
\hyperlink{it.unisa.sesa.repominer.db.entities.Metric.NUM_REVISION_NAME}{{\bf NUM\_REVISION\_NAME}} \\
\hyperlink{it.unisa.sesa.repominer.db.entities.Metric.REVISION_DESCRIPTION}{{\bf REVISION\_DESCRIPTION}} \\
\hyperlink{it.unisa.sesa.repominer.db.entities.Metric.REVISION_NAME}{{\bf REVISION\_NAME}} \\
\hyperlink{it.unisa.sesa.repominer.db.entities.Metric.SUM_LINES_DESCRIPTION}{{\bf SUM\_LINES\_DESCRIPTION}} \\
\hyperlink{it.unisa.sesa.repominer.db.entities.Metric.SUM_LINES_NAME}{{\bf SUM\_LINES\_NAME}} \\
\end{verse}
}
\subsubsection{Constructor summary}{
\begin{verse}
\hyperlink{it.unisa.sesa.repominer.db.entities.Metric()}{{\bf Metric()}} \\
\hyperlink{it.unisa.sesa.repominer.db.entities.Metric(java.lang.Integer, java.lang.String, java.lang.String)}{{\bf Metric(Integer, String, String)}} \\
\end{verse}
}
\subsubsection{Method summary}{
\begin{verse}
\hyperlink{it.unisa.sesa.repominer.db.entities.Metric.equals(java.lang.Object)}{{\bf equals(Object)}} \\
\hyperlink{it.unisa.sesa.repominer.db.entities.Metric.getDescription()}{{\bf getDescription()}} \\
\hyperlink{it.unisa.sesa.repominer.db.entities.Metric.getId()}{{\bf getId()}} \\
\hyperlink{it.unisa.sesa.repominer.db.entities.Metric.getName()}{{\bf getName()}} \\
\hyperlink{it.unisa.sesa.repominer.db.entities.Metric.hashCode()}{{\bf hashCode()}} \\
\hyperlink{it.unisa.sesa.repominer.db.entities.Metric.setDescription(java.lang.String)}{{\bf setDescription(String)}} \\
\hyperlink{it.unisa.sesa.repominer.db.entities.Metric.setId(java.lang.Integer)}{{\bf setId(Integer)}} \\
\hyperlink{it.unisa.sesa.repominer.db.entities.Metric.setName(java.lang.String)}{{\bf setName(String)}} \\
\hyperlink{it.unisa.sesa.repominer.db.entities.Metric.toString()}{{\bf toString()}} \\
\end{verse}
}
\subsubsection{Fields}{
\begin{itemize}
\item{
\index{BCCM\_DESCRIPTION}
\label{it.unisa.sesa.repominer.db.entities.Metric.BCCM_DESCRIPTION}\hypertarget{it.unisa.sesa.repominer.db.entities.Metric.BCCM_DESCRIPTION}{public static final java.lang.String {\bf  BCCM\_DESCRIPTION}}
}
\item{
\index{BCCM\_NAME}
\label{it.unisa.sesa.repominer.db.entities.Metric.BCCM_NAME}\hypertarget{it.unisa.sesa.repominer.db.entities.Metric.BCCM_NAME}{public static final java.lang.String {\bf  BCCM\_NAME}}
}
\item{
\index{NAUTH\_NAME}
\label{it.unisa.sesa.repominer.db.entities.Metric.NAUTH_NAME}\hypertarget{it.unisa.sesa.repominer.db.entities.Metric.NAUTH_NAME}{public static final java.lang.String {\bf  NAUTH\_NAME}}
}
\item{
\index{NAUTH\_DESCRIPTION}
\label{it.unisa.sesa.repominer.db.entities.Metric.NAUTH_DESCRIPTION}\hypertarget{it.unisa.sesa.repominer.db.entities.Metric.NAUTH_DESCRIPTION}{public static final java.lang.String {\bf  NAUTH\_DESCRIPTION}}
}
\item{
\index{REVISION\_NAME}
\label{it.unisa.sesa.repominer.db.entities.Metric.REVISION_NAME}\hypertarget{it.unisa.sesa.repominer.db.entities.Metric.REVISION_NAME}{public static final java.lang.String {\bf  REVISION\_NAME}}
}
\item{
\index{REVISION\_DESCRIPTION}
\label{it.unisa.sesa.repominer.db.entities.Metric.REVISION_DESCRIPTION}\hypertarget{it.unisa.sesa.repominer.db.entities.Metric.REVISION_DESCRIPTION}{public static final java.lang.String {\bf  REVISION\_DESCRIPTION}}
}
\item{
\index{CHANGE\_SET\_SIZE\_NAME}
\label{it.unisa.sesa.repominer.db.entities.Metric.CHANGE_SET_SIZE_NAME}\hypertarget{it.unisa.sesa.repominer.db.entities.Metric.CHANGE_SET_SIZE_NAME}{public static final java.lang.String {\bf  CHANGE\_SET\_SIZE\_NAME}}
}
\item{
\index{CHANGE\_SET\_SIZE\_DESCRIPTION}
\label{it.unisa.sesa.repominer.db.entities.Metric.CHANGE_SET_SIZE_DESCRIPTION}\hypertarget{it.unisa.sesa.repominer.db.entities.Metric.CHANGE_SET_SIZE_DESCRIPTION}{public static final java.lang.String {\bf  CHANGE\_SET\_SIZE\_DESCRIPTION}}
}
\item{
\index{NCHANGE\_NAME}
\label{it.unisa.sesa.repominer.db.entities.Metric.NCHANGE_NAME}\hypertarget{it.unisa.sesa.repominer.db.entities.Metric.NCHANGE_NAME}{public static final java.lang.String {\bf  NCHANGE\_NAME}}
}
\item{
\index{NCHANGE\_DESCRIPTION}
\label{it.unisa.sesa.repominer.db.entities.Metric.NCHANGE_DESCRIPTION}\hypertarget{it.unisa.sesa.repominer.db.entities.Metric.NCHANGE_DESCRIPTION}{public static final java.lang.String {\bf  NCHANGE\_DESCRIPTION}}
}
\item{
\index{NREF\_NAME}
\label{it.unisa.sesa.repominer.db.entities.Metric.NREF_NAME}\hypertarget{it.unisa.sesa.repominer.db.entities.Metric.NREF_NAME}{public static final java.lang.String {\bf  NREF\_NAME}}
}
\item{
\index{NREF\_DESCRIPTION}
\label{it.unisa.sesa.repominer.db.entities.Metric.NREF_DESCRIPTION}\hypertarget{it.unisa.sesa.repominer.db.entities.Metric.NREF_DESCRIPTION}{public static final java.lang.String {\bf  NREF\_DESCRIPTION}}
}
\item{
\index{NFIX\_NAME}
\label{it.unisa.sesa.repominer.db.entities.Metric.NFIX_NAME}\hypertarget{it.unisa.sesa.repominer.db.entities.Metric.NFIX_NAME}{public static final java.lang.String {\bf  NFIX\_NAME}}
}
\item{
\index{NFIX\_DESCRIPTION}
\label{it.unisa.sesa.repominer.db.entities.Metric.NFIX_DESCRIPTION}\hypertarget{it.unisa.sesa.repominer.db.entities.Metric.NFIX_DESCRIPTION}{public static final java.lang.String {\bf  NFIX\_DESCRIPTION}}
}
\item{
\index{SUM\_LINES\_NAME}
\label{it.unisa.sesa.repominer.db.entities.Metric.SUM_LINES_NAME}\hypertarget{it.unisa.sesa.repominer.db.entities.Metric.SUM_LINES_NAME}{public static final java.lang.String {\bf  SUM\_LINES\_NAME}}
}
\item{
\index{SUM\_LINES\_DESCRIPTION}
\label{it.unisa.sesa.repominer.db.entities.Metric.SUM_LINES_DESCRIPTION}\hypertarget{it.unisa.sesa.repominer.db.entities.Metric.SUM_LINES_DESCRIPTION}{public static final java.lang.String {\bf  SUM\_LINES\_DESCRIPTION}}
}
\item{
\index{MEAN\_LINES\_NAME}
\label{it.unisa.sesa.repominer.db.entities.Metric.MEAN_LINES_NAME}\hypertarget{it.unisa.sesa.repominer.db.entities.Metric.MEAN_LINES_NAME}{public static final java.lang.String {\bf  MEAN\_LINES\_NAME}}
}
\item{
\index{MEAN\_LINES\_DESCRIPTION}
\label{it.unisa.sesa.repominer.db.entities.Metric.MEAN_LINES_DESCRIPTION}\hypertarget{it.unisa.sesa.repominer.db.entities.Metric.MEAN_LINES_DESCRIPTION}{public static final java.lang.String {\bf  MEAN\_LINES\_DESCRIPTION}}
}
\item{
\index{MAX\_LINES\_NAME}
\label{it.unisa.sesa.repominer.db.entities.Metric.MAX_LINES_NAME}\hypertarget{it.unisa.sesa.repominer.db.entities.Metric.MAX_LINES_NAME}{public static final java.lang.String {\bf  MAX\_LINES\_NAME}}
}
\item{
\index{MAX\_LINES\_DESCRIPTION}
\label{it.unisa.sesa.repominer.db.entities.Metric.MAX_LINES_DESCRIPTION}\hypertarget{it.unisa.sesa.repominer.db.entities.Metric.MAX_LINES_DESCRIPTION}{public static final java.lang.String {\bf  MAX\_LINES\_DESCRIPTION}}
}
\item{
\index{ECCM\_NAME}
\label{it.unisa.sesa.repominer.db.entities.Metric.ECCM_NAME}\hypertarget{it.unisa.sesa.repominer.db.entities.Metric.ECCM_NAME}{public static final java.lang.String {\bf  ECCM\_NAME}}
}
\item{
\index{ECCM\_DESCRIPTION}
\label{it.unisa.sesa.repominer.db.entities.Metric.ECCM_DESCRIPTION}\hypertarget{it.unisa.sesa.repominer.db.entities.Metric.ECCM_DESCRIPTION}{public static final java.lang.String {\bf  ECCM\_DESCRIPTION}}
}
\item{
\index{ECCM\_STATIC\_DESCRIPTION}
\label{it.unisa.sesa.repominer.db.entities.Metric.ECCM_STATIC_DESCRIPTION}\hypertarget{it.unisa.sesa.repominer.db.entities.Metric.ECCM_STATIC_DESCRIPTION}{public static final java.lang.String {\bf  ECCM\_STATIC\_DESCRIPTION}}
}
\item{
\index{ECCM\_STATIC\_NAME}
\label{it.unisa.sesa.repominer.db.entities.Metric.ECCM_STATIC_NAME}\hypertarget{it.unisa.sesa.repominer.db.entities.Metric.ECCM_STATIC_NAME}{public static final java.lang.String {\bf  ECCM\_STATIC\_NAME}}
}
\item{
\index{NUM\_REVISION\_NAME}
\label{it.unisa.sesa.repominer.db.entities.Metric.NUM_REVISION_NAME}\hypertarget{it.unisa.sesa.repominer.db.entities.Metric.NUM_REVISION_NAME}{public static final java.lang.String {\bf  NUM\_REVISION\_NAME}}
}
\item{
\index{NUM\_REVISION\_DESCRIPTION}
\label{it.unisa.sesa.repominer.db.entities.Metric.NUM_REVISION_DESCRIPTION}\hypertarget{it.unisa.sesa.repominer.db.entities.Metric.NUM_REVISION_DESCRIPTION}{public static final java.lang.String {\bf  NUM\_REVISION\_DESCRIPTION}}
}
\end{itemize}
}
\subsubsection{Constructors}{
\vskip -2em
\begin{itemize}
\item{ 
\index{Metric()}
\hypertarget{it.unisa.sesa.repominer.db.entities.Metric()}{{\bf  Metric}\\}
\texttt{public\ {\bf  Metric}()
\label{it.unisa.sesa.repominer.db.entities.Metric()}}%end signature
}%end item
\item{ 
\index{Metric(Integer, String, String)}
\hypertarget{it.unisa.sesa.repominer.db.entities.Metric(java.lang.Integer, java.lang.String, java.lang.String)}{{\bf  Metric}\\}
\texttt{public\ {\bf  Metric}(\texttt{java.lang.Integer} {\bf  id},
\texttt{java.lang.String} {\bf  name},
\texttt{java.lang.String} {\bf  description})
\label{it.unisa.sesa.repominer.db.entities.Metric(java.lang.Integer, java.lang.String, java.lang.String)}}%end signature
}%end item
\end{itemize}
}
\subsubsection{Methods}{
\vskip -2em
\begin{itemize}
\item{ 
\index{equals(Object)}
\hypertarget{it.unisa.sesa.repominer.db.entities.Metric.equals(java.lang.Object)}{{\bf  equals}\\}
\texttt{public boolean\ {\bf  equals}(\texttt{java.lang.Object} {\bf  arg0})
\label{it.unisa.sesa.repominer.db.entities.Metric.equals(java.lang.Object)}}%end signature
}%end item
\item{ 
\index{getDescription()}
\hypertarget{it.unisa.sesa.repominer.db.entities.Metric.getDescription()}{{\bf  getDescription}\\}
\texttt{public java.lang.String\ {\bf  getDescription}()
\label{it.unisa.sesa.repominer.db.entities.Metric.getDescription()}}%end signature
}%end item
\item{ 
\index{getId()}
\hypertarget{it.unisa.sesa.repominer.db.entities.Metric.getId()}{{\bf  getId}\\}
\texttt{public java.lang.Integer\ {\bf  getId}()
\label{it.unisa.sesa.repominer.db.entities.Metric.getId()}}%end signature
}%end item
\item{ 
\index{getName()}
\hypertarget{it.unisa.sesa.repominer.db.entities.Metric.getName()}{{\bf  getName}\\}
\texttt{public java.lang.String\ {\bf  getName}()
\label{it.unisa.sesa.repominer.db.entities.Metric.getName()}}%end signature
}%end item
\item{ 
\index{hashCode()}
\hypertarget{it.unisa.sesa.repominer.db.entities.Metric.hashCode()}{{\bf  hashCode}\\}
\texttt{public native int\ {\bf  hashCode}()
\label{it.unisa.sesa.repominer.db.entities.Metric.hashCode()}}%end signature
}%end item
\item{ 
\index{setDescription(String)}
\hypertarget{it.unisa.sesa.repominer.db.entities.Metric.setDescription(java.lang.String)}{{\bf  setDescription}\\}
\texttt{public void\ {\bf  setDescription}(\texttt{java.lang.String} {\bf  description})
\label{it.unisa.sesa.repominer.db.entities.Metric.setDescription(java.lang.String)}}%end signature
}%end item
\item{ 
\index{setId(Integer)}
\hypertarget{it.unisa.sesa.repominer.db.entities.Metric.setId(java.lang.Integer)}{{\bf  setId}\\}
\texttt{public void\ {\bf  setId}(\texttt{java.lang.Integer} {\bf  id})
\label{it.unisa.sesa.repominer.db.entities.Metric.setId(java.lang.Integer)}}%end signature
}%end item
\item{ 
\index{setName(String)}
\hypertarget{it.unisa.sesa.repominer.db.entities.Metric.setName(java.lang.String)}{{\bf  setName}\\}
\texttt{public void\ {\bf  setName}(\texttt{java.lang.String} {\bf  name})
\label{it.unisa.sesa.repominer.db.entities.Metric.setName(java.lang.String)}}%end signature
}%end item
\item{ 
\index{toString()}
\hypertarget{it.unisa.sesa.repominer.db.entities.Metric.toString()}{{\bf  toString}\\}
\texttt{public java.lang.String\ {\bf  toString}()
\label{it.unisa.sesa.repominer.db.entities.Metric.toString()}}%end signature
}%end item
\end{itemize}
}
}
\subsection{\label{it.unisa.sesa.repominer.db.entities.MetricMethod}Class MetricMethod}{
\hypertarget{it.unisa.sesa.repominer.db.entities.MetricMethod}{}\vskip .1in 
\subsubsection{Declaration}{
\small public class MetricMethod
\\ {\bf  extends} it.unisa.sesa.repominer.db.entities.Metric
\refdefined{it.unisa.sesa.repominer.db.entities.Metric}}
\subsubsection{Constructor summary}{
\begin{verse}
\hyperlink{it.unisa.sesa.repominer.db.entities.MetricMethod()}{{\bf MetricMethod()}} \\
\hyperlink{it.unisa.sesa.repominer.db.entities.MetricMethod(java.lang.Integer, java.lang.Integer, java.lang.String)}{{\bf MetricMethod(Integer, Integer, String)}} \\
\end{verse}
}
\subsubsection{Method summary}{
\begin{verse}
\hyperlink{it.unisa.sesa.repominer.db.entities.MetricMethod.equals(java.lang.Object)}{{\bf equals(Object)}} \\
\hyperlink{it.unisa.sesa.repominer.db.entities.MetricMethod.getMethodId()}{{\bf getMethodId()}} \\
\hyperlink{it.unisa.sesa.repominer.db.entities.MetricMethod.getMetricId()}{{\bf getMetricId()}} \\
\hyperlink{it.unisa.sesa.repominer.db.entities.MetricMethod.getValue()}{{\bf getValue()}} \\
\hyperlink{it.unisa.sesa.repominer.db.entities.MetricMethod.hashCode()}{{\bf hashCode()}} \\
\hyperlink{it.unisa.sesa.repominer.db.entities.MetricMethod.setMethodId(java.lang.Integer)}{{\bf setMethodId(Integer)}} \\
\hyperlink{it.unisa.sesa.repominer.db.entities.MetricMethod.setMetricId(java.lang.Integer)}{{\bf setMetricId(Integer)}} \\
\hyperlink{it.unisa.sesa.repominer.db.entities.MetricMethod.setValue(java.lang.String)}{{\bf setValue(String)}} \\
\hyperlink{it.unisa.sesa.repominer.db.entities.MetricMethod.toString()}{{\bf toString()}} \\
\end{verse}
}
\subsubsection{Constructors}{
\vskip -2em
\begin{itemize}
\item{ 
\index{MetricMethod()}
\hypertarget{it.unisa.sesa.repominer.db.entities.MetricMethod()}{{\bf  MetricMethod}\\}
\texttt{public\ {\bf  MetricMethod}()
\label{it.unisa.sesa.repominer.db.entities.MetricMethod()}}%end signature
}%end item
\item{ 
\index{MetricMethod(Integer, Integer, String)}
\hypertarget{it.unisa.sesa.repominer.db.entities.MetricMethod(java.lang.Integer, java.lang.Integer, java.lang.String)}{{\bf  MetricMethod}\\}
\texttt{public\ {\bf  MetricMethod}(\texttt{java.lang.Integer} {\bf  methodId},
\texttt{java.lang.Integer} {\bf  metricId},
\texttt{java.lang.String} {\bf  value})
\label{it.unisa.sesa.repominer.db.entities.MetricMethod(java.lang.Integer, java.lang.Integer, java.lang.String)}}%end signature
}%end item
\end{itemize}
}
\subsubsection{Methods}{
\vskip -2em
\begin{itemize}
\item{ 
\index{equals(Object)}
\hypertarget{it.unisa.sesa.repominer.db.entities.MetricMethod.equals(java.lang.Object)}{{\bf  equals}\\}
\texttt{public boolean\ {\bf  equals}(\texttt{java.lang.Object} {\bf  arg0})
\label{it.unisa.sesa.repominer.db.entities.MetricMethod.equals(java.lang.Object)}}%end signature
}%end item
\item{ 
\index{getMethodId()}
\hypertarget{it.unisa.sesa.repominer.db.entities.MetricMethod.getMethodId()}{{\bf  getMethodId}\\}
\texttt{public java.lang.Integer\ {\bf  getMethodId}()
\label{it.unisa.sesa.repominer.db.entities.MetricMethod.getMethodId()}}%end signature
}%end item
\item{ 
\index{getMetricId()}
\hypertarget{it.unisa.sesa.repominer.db.entities.MetricMethod.getMetricId()}{{\bf  getMetricId}\\}
\texttt{public java.lang.Integer\ {\bf  getMetricId}()
\label{it.unisa.sesa.repominer.db.entities.MetricMethod.getMetricId()}}%end signature
}%end item
\item{ 
\index{getValue()}
\hypertarget{it.unisa.sesa.repominer.db.entities.MetricMethod.getValue()}{{\bf  getValue}\\}
\texttt{public java.lang.String\ {\bf  getValue}()
\label{it.unisa.sesa.repominer.db.entities.MetricMethod.getValue()}}%end signature
}%end item
\item{ 
\index{hashCode()}
\hypertarget{it.unisa.sesa.repominer.db.entities.MetricMethod.hashCode()}{{\bf  hashCode}\\}
\texttt{public native int\ {\bf  hashCode}()
\label{it.unisa.sesa.repominer.db.entities.MetricMethod.hashCode()}}%end signature
}%end item
\item{ 
\index{setMethodId(Integer)}
\hypertarget{it.unisa.sesa.repominer.db.entities.MetricMethod.setMethodId(java.lang.Integer)}{{\bf  setMethodId}\\}
\texttt{public void\ {\bf  setMethodId}(\texttt{java.lang.Integer} {\bf  methodId})
\label{it.unisa.sesa.repominer.db.entities.MetricMethod.setMethodId(java.lang.Integer)}}%end signature
}%end item
\item{ 
\index{setMetricId(Integer)}
\hypertarget{it.unisa.sesa.repominer.db.entities.MetricMethod.setMetricId(java.lang.Integer)}{{\bf  setMetricId}\\}
\texttt{public void\ {\bf  setMetricId}(\texttt{java.lang.Integer} {\bf  metricId})
\label{it.unisa.sesa.repominer.db.entities.MetricMethod.setMetricId(java.lang.Integer)}}%end signature
}%end item
\item{ 
\index{setValue(String)}
\hypertarget{it.unisa.sesa.repominer.db.entities.MetricMethod.setValue(java.lang.String)}{{\bf  setValue}\\}
\texttt{public void\ {\bf  setValue}(\texttt{java.lang.String} {\bf  value})
\label{it.unisa.sesa.repominer.db.entities.MetricMethod.setValue(java.lang.String)}}%end signature
}%end item
\item{ 
\index{toString()}
\hypertarget{it.unisa.sesa.repominer.db.entities.MetricMethod.toString()}{{\bf  toString}\\}
\texttt{public java.lang.String\ {\bf  toString}()
\label{it.unisa.sesa.repominer.db.entities.MetricMethod.toString()}}%end signature
}%end item
\end{itemize}
}
\subsubsection{Members inherited from class Metric }{
\texttt{it.unisa.sesa.repominer.db.entities.Metric} {\small 
\refdefined{it.unisa.sesa.repominer.db.entities.Metric}}
{\small 

\vskip -2em
\begin{itemize}
\item{\vskip -1.5ex 
\texttt{public static final {\bf  BCCM\_DESCRIPTION}}%end signature
}%end item
\item{\vskip -1.5ex 
\texttt{public static final {\bf  BCCM\_NAME}}%end signature
}%end item
\item{\vskip -1.5ex 
\texttt{public static final {\bf  CHANGE\_SET\_SIZE\_DESCRIPTION}}%end signature
}%end item
\item{\vskip -1.5ex 
\texttt{public static final {\bf  CHANGE\_SET\_SIZE\_NAME}}%end signature
}%end item
\item{\vskip -1.5ex 
\texttt{public static final {\bf  ECCM\_DESCRIPTION}}%end signature
}%end item
\item{\vskip -1.5ex 
\texttt{public static final {\bf  ECCM\_NAME}}%end signature
}%end item
\item{\vskip -1.5ex 
\texttt{public static final {\bf  ECCM\_STATIC\_DESCRIPTION}}%end signature
}%end item
\item{\vskip -1.5ex 
\texttt{public static final {\bf  ECCM\_STATIC\_NAME}}%end signature
}%end item
\item{\vskip -1.5ex 
\texttt{public boolean {\bf  equals}(\texttt{java.lang.Object} {\bf  obj})
}%end signature
}%end item
\item{\vskip -1.5ex 
\texttt{public String {\bf  getDescription}()
}%end signature
}%end item
\item{\vskip -1.5ex 
\texttt{public Integer {\bf  getId}()
}%end signature
}%end item
\item{\vskip -1.5ex 
\texttt{public String {\bf  getName}()
}%end signature
}%end item
\item{\vskip -1.5ex 
\texttt{public int {\bf  hashCode}()
}%end signature
}%end item
\item{\vskip -1.5ex 
\texttt{public static final {\bf  MAX\_LINES\_DESCRIPTION}}%end signature
}%end item
\item{\vskip -1.5ex 
\texttt{public static final {\bf  MAX\_LINES\_NAME}}%end signature
}%end item
\item{\vskip -1.5ex 
\texttt{public static final {\bf  MEAN\_LINES\_DESCRIPTION}}%end signature
}%end item
\item{\vskip -1.5ex 
\texttt{public static final {\bf  MEAN\_LINES\_NAME}}%end signature
}%end item
\item{\vskip -1.5ex 
\texttt{public static final {\bf  NAUTH\_DESCRIPTION}}%end signature
}%end item
\item{\vskip -1.5ex 
\texttt{public static final {\bf  NAUTH\_NAME}}%end signature
}%end item
\item{\vskip -1.5ex 
\texttt{public static final {\bf  NCHANGE\_DESCRIPTION}}%end signature
}%end item
\item{\vskip -1.5ex 
\texttt{public static final {\bf  NCHANGE\_NAME}}%end signature
}%end item
\item{\vskip -1.5ex 
\texttt{public static final {\bf  NFIX\_DESCRIPTION}}%end signature
}%end item
\item{\vskip -1.5ex 
\texttt{public static final {\bf  NFIX\_NAME}}%end signature
}%end item
\item{\vskip -1.5ex 
\texttt{public static final {\bf  NREF\_DESCRIPTION}}%end signature
}%end item
\item{\vskip -1.5ex 
\texttt{public static final {\bf  NREF\_NAME}}%end signature
}%end item
\item{\vskip -1.5ex 
\texttt{public static final {\bf  NUM\_REVISION\_DESCRIPTION}}%end signature
}%end item
\item{\vskip -1.5ex 
\texttt{public static final {\bf  NUM\_REVISION\_NAME}}%end signature
}%end item
\item{\vskip -1.5ex 
\texttt{public static final {\bf  REVISION\_DESCRIPTION}}%end signature
}%end item
\item{\vskip -1.5ex 
\texttt{public static final {\bf  REVISION\_NAME}}%end signature
}%end item
\item{\vskip -1.5ex 
\texttt{public void {\bf  setDescription}(\texttt{java.lang.String} {\bf  description})
}%end signature
}%end item
\item{\vskip -1.5ex 
\texttt{public void {\bf  setId}(\texttt{java.lang.Integer} {\bf  id})
}%end signature
}%end item
\item{\vskip -1.5ex 
\texttt{public void {\bf  setName}(\texttt{java.lang.String} {\bf  name})
}%end signature
}%end item
\item{\vskip -1.5ex 
\texttt{public static final {\bf  SUM\_LINES\_DESCRIPTION}}%end signature
}%end item
\item{\vskip -1.5ex 
\texttt{public static final {\bf  SUM\_LINES\_NAME}}%end signature
}%end item
\item{\vskip -1.5ex 
\texttt{public String {\bf  toString}()
}%end signature
}%end item
\end{itemize}
}
}
\subsection{\label{it.unisa.sesa.repominer.db.entities.PackageMetric}Class PackageMetric}{
\hypertarget{it.unisa.sesa.repominer.db.entities.PackageMetric}{}\vskip .1in 
\subsubsection{Declaration}{
\small public class PackageMetric
\\ {\bf  extends} it.unisa.sesa.repominer.db.entities.Metric
\refdefined{it.unisa.sesa.repominer.db.entities.Metric}}
\subsubsection{Constructor summary}{
\begin{verse}
\hyperlink{it.unisa.sesa.repominer.db.entities.PackageMetric()}{{\bf PackageMetric()}} \\
\hyperlink{it.unisa.sesa.repominer.db.entities.PackageMetric(java.lang.Integer, java.lang.String, java.lang.String, java.lang.Integer, java.lang.Integer, java.lang.Double, java.util.Date, java.util.Date)}{{\bf PackageMetric(Integer, String, String, Integer, Integer, Double, Date, Date)}} \\
\end{verse}
}
\subsubsection{Method summary}{
\begin{verse}
\hyperlink{it.unisa.sesa.repominer.db.entities.PackageMetric.equals(java.lang.Object)}{{\bf equals(Object)}} \\
\hyperlink{it.unisa.sesa.repominer.db.entities.PackageMetric.getEnd()}{{\bf getEnd()}} \\
\hyperlink{it.unisa.sesa.repominer.db.entities.PackageMetric.getMetricId()}{{\bf getMetricId()}} \\
\hyperlink{it.unisa.sesa.repominer.db.entities.PackageMetric.getPackageId()}{{\bf getPackageId()}} \\
\hyperlink{it.unisa.sesa.repominer.db.entities.PackageMetric.getStart()}{{\bf getStart()}} \\
\hyperlink{it.unisa.sesa.repominer.db.entities.PackageMetric.getValue()}{{\bf getValue()}} \\
\hyperlink{it.unisa.sesa.repominer.db.entities.PackageMetric.hashCode()}{{\bf hashCode()}} \\
\hyperlink{it.unisa.sesa.repominer.db.entities.PackageMetric.setEnd(java.util.Date)}{{\bf setEnd(Date)}} \\
\hyperlink{it.unisa.sesa.repominer.db.entities.PackageMetric.setMetricId(java.lang.Integer)}{{\bf setMetricId(Integer)}} \\
\hyperlink{it.unisa.sesa.repominer.db.entities.PackageMetric.setPackageId(java.lang.Integer)}{{\bf setPackageId(Integer)}} \\
\hyperlink{it.unisa.sesa.repominer.db.entities.PackageMetric.setStart(java.util.Date)}{{\bf setStart(Date)}} \\
\hyperlink{it.unisa.sesa.repominer.db.entities.PackageMetric.setValue(java.lang.Double)}{{\bf setValue(Double)}} \\
\hyperlink{it.unisa.sesa.repominer.db.entities.PackageMetric.toString()}{{\bf toString()}} \\
\end{verse}
}
\subsubsection{Constructors}{
\vskip -2em
\begin{itemize}
\item{ 
\index{PackageMetric()}
\hypertarget{it.unisa.sesa.repominer.db.entities.PackageMetric()}{{\bf  PackageMetric}\\}
\texttt{public\ {\bf  PackageMetric}()
\label{it.unisa.sesa.repominer.db.entities.PackageMetric()}}%end signature
}%end item
\item{ 
\index{PackageMetric(Integer, String, String, Integer, Integer, Double, Date, Date)}
\hypertarget{it.unisa.sesa.repominer.db.entities.PackageMetric(java.lang.Integer, java.lang.String, java.lang.String, java.lang.Integer, java.lang.Integer, java.lang.Double, java.util.Date, java.util.Date)}{{\bf  PackageMetric}\\}
\texttt{public\ {\bf  PackageMetric}(\texttt{java.lang.Integer} {\bf  id},
\texttt{java.lang.String} {\bf  name},
\texttt{java.lang.String} {\bf  description},
\texttt{java.lang.Integer} {\bf  packageId},
\texttt{java.lang.Integer} {\bf  metricId},
\texttt{java.lang.Double} {\bf  value},
\texttt{java.util.Date} {\bf  startDate},
\texttt{java.util.Date} {\bf  endDate})
\label{it.unisa.sesa.repominer.db.entities.PackageMetric(java.lang.Integer, java.lang.String, java.lang.String, java.lang.Integer, java.lang.Integer, java.lang.Double, java.util.Date, java.util.Date)}}%end signature
}%end item
\end{itemize}
}
\subsubsection{Methods}{
\vskip -2em
\begin{itemize}
\item{ 
\index{equals(Object)}
\hypertarget{it.unisa.sesa.repominer.db.entities.PackageMetric.equals(java.lang.Object)}{{\bf  equals}\\}
\texttt{public boolean\ {\bf  equals}(\texttt{java.lang.Object} {\bf  arg0})
\label{it.unisa.sesa.repominer.db.entities.PackageMetric.equals(java.lang.Object)}}%end signature
}%end item
\item{ 
\index{getEnd()}
\hypertarget{it.unisa.sesa.repominer.db.entities.PackageMetric.getEnd()}{{\bf  getEnd}\\}
\texttt{public java.util.Date\ {\bf  getEnd}()
\label{it.unisa.sesa.repominer.db.entities.PackageMetric.getEnd()}}%end signature
}%end item
\item{ 
\index{getMetricId()}
\hypertarget{it.unisa.sesa.repominer.db.entities.PackageMetric.getMetricId()}{{\bf  getMetricId}\\}
\texttt{public java.lang.Integer\ {\bf  getMetricId}()
\label{it.unisa.sesa.repominer.db.entities.PackageMetric.getMetricId()}}%end signature
}%end item
\item{ 
\index{getPackageId()}
\hypertarget{it.unisa.sesa.repominer.db.entities.PackageMetric.getPackageId()}{{\bf  getPackageId}\\}
\texttt{public java.lang.Integer\ {\bf  getPackageId}()
\label{it.unisa.sesa.repominer.db.entities.PackageMetric.getPackageId()}}%end signature
}%end item
\item{ 
\index{getStart()}
\hypertarget{it.unisa.sesa.repominer.db.entities.PackageMetric.getStart()}{{\bf  getStart}\\}
\texttt{public java.util.Date\ {\bf  getStart}()
\label{it.unisa.sesa.repominer.db.entities.PackageMetric.getStart()}}%end signature
}%end item
\item{ 
\index{getValue()}
\hypertarget{it.unisa.sesa.repominer.db.entities.PackageMetric.getValue()}{{\bf  getValue}\\}
\texttt{public java.lang.Double\ {\bf  getValue}()
\label{it.unisa.sesa.repominer.db.entities.PackageMetric.getValue()}}%end signature
}%end item
\item{ 
\index{hashCode()}
\hypertarget{it.unisa.sesa.repominer.db.entities.PackageMetric.hashCode()}{{\bf  hashCode}\\}
\texttt{public native int\ {\bf  hashCode}()
\label{it.unisa.sesa.repominer.db.entities.PackageMetric.hashCode()}}%end signature
}%end item
\item{ 
\index{setEnd(Date)}
\hypertarget{it.unisa.sesa.repominer.db.entities.PackageMetric.setEnd(java.util.Date)}{{\bf  setEnd}\\}
\texttt{public void\ {\bf  setEnd}(\texttt{java.util.Date} {\bf  endDate})
\label{it.unisa.sesa.repominer.db.entities.PackageMetric.setEnd(java.util.Date)}}%end signature
}%end item
\item{ 
\index{setMetricId(Integer)}
\hypertarget{it.unisa.sesa.repominer.db.entities.PackageMetric.setMetricId(java.lang.Integer)}{{\bf  setMetricId}\\}
\texttt{public void\ {\bf  setMetricId}(\texttt{java.lang.Integer} {\bf  metricId})
\label{it.unisa.sesa.repominer.db.entities.PackageMetric.setMetricId(java.lang.Integer)}}%end signature
}%end item
\item{ 
\index{setPackageId(Integer)}
\hypertarget{it.unisa.sesa.repominer.db.entities.PackageMetric.setPackageId(java.lang.Integer)}{{\bf  setPackageId}\\}
\texttt{public void\ {\bf  setPackageId}(\texttt{java.lang.Integer} {\bf  packageId})
\label{it.unisa.sesa.repominer.db.entities.PackageMetric.setPackageId(java.lang.Integer)}}%end signature
}%end item
\item{ 
\index{setStart(Date)}
\hypertarget{it.unisa.sesa.repominer.db.entities.PackageMetric.setStart(java.util.Date)}{{\bf  setStart}\\}
\texttt{public void\ {\bf  setStart}(\texttt{java.util.Date} {\bf  startDate})
\label{it.unisa.sesa.repominer.db.entities.PackageMetric.setStart(java.util.Date)}}%end signature
}%end item
\item{ 
\index{setValue(Double)}
\hypertarget{it.unisa.sesa.repominer.db.entities.PackageMetric.setValue(java.lang.Double)}{{\bf  setValue}\\}
\texttt{public void\ {\bf  setValue}(\texttt{java.lang.Double} {\bf  value})
\label{it.unisa.sesa.repominer.db.entities.PackageMetric.setValue(java.lang.Double)}}%end signature
}%end item
\item{ 
\index{toString()}
\hypertarget{it.unisa.sesa.repominer.db.entities.PackageMetric.toString()}{{\bf  toString}\\}
\texttt{public java.lang.String\ {\bf  toString}()
\label{it.unisa.sesa.repominer.db.entities.PackageMetric.toString()}}%end signature
}%end item
\end{itemize}
}
\subsubsection{Members inherited from class Metric }{
\texttt{it.unisa.sesa.repominer.db.entities.Metric} {\small 
\refdefined{it.unisa.sesa.repominer.db.entities.Metric}}
{\small 

\vskip -2em
\begin{itemize}
\item{\vskip -1.5ex 
\texttt{public static final {\bf  BCCM\_DESCRIPTION}}%end signature
}%end item
\item{\vskip -1.5ex 
\texttt{public static final {\bf  BCCM\_NAME}}%end signature
}%end item
\item{\vskip -1.5ex 
\texttt{public static final {\bf  CHANGE\_SET\_SIZE\_DESCRIPTION}}%end signature
}%end item
\item{\vskip -1.5ex 
\texttt{public static final {\bf  CHANGE\_SET\_SIZE\_NAME}}%end signature
}%end item
\item{\vskip -1.5ex 
\texttt{public static final {\bf  ECCM\_DESCRIPTION}}%end signature
}%end item
\item{\vskip -1.5ex 
\texttt{public static final {\bf  ECCM\_NAME}}%end signature
}%end item
\item{\vskip -1.5ex 
\texttt{public static final {\bf  ECCM\_STATIC\_DESCRIPTION}}%end signature
}%end item
\item{\vskip -1.5ex 
\texttt{public static final {\bf  ECCM\_STATIC\_NAME}}%end signature
}%end item
\item{\vskip -1.5ex 
\texttt{public boolean {\bf  equals}(\texttt{java.lang.Object} {\bf  obj})
}%end signature
}%end item
\item{\vskip -1.5ex 
\texttt{public String {\bf  getDescription}()
}%end signature
}%end item
\item{\vskip -1.5ex 
\texttt{public Integer {\bf  getId}()
}%end signature
}%end item
\item{\vskip -1.5ex 
\texttt{public String {\bf  getName}()
}%end signature
}%end item
\item{\vskip -1.5ex 
\texttt{public int {\bf  hashCode}()
}%end signature
}%end item
\item{\vskip -1.5ex 
\texttt{public static final {\bf  MAX\_LINES\_DESCRIPTION}}%end signature
}%end item
\item{\vskip -1.5ex 
\texttt{public static final {\bf  MAX\_LINES\_NAME}}%end signature
}%end item
\item{\vskip -1.5ex 
\texttt{public static final {\bf  MEAN\_LINES\_DESCRIPTION}}%end signature
}%end item
\item{\vskip -1.5ex 
\texttt{public static final {\bf  MEAN\_LINES\_NAME}}%end signature
}%end item
\item{\vskip -1.5ex 
\texttt{public static final {\bf  NAUTH\_DESCRIPTION}}%end signature
}%end item
\item{\vskip -1.5ex 
\texttt{public static final {\bf  NAUTH\_NAME}}%end signature
}%end item
\item{\vskip -1.5ex 
\texttt{public static final {\bf  NCHANGE\_DESCRIPTION}}%end signature
}%end item
\item{\vskip -1.5ex 
\texttt{public static final {\bf  NCHANGE\_NAME}}%end signature
}%end item
\item{\vskip -1.5ex 
\texttt{public static final {\bf  NFIX\_DESCRIPTION}}%end signature
}%end item
\item{\vskip -1.5ex 
\texttt{public static final {\bf  NFIX\_NAME}}%end signature
}%end item
\item{\vskip -1.5ex 
\texttt{public static final {\bf  NREF\_DESCRIPTION}}%end signature
}%end item
\item{\vskip -1.5ex 
\texttt{public static final {\bf  NREF\_NAME}}%end signature
}%end item
\item{\vskip -1.5ex 
\texttt{public static final {\bf  NUM\_REVISION\_DESCRIPTION}}%end signature
}%end item
\item{\vskip -1.5ex 
\texttt{public static final {\bf  NUM\_REVISION\_NAME}}%end signature
}%end item
\item{\vskip -1.5ex 
\texttt{public static final {\bf  REVISION\_DESCRIPTION}}%end signature
}%end item
\item{\vskip -1.5ex 
\texttt{public static final {\bf  REVISION\_NAME}}%end signature
}%end item
\item{\vskip -1.5ex 
\texttt{public void {\bf  setDescription}(\texttt{java.lang.String} {\bf  description})
}%end signature
}%end item
\item{\vskip -1.5ex 
\texttt{public void {\bf  setId}(\texttt{java.lang.Integer} {\bf  id})
}%end signature
}%end item
\item{\vskip -1.5ex 
\texttt{public void {\bf  setName}(\texttt{java.lang.String} {\bf  name})
}%end signature
}%end item
\item{\vskip -1.5ex 
\texttt{public static final {\bf  SUM\_LINES\_DESCRIPTION}}%end signature
}%end item
\item{\vskip -1.5ex 
\texttt{public static final {\bf  SUM\_LINES\_NAME}}%end signature
}%end item
\item{\vskip -1.5ex 
\texttt{public String {\bf  toString}()
}%end signature
}%end item
\end{itemize}
}
}
\subsection{\label{it.unisa.sesa.repominer.db.entities.Project}Class Project}{
\hypertarget{it.unisa.sesa.repominer.db.entities.Project}{}\vskip .1in 
\subsubsection{Declaration}{
\small public class Project
\\ {\bf  extends} java.lang.Object
\refdefined{java.lang.Object}}
\subsubsection{Constructor summary}{
\begin{verse}
\hyperlink{it.unisa.sesa.repominer.db.entities.Project()}{{\bf Project()}} \\
\hyperlink{it.unisa.sesa.repominer.db.entities.Project(java.lang.Integer, java.lang.String, java.lang.String, java.lang.String)}{{\bf Project(Integer, String, String, String)}} \\
\end{verse}
}
\subsubsection{Method summary}{
\begin{verse}
\hyperlink{it.unisa.sesa.repominer.db.entities.Project.equals(java.lang.Object)}{{\bf equals(Object)}} \\
\hyperlink{it.unisa.sesa.repominer.db.entities.Project.getBugtrackerUrl()}{{\bf getBugtrackerUrl()}} \\
\hyperlink{it.unisa.sesa.repominer.db.entities.Project.getId()}{{\bf getId()}} \\
\hyperlink{it.unisa.sesa.repominer.db.entities.Project.getName()}{{\bf getName()}} \\
\hyperlink{it.unisa.sesa.repominer.db.entities.Project.getVersioningUrl()}{{\bf getVersioningUrl()}} \\
\hyperlink{it.unisa.sesa.repominer.db.entities.Project.hashCode()}{{\bf hashCode()}} \\
\hyperlink{it.unisa.sesa.repominer.db.entities.Project.setBugtrackerUrl(java.lang.String)}{{\bf setBugtrackerUrl(String)}} \\
\hyperlink{it.unisa.sesa.repominer.db.entities.Project.setId(java.lang.Integer)}{{\bf setId(Integer)}} \\
\hyperlink{it.unisa.sesa.repominer.db.entities.Project.setName(java.lang.String)}{{\bf setName(String)}} \\
\hyperlink{it.unisa.sesa.repominer.db.entities.Project.setVersioningUrl(java.lang.String)}{{\bf setVersioningUrl(String)}} \\
\hyperlink{it.unisa.sesa.repominer.db.entities.Project.toString()}{{\bf toString()}} \\
\end{verse}
}
\subsubsection{Constructors}{
\vskip -2em
\begin{itemize}
\item{ 
\index{Project()}
\hypertarget{it.unisa.sesa.repominer.db.entities.Project()}{{\bf  Project}\\}
\texttt{public\ {\bf  Project}()
\label{it.unisa.sesa.repominer.db.entities.Project()}}%end signature
}%end item
\item{ 
\index{Project(Integer, String, String, String)}
\hypertarget{it.unisa.sesa.repominer.db.entities.Project(java.lang.Integer, java.lang.String, java.lang.String, java.lang.String)}{{\bf  Project}\\}
\texttt{public\ {\bf  Project}(\texttt{java.lang.Integer} {\bf  id},
\texttt{java.lang.String} {\bf  name},
\texttt{java.lang.String} {\bf  versioningUrl},
\texttt{java.lang.String} {\bf  bugtrackerUrl})
\label{it.unisa.sesa.repominer.db.entities.Project(java.lang.Integer, java.lang.String, java.lang.String, java.lang.String)}}%end signature
}%end item
\end{itemize}
}
\subsubsection{Methods}{
\vskip -2em
\begin{itemize}
\item{ 
\index{equals(Object)}
\hypertarget{it.unisa.sesa.repominer.db.entities.Project.equals(java.lang.Object)}{{\bf  equals}\\}
\texttt{public boolean\ {\bf  equals}(\texttt{java.lang.Object} {\bf  arg0})
\label{it.unisa.sesa.repominer.db.entities.Project.equals(java.lang.Object)}}%end signature
}%end item
\item{ 
\index{getBugtrackerUrl()}
\hypertarget{it.unisa.sesa.repominer.db.entities.Project.getBugtrackerUrl()}{{\bf  getBugtrackerUrl}\\}
\texttt{public java.lang.String\ {\bf  getBugtrackerUrl}()
\label{it.unisa.sesa.repominer.db.entities.Project.getBugtrackerUrl()}}%end signature
}%end item
\item{ 
\index{getId()}
\hypertarget{it.unisa.sesa.repominer.db.entities.Project.getId()}{{\bf  getId}\\}
\texttt{public java.lang.Integer\ {\bf  getId}()
\label{it.unisa.sesa.repominer.db.entities.Project.getId()}}%end signature
}%end item
\item{ 
\index{getName()}
\hypertarget{it.unisa.sesa.repominer.db.entities.Project.getName()}{{\bf  getName}\\}
\texttt{public java.lang.String\ {\bf  getName}()
\label{it.unisa.sesa.repominer.db.entities.Project.getName()}}%end signature
}%end item
\item{ 
\index{getVersioningUrl()}
\hypertarget{it.unisa.sesa.repominer.db.entities.Project.getVersioningUrl()}{{\bf  getVersioningUrl}\\}
\texttt{public java.lang.String\ {\bf  getVersioningUrl}()
\label{it.unisa.sesa.repominer.db.entities.Project.getVersioningUrl()}}%end signature
}%end item
\item{ 
\index{hashCode()}
\hypertarget{it.unisa.sesa.repominer.db.entities.Project.hashCode()}{{\bf  hashCode}\\}
\texttt{public native int\ {\bf  hashCode}()
\label{it.unisa.sesa.repominer.db.entities.Project.hashCode()}}%end signature
}%end item
\item{ 
\index{setBugtrackerUrl(String)}
\hypertarget{it.unisa.sesa.repominer.db.entities.Project.setBugtrackerUrl(java.lang.String)}{{\bf  setBugtrackerUrl}\\}
\texttt{public void\ {\bf  setBugtrackerUrl}(\texttt{java.lang.String} {\bf  bugtrackerUrl})
\label{it.unisa.sesa.repominer.db.entities.Project.setBugtrackerUrl(java.lang.String)}}%end signature
}%end item
\item{ 
\index{setId(Integer)}
\hypertarget{it.unisa.sesa.repominer.db.entities.Project.setId(java.lang.Integer)}{{\bf  setId}\\}
\texttt{public void\ {\bf  setId}(\texttt{java.lang.Integer} {\bf  id})
\label{it.unisa.sesa.repominer.db.entities.Project.setId(java.lang.Integer)}}%end signature
}%end item
\item{ 
\index{setName(String)}
\hypertarget{it.unisa.sesa.repominer.db.entities.Project.setName(java.lang.String)}{{\bf  setName}\\}
\texttt{public void\ {\bf  setName}(\texttt{java.lang.String} {\bf  name})
\label{it.unisa.sesa.repominer.db.entities.Project.setName(java.lang.String)}}%end signature
}%end item
\item{ 
\index{setVersioningUrl(String)}
\hypertarget{it.unisa.sesa.repominer.db.entities.Project.setVersioningUrl(java.lang.String)}{{\bf  setVersioningUrl}\\}
\texttt{public void\ {\bf  setVersioningUrl}(\texttt{java.lang.String} {\bf  versioningUrl})
\label{it.unisa.sesa.repominer.db.entities.Project.setVersioningUrl(java.lang.String)}}%end signature
}%end item
\item{ 
\index{toString()}
\hypertarget{it.unisa.sesa.repominer.db.entities.Project.toString()}{{\bf  toString}\\}
\texttt{public java.lang.String\ {\bf  toString}()
\label{it.unisa.sesa.repominer.db.entities.Project.toString()}}%end signature
}%end item
\end{itemize}
}
}
\subsection{\label{it.unisa.sesa.repominer.db.entities.ProjectMetric}Class ProjectMetric}{
\hypertarget{it.unisa.sesa.repominer.db.entities.ProjectMetric}{}\vskip .1in 
\subsubsection{Declaration}{
\small public class ProjectMetric
\\ {\bf  extends} it.unisa.sesa.repominer.db.entities.Metric
\refdefined{it.unisa.sesa.repominer.db.entities.Metric}}
\subsubsection{Constructor summary}{
\begin{verse}
\hyperlink{it.unisa.sesa.repominer.db.entities.ProjectMetric()}{{\bf ProjectMetric()}} \\
\hyperlink{it.unisa.sesa.repominer.db.entities.ProjectMetric(java.lang.Integer, java.lang.String, java.lang.String, java.lang.Integer, java.lang.Integer, java.lang.Double, java.util.Date, java.util.Date)}{{\bf ProjectMetric(Integer, String, String, Integer, Integer, Double, Date, Date)}} \\
\end{verse}
}
\subsubsection{Method summary}{
\begin{verse}
\hyperlink{it.unisa.sesa.repominer.db.entities.ProjectMetric.equals(java.lang.Object)}{{\bf equals(Object)}} \\
\hyperlink{it.unisa.sesa.repominer.db.entities.ProjectMetric.getEnd()}{{\bf getEnd()}} \\
\hyperlink{it.unisa.sesa.repominer.db.entities.ProjectMetric.getMetricId()}{{\bf getMetricId()}} \\
\hyperlink{it.unisa.sesa.repominer.db.entities.ProjectMetric.getProjectId()}{{\bf getProjectId()}} \\
\hyperlink{it.unisa.sesa.repominer.db.entities.ProjectMetric.getStart()}{{\bf getStart()}} \\
\hyperlink{it.unisa.sesa.repominer.db.entities.ProjectMetric.getValue()}{{\bf getValue()}} \\
\hyperlink{it.unisa.sesa.repominer.db.entities.ProjectMetric.hashCode()}{{\bf hashCode()}} \\
\hyperlink{it.unisa.sesa.repominer.db.entities.ProjectMetric.setEnd(java.util.Date)}{{\bf setEnd(Date)}} \\
\hyperlink{it.unisa.sesa.repominer.db.entities.ProjectMetric.setMetricId(java.lang.Integer)}{{\bf setMetricId(Integer)}} \\
\hyperlink{it.unisa.sesa.repominer.db.entities.ProjectMetric.setProjectId(java.lang.Integer)}{{\bf setProjectId(Integer)}} \\
\hyperlink{it.unisa.sesa.repominer.db.entities.ProjectMetric.setStart(java.util.Date)}{{\bf setStart(Date)}} \\
\hyperlink{it.unisa.sesa.repominer.db.entities.ProjectMetric.setValue(java.lang.Double)}{{\bf setValue(Double)}} \\
\hyperlink{it.unisa.sesa.repominer.db.entities.ProjectMetric.toString()}{{\bf toString()}} \\
\end{verse}
}
\subsubsection{Constructors}{
\vskip -2em
\begin{itemize}
\item{ 
\index{ProjectMetric()}
\hypertarget{it.unisa.sesa.repominer.db.entities.ProjectMetric()}{{\bf  ProjectMetric}\\}
\texttt{public\ {\bf  ProjectMetric}()
\label{it.unisa.sesa.repominer.db.entities.ProjectMetric()}}%end signature
}%end item
\item{ 
\index{ProjectMetric(Integer, String, String, Integer, Integer, Double, Date, Date)}
\hypertarget{it.unisa.sesa.repominer.db.entities.ProjectMetric(java.lang.Integer, java.lang.String, java.lang.String, java.lang.Integer, java.lang.Integer, java.lang.Double, java.util.Date, java.util.Date)}{{\bf  ProjectMetric}\\}
\texttt{public\ {\bf  ProjectMetric}(\texttt{java.lang.Integer} {\bf  id},
\texttt{java.lang.String} {\bf  name},
\texttt{java.lang.String} {\bf  description},
\texttt{java.lang.Integer} {\bf  projectId},
\texttt{java.lang.Integer} {\bf  metricId},
\texttt{java.lang.Double} {\bf  value},
\texttt{java.util.Date} {\bf  start},
\texttt{java.util.Date} {\bf  end})
\label{it.unisa.sesa.repominer.db.entities.ProjectMetric(java.lang.Integer, java.lang.String, java.lang.String, java.lang.Integer, java.lang.Integer, java.lang.Double, java.util.Date, java.util.Date)}}%end signature
}%end item
\end{itemize}
}
\subsubsection{Methods}{
\vskip -2em
\begin{itemize}
\item{ 
\index{equals(Object)}
\hypertarget{it.unisa.sesa.repominer.db.entities.ProjectMetric.equals(java.lang.Object)}{{\bf  equals}\\}
\texttt{public boolean\ {\bf  equals}(\texttt{java.lang.Object} {\bf  arg0})
\label{it.unisa.sesa.repominer.db.entities.ProjectMetric.equals(java.lang.Object)}}%end signature
}%end item
\item{ 
\index{getEnd()}
\hypertarget{it.unisa.sesa.repominer.db.entities.ProjectMetric.getEnd()}{{\bf  getEnd}\\}
\texttt{public java.util.Date\ {\bf  getEnd}()
\label{it.unisa.sesa.repominer.db.entities.ProjectMetric.getEnd()}}%end signature
}%end item
\item{ 
\index{getMetricId()}
\hypertarget{it.unisa.sesa.repominer.db.entities.ProjectMetric.getMetricId()}{{\bf  getMetricId}\\}
\texttt{public java.lang.Integer\ {\bf  getMetricId}()
\label{it.unisa.sesa.repominer.db.entities.ProjectMetric.getMetricId()}}%end signature
}%end item
\item{ 
\index{getProjectId()}
\hypertarget{it.unisa.sesa.repominer.db.entities.ProjectMetric.getProjectId()}{{\bf  getProjectId}\\}
\texttt{public java.lang.Integer\ {\bf  getProjectId}()
\label{it.unisa.sesa.repominer.db.entities.ProjectMetric.getProjectId()}}%end signature
}%end item
\item{ 
\index{getStart()}
\hypertarget{it.unisa.sesa.repominer.db.entities.ProjectMetric.getStart()}{{\bf  getStart}\\}
\texttt{public java.util.Date\ {\bf  getStart}()
\label{it.unisa.sesa.repominer.db.entities.ProjectMetric.getStart()}}%end signature
}%end item
\item{ 
\index{getValue()}
\hypertarget{it.unisa.sesa.repominer.db.entities.ProjectMetric.getValue()}{{\bf  getValue}\\}
\texttt{public java.lang.Double\ {\bf  getValue}()
\label{it.unisa.sesa.repominer.db.entities.ProjectMetric.getValue()}}%end signature
}%end item
\item{ 
\index{hashCode()}
\hypertarget{it.unisa.sesa.repominer.db.entities.ProjectMetric.hashCode()}{{\bf  hashCode}\\}
\texttt{public native int\ {\bf  hashCode}()
\label{it.unisa.sesa.repominer.db.entities.ProjectMetric.hashCode()}}%end signature
}%end item
\item{ 
\index{setEnd(Date)}
\hypertarget{it.unisa.sesa.repominer.db.entities.ProjectMetric.setEnd(java.util.Date)}{{\bf  setEnd}\\}
\texttt{public void\ {\bf  setEnd}(\texttt{java.util.Date} {\bf  end})
\label{it.unisa.sesa.repominer.db.entities.ProjectMetric.setEnd(java.util.Date)}}%end signature
}%end item
\item{ 
\index{setMetricId(Integer)}
\hypertarget{it.unisa.sesa.repominer.db.entities.ProjectMetric.setMetricId(java.lang.Integer)}{{\bf  setMetricId}\\}
\texttt{public void\ {\bf  setMetricId}(\texttt{java.lang.Integer} {\bf  metricId})
\label{it.unisa.sesa.repominer.db.entities.ProjectMetric.setMetricId(java.lang.Integer)}}%end signature
}%end item
\item{ 
\index{setProjectId(Integer)}
\hypertarget{it.unisa.sesa.repominer.db.entities.ProjectMetric.setProjectId(java.lang.Integer)}{{\bf  setProjectId}\\}
\texttt{public void\ {\bf  setProjectId}(\texttt{java.lang.Integer} {\bf  projectId})
\label{it.unisa.sesa.repominer.db.entities.ProjectMetric.setProjectId(java.lang.Integer)}}%end signature
}%end item
\item{ 
\index{setStart(Date)}
\hypertarget{it.unisa.sesa.repominer.db.entities.ProjectMetric.setStart(java.util.Date)}{{\bf  setStart}\\}
\texttt{public void\ {\bf  setStart}(\texttt{java.util.Date} {\bf  start})
\label{it.unisa.sesa.repominer.db.entities.ProjectMetric.setStart(java.util.Date)}}%end signature
}%end item
\item{ 
\index{setValue(Double)}
\hypertarget{it.unisa.sesa.repominer.db.entities.ProjectMetric.setValue(java.lang.Double)}{{\bf  setValue}\\}
\texttt{public void\ {\bf  setValue}(\texttt{java.lang.Double} {\bf  value})
\label{it.unisa.sesa.repominer.db.entities.ProjectMetric.setValue(java.lang.Double)}}%end signature
}%end item
\item{ 
\index{toString()}
\hypertarget{it.unisa.sesa.repominer.db.entities.ProjectMetric.toString()}{{\bf  toString}\\}
\texttt{public java.lang.String\ {\bf  toString}()
\label{it.unisa.sesa.repominer.db.entities.ProjectMetric.toString()}}%end signature
}%end item
\end{itemize}
}
\subsubsection{Members inherited from class Metric }{
\texttt{it.unisa.sesa.repominer.db.entities.Metric} {\small 
\refdefined{it.unisa.sesa.repominer.db.entities.Metric}}
{\small 

\vskip -2em
\begin{itemize}
\item{\vskip -1.5ex 
\texttt{public static final {\bf  BCCM\_DESCRIPTION}}%end signature
}%end item
\item{\vskip -1.5ex 
\texttt{public static final {\bf  BCCM\_NAME}}%end signature
}%end item
\item{\vskip -1.5ex 
\texttt{public static final {\bf  CHANGE\_SET\_SIZE\_DESCRIPTION}}%end signature
}%end item
\item{\vskip -1.5ex 
\texttt{public static final {\bf  CHANGE\_SET\_SIZE\_NAME}}%end signature
}%end item
\item{\vskip -1.5ex 
\texttt{public static final {\bf  ECCM\_DESCRIPTION}}%end signature
}%end item
\item{\vskip -1.5ex 
\texttt{public static final {\bf  ECCM\_NAME}}%end signature
}%end item
\item{\vskip -1.5ex 
\texttt{public static final {\bf  ECCM\_STATIC\_DESCRIPTION}}%end signature
}%end item
\item{\vskip -1.5ex 
\texttt{public static final {\bf  ECCM\_STATIC\_NAME}}%end signature
}%end item
\item{\vskip -1.5ex 
\texttt{public boolean {\bf  equals}(\texttt{java.lang.Object} {\bf  obj})
}%end signature
}%end item
\item{\vskip -1.5ex 
\texttt{public String {\bf  getDescription}()
}%end signature
}%end item
\item{\vskip -1.5ex 
\texttt{public Integer {\bf  getId}()
}%end signature
}%end item
\item{\vskip -1.5ex 
\texttt{public String {\bf  getName}()
}%end signature
}%end item
\item{\vskip -1.5ex 
\texttt{public int {\bf  hashCode}()
}%end signature
}%end item
\item{\vskip -1.5ex 
\texttt{public static final {\bf  MAX\_LINES\_DESCRIPTION}}%end signature
}%end item
\item{\vskip -1.5ex 
\texttt{public static final {\bf  MAX\_LINES\_NAME}}%end signature
}%end item
\item{\vskip -1.5ex 
\texttt{public static final {\bf  MEAN\_LINES\_DESCRIPTION}}%end signature
}%end item
\item{\vskip -1.5ex 
\texttt{public static final {\bf  MEAN\_LINES\_NAME}}%end signature
}%end item
\item{\vskip -1.5ex 
\texttt{public static final {\bf  NAUTH\_DESCRIPTION}}%end signature
}%end item
\item{\vskip -1.5ex 
\texttt{public static final {\bf  NAUTH\_NAME}}%end signature
}%end item
\item{\vskip -1.5ex 
\texttt{public static final {\bf  NCHANGE\_DESCRIPTION}}%end signature
}%end item
\item{\vskip -1.5ex 
\texttt{public static final {\bf  NCHANGE\_NAME}}%end signature
}%end item
\item{\vskip -1.5ex 
\texttt{public static final {\bf  NFIX\_DESCRIPTION}}%end signature
}%end item
\item{\vskip -1.5ex 
\texttt{public static final {\bf  NFIX\_NAME}}%end signature
}%end item
\item{\vskip -1.5ex 
\texttt{public static final {\bf  NREF\_DESCRIPTION}}%end signature
}%end item
\item{\vskip -1.5ex 
\texttt{public static final {\bf  NREF\_NAME}}%end signature
}%end item
\item{\vskip -1.5ex 
\texttt{public static final {\bf  NUM\_REVISION\_DESCRIPTION}}%end signature
}%end item
\item{\vskip -1.5ex 
\texttt{public static final {\bf  NUM\_REVISION\_NAME}}%end signature
}%end item
\item{\vskip -1.5ex 
\texttt{public static final {\bf  REVISION\_DESCRIPTION}}%end signature
}%end item
\item{\vskip -1.5ex 
\texttt{public static final {\bf  REVISION\_NAME}}%end signature
}%end item
\item{\vskip -1.5ex 
\texttt{public void {\bf  setDescription}(\texttt{java.lang.String} {\bf  description})
}%end signature
}%end item
\item{\vskip -1.5ex 
\texttt{public void {\bf  setId}(\texttt{java.lang.Integer} {\bf  id})
}%end signature
}%end item
\item{\vskip -1.5ex 
\texttt{public void {\bf  setName}(\texttt{java.lang.String} {\bf  name})
}%end signature
}%end item
\item{\vskip -1.5ex 
\texttt{public static final {\bf  SUM\_LINES\_DESCRIPTION}}%end signature
}%end item
\item{\vskip -1.5ex 
\texttt{public static final {\bf  SUM\_LINES\_NAME}}%end signature
}%end item
\item{\vskip -1.5ex 
\texttt{public String {\bf  toString}()
}%end signature
}%end item
\end{itemize}
}
}
\subsection{\label{it.unisa.sesa.repominer.db.entities.Review}Class Review}{
\hypertarget{it.unisa.sesa.repominer.db.entities.Review}{}\vskip .1in 
\subsubsection{Declaration}{
\small public class Review
\\ {\bf  extends} java.lang.Object
\refdefined{java.lang.Object}}
\subsubsection{Constructor summary}{
\begin{verse}
\hyperlink{it.unisa.sesa.repominer.db.entities.Review()}{{\bf Review()}} \\
\hyperlink{it.unisa.sesa.repominer.db.entities.Review(java.lang.String, java.lang.String, java.lang.String, java.lang.String, java.lang.String, java.lang.String, java.util.Date)}{{\bf Review(String, String, String, String, String, String, Date)}} \\
\end{verse}
}
\subsubsection{Method summary}{
\begin{verse}
\hyperlink{it.unisa.sesa.repominer.db.entities.Review.equals(java.lang.Object)}{{\bf equals(Object)}} \\
\hyperlink{it.unisa.sesa.repominer.db.entities.Review.getAuthor()}{{\bf getAuthor()}} \\
\hyperlink{it.unisa.sesa.repominer.db.entities.Review.getDate()}{{\bf getDate()}} \\
\hyperlink{it.unisa.sesa.repominer.db.entities.Review.getNameApp()}{{\bf getNameApp()}} \\
\hyperlink{it.unisa.sesa.repominer.db.entities.Review.getRating()}{{\bf getRating()}} \\
\hyperlink{it.unisa.sesa.repominer.db.entities.Review.getReview()}{{\bf getReview()}} \\
\hyperlink{it.unisa.sesa.repominer.db.entities.Review.getTitle()}{{\bf getTitle()}} \\
\hyperlink{it.unisa.sesa.repominer.db.entities.Review.getVersioningUrl()}{{\bf getVersioningUrl()}} \\
\hyperlink{it.unisa.sesa.repominer.db.entities.Review.hashCode()}{{\bf hashCode()}} \\
\hyperlink{it.unisa.sesa.repominer.db.entities.Review.setAuthor(java.lang.String)}{{\bf setAuthor(String)}} \\
\hyperlink{it.unisa.sesa.repominer.db.entities.Review.setDate(java.util.Date)}{{\bf setDate(Date)}} \\
\hyperlink{it.unisa.sesa.repominer.db.entities.Review.setNameApp(java.lang.String)}{{\bf setNameApp(String)}} \\
\hyperlink{it.unisa.sesa.repominer.db.entities.Review.setRating(java.lang.String)}{{\bf setRating(String)}} \\
\hyperlink{it.unisa.sesa.repominer.db.entities.Review.setReview(java.lang.String)}{{\bf setReview(String)}} \\
\hyperlink{it.unisa.sesa.repominer.db.entities.Review.setTitle(java.lang.String)}{{\bf setTitle(String)}} \\
\hyperlink{it.unisa.sesa.repominer.db.entities.Review.setVersioningUrl(java.lang.String)}{{\bf setVersioningUrl(String)}} \\
\hyperlink{it.unisa.sesa.repominer.db.entities.Review.toString()}{{\bf toString()}} \\
\end{verse}
}
\subsubsection{Constructors}{
\vskip -2em
\begin{itemize}
\item{ 
\index{Review()}
\hypertarget{it.unisa.sesa.repominer.db.entities.Review()}{{\bf  Review}\\}
\texttt{public\ {\bf  Review}()
\label{it.unisa.sesa.repominer.db.entities.Review()}}%end signature
}%end item
\item{ 
\index{Review(String, String, String, String, String, String, Date)}
\hypertarget{it.unisa.sesa.repominer.db.entities.Review(java.lang.String, java.lang.String, java.lang.String, java.lang.String, java.lang.String, java.lang.String, java.util.Date)}{{\bf  Review}\\}
\texttt{public\ {\bf  Review}(\texttt{java.lang.String} {\bf  versioningUrl},
\texttt{java.lang.String} {\bf  nameApp},
\texttt{java.lang.String} {\bf  author},
\texttt{java.lang.String} {\bf  title},
\texttt{java.lang.String} {\bf  review},
\texttt{java.lang.String} {\bf  rating},
\texttt{java.util.Date} {\bf  date})
\label{it.unisa.sesa.repominer.db.entities.Review(java.lang.String, java.lang.String, java.lang.String, java.lang.String, java.lang.String, java.lang.String, java.util.Date)}}%end signature
}%end item
\end{itemize}
}
\subsubsection{Methods}{
\vskip -2em
\begin{itemize}
\item{ 
\index{equals(Object)}
\hypertarget{it.unisa.sesa.repominer.db.entities.Review.equals(java.lang.Object)}{{\bf  equals}\\}
\texttt{public boolean\ {\bf  equals}(\texttt{java.lang.Object} {\bf  arg0})
\label{it.unisa.sesa.repominer.db.entities.Review.equals(java.lang.Object)}}%end signature
}%end item
\item{ 
\index{getAuthor()}
\hypertarget{it.unisa.sesa.repominer.db.entities.Review.getAuthor()}{{\bf  getAuthor}\\}
\texttt{public java.lang.String\ {\bf  getAuthor}()
\label{it.unisa.sesa.repominer.db.entities.Review.getAuthor()}}%end signature
}%end item
\item{ 
\index{getDate()}
\hypertarget{it.unisa.sesa.repominer.db.entities.Review.getDate()}{{\bf  getDate}\\}
\texttt{public java.util.Date\ {\bf  getDate}()
\label{it.unisa.sesa.repominer.db.entities.Review.getDate()}}%end signature
}%end item
\item{ 
\index{getNameApp()}
\hypertarget{it.unisa.sesa.repominer.db.entities.Review.getNameApp()}{{\bf  getNameApp}\\}
\texttt{public java.lang.String\ {\bf  getNameApp}()
\label{it.unisa.sesa.repominer.db.entities.Review.getNameApp()}}%end signature
}%end item
\item{ 
\index{getRating()}
\hypertarget{it.unisa.sesa.repominer.db.entities.Review.getRating()}{{\bf  getRating}\\}
\texttt{public java.lang.String\ {\bf  getRating}()
\label{it.unisa.sesa.repominer.db.entities.Review.getRating()}}%end signature
}%end item
\item{ 
\index{getReview()}
\hypertarget{it.unisa.sesa.repominer.db.entities.Review.getReview()}{{\bf  getReview}\\}
\texttt{public java.lang.String\ {\bf  getReview}()
\label{it.unisa.sesa.repominer.db.entities.Review.getReview()}}%end signature
}%end item
\item{ 
\index{getTitle()}
\hypertarget{it.unisa.sesa.repominer.db.entities.Review.getTitle()}{{\bf  getTitle}\\}
\texttt{public java.lang.String\ {\bf  getTitle}()
\label{it.unisa.sesa.repominer.db.entities.Review.getTitle()}}%end signature
}%end item
\item{ 
\index{getVersioningUrl()}
\hypertarget{it.unisa.sesa.repominer.db.entities.Review.getVersioningUrl()}{{\bf  getVersioningUrl}\\}
\texttt{public java.lang.String\ {\bf  getVersioningUrl}()
\label{it.unisa.sesa.repominer.db.entities.Review.getVersioningUrl()}}%end signature
}%end item
\item{ 
\index{hashCode()}
\hypertarget{it.unisa.sesa.repominer.db.entities.Review.hashCode()}{{\bf  hashCode}\\}
\texttt{public native int\ {\bf  hashCode}()
\label{it.unisa.sesa.repominer.db.entities.Review.hashCode()}}%end signature
}%end item
\item{ 
\index{setAuthor(String)}
\hypertarget{it.unisa.sesa.repominer.db.entities.Review.setAuthor(java.lang.String)}{{\bf  setAuthor}\\}
\texttt{public void\ {\bf  setAuthor}(\texttt{java.lang.String} {\bf  author})
\label{it.unisa.sesa.repominer.db.entities.Review.setAuthor(java.lang.String)}}%end signature
}%end item
\item{ 
\index{setDate(Date)}
\hypertarget{it.unisa.sesa.repominer.db.entities.Review.setDate(java.util.Date)}{{\bf  setDate}\\}
\texttt{public void\ {\bf  setDate}(\texttt{java.util.Date} {\bf  date})
\label{it.unisa.sesa.repominer.db.entities.Review.setDate(java.util.Date)}}%end signature
}%end item
\item{ 
\index{setNameApp(String)}
\hypertarget{it.unisa.sesa.repominer.db.entities.Review.setNameApp(java.lang.String)}{{\bf  setNameApp}\\}
\texttt{public void\ {\bf  setNameApp}(\texttt{java.lang.String} {\bf  nameApp})
\label{it.unisa.sesa.repominer.db.entities.Review.setNameApp(java.lang.String)}}%end signature
}%end item
\item{ 
\index{setRating(String)}
\hypertarget{it.unisa.sesa.repominer.db.entities.Review.setRating(java.lang.String)}{{\bf  setRating}\\}
\texttt{public void\ {\bf  setRating}(\texttt{java.lang.String} {\bf  rating})
\label{it.unisa.sesa.repominer.db.entities.Review.setRating(java.lang.String)}}%end signature
}%end item
\item{ 
\index{setReview(String)}
\hypertarget{it.unisa.sesa.repominer.db.entities.Review.setReview(java.lang.String)}{{\bf  setReview}\\}
\texttt{public void\ {\bf  setReview}(\texttt{java.lang.String} {\bf  review})
\label{it.unisa.sesa.repominer.db.entities.Review.setReview(java.lang.String)}}%end signature
}%end item
\item{ 
\index{setTitle(String)}
\hypertarget{it.unisa.sesa.repominer.db.entities.Review.setTitle(java.lang.String)}{{\bf  setTitle}\\}
\texttt{public void\ {\bf  setTitle}(\texttt{java.lang.String} {\bf  title})
\label{it.unisa.sesa.repominer.db.entities.Review.setTitle(java.lang.String)}}%end signature
}%end item
\item{ 
\index{setVersioningUrl(String)}
\hypertarget{it.unisa.sesa.repominer.db.entities.Review.setVersioningUrl(java.lang.String)}{{\bf  setVersioningUrl}\\}
\texttt{public void\ {\bf  setVersioningUrl}(\texttt{java.lang.String} {\bf  versioningUrl})
\label{it.unisa.sesa.repominer.db.entities.Review.setVersioningUrl(java.lang.String)}}%end signature
}%end item
\item{ 
\index{toString()}
\hypertarget{it.unisa.sesa.repominer.db.entities.Review.toString()}{{\bf  toString}\\}
\texttt{public java.lang.String\ {\bf  toString}()
\label{it.unisa.sesa.repominer.db.entities.Review.toString()}}%end signature
}%end item
\end{itemize}
}
}
\subsection{\label{it.unisa.sesa.repominer.db.entities.SourceContainer}Class SourceContainer}{
\hypertarget{it.unisa.sesa.repominer.db.entities.SourceContainer}{}\vskip .1in 
\subsubsection{Declaration}{
\small public class SourceContainer
\\ {\bf  extends} java.lang.Object
\refdefined{java.lang.Object}}
\subsubsection{Constructor summary}{
\begin{verse}
\hyperlink{it.unisa.sesa.repominer.db.entities.SourceContainer()}{{\bf SourceContainer()}} \\
\hyperlink{it.unisa.sesa.repominer.db.entities.SourceContainer(java.lang.Integer, java.lang.Integer, java.lang.Integer)}{{\bf SourceContainer(Integer, Integer, Integer)}} \\
\end{verse}
}
\subsubsection{Method summary}{
\begin{verse}
\hyperlink{it.unisa.sesa.repominer.db.entities.SourceContainer.equals(java.lang.Object)}{{\bf equals(Object)}} \\
\hyperlink{it.unisa.sesa.repominer.db.entities.SourceContainer.getId()}{{\bf getId()}} \\
\hyperlink{it.unisa.sesa.repominer.db.entities.SourceContainer.getImportId()}{{\bf getImportId()}} \\
\hyperlink{it.unisa.sesa.repominer.db.entities.SourceContainer.getName()}{{\bf getName()}} \\
\hyperlink{it.unisa.sesa.repominer.db.entities.SourceContainer.getProjectId()}{{\bf getProjectId()}} \\
\hyperlink{it.unisa.sesa.repominer.db.entities.SourceContainer.hashCode()}{{\bf hashCode()}} \\
\hyperlink{it.unisa.sesa.repominer.db.entities.SourceContainer.setId(java.lang.Integer)}{{\bf setId(Integer)}} \\
\hyperlink{it.unisa.sesa.repominer.db.entities.SourceContainer.setImportId(java.lang.Integer)}{{\bf setImportId(Integer)}} \\
\hyperlink{it.unisa.sesa.repominer.db.entities.SourceContainer.setName(java.lang.String)}{{\bf setName(String)}} \\
\hyperlink{it.unisa.sesa.repominer.db.entities.SourceContainer.setProjectId(java.lang.Integer)}{{\bf setProjectId(Integer)}} \\
\hyperlink{it.unisa.sesa.repominer.db.entities.SourceContainer.toString()}{{\bf toString()}} \\
\end{verse}
}
\subsubsection{Constructors}{
\vskip -2em
\begin{itemize}
\item{ 
\index{SourceContainer()}
\hypertarget{it.unisa.sesa.repominer.db.entities.SourceContainer()}{{\bf  SourceContainer}\\}
\texttt{public\ {\bf  SourceContainer}()
\label{it.unisa.sesa.repominer.db.entities.SourceContainer()}}%end signature
}%end item
\item{ 
\index{SourceContainer(Integer, Integer, Integer)}
\hypertarget{it.unisa.sesa.repominer.db.entities.SourceContainer(java.lang.Integer, java.lang.Integer, java.lang.Integer)}{{\bf  SourceContainer}\\}
\texttt{public\ {\bf  SourceContainer}(\texttt{java.lang.Integer} {\bf  id},
\texttt{java.lang.Integer} {\bf  projectId},
\texttt{java.lang.Integer} {\bf  importId})
\label{it.unisa.sesa.repominer.db.entities.SourceContainer(java.lang.Integer, java.lang.Integer, java.lang.Integer)}}%end signature
}%end item
\end{itemize}
}
\subsubsection{Methods}{
\vskip -2em
\begin{itemize}
\item{ 
\index{equals(Object)}
\hypertarget{it.unisa.sesa.repominer.db.entities.SourceContainer.equals(java.lang.Object)}{{\bf  equals}\\}
\texttt{public boolean\ {\bf  equals}(\texttt{java.lang.Object} {\bf  arg0})
\label{it.unisa.sesa.repominer.db.entities.SourceContainer.equals(java.lang.Object)}}%end signature
}%end item
\item{ 
\index{getId()}
\hypertarget{it.unisa.sesa.repominer.db.entities.SourceContainer.getId()}{{\bf  getId}\\}
\texttt{public java.lang.Integer\ {\bf  getId}()
\label{it.unisa.sesa.repominer.db.entities.SourceContainer.getId()}}%end signature
}%end item
\item{ 
\index{getImportId()}
\hypertarget{it.unisa.sesa.repominer.db.entities.SourceContainer.getImportId()}{{\bf  getImportId}\\}
\texttt{public java.lang.Integer\ {\bf  getImportId}()
\label{it.unisa.sesa.repominer.db.entities.SourceContainer.getImportId()}}%end signature
}%end item
\item{ 
\index{getName()}
\hypertarget{it.unisa.sesa.repominer.db.entities.SourceContainer.getName()}{{\bf  getName}\\}
\texttt{public java.lang.String\ {\bf  getName}()
\label{it.unisa.sesa.repominer.db.entities.SourceContainer.getName()}}%end signature
}%end item
\item{ 
\index{getProjectId()}
\hypertarget{it.unisa.sesa.repominer.db.entities.SourceContainer.getProjectId()}{{\bf  getProjectId}\\}
\texttt{public java.lang.Integer\ {\bf  getProjectId}()
\label{it.unisa.sesa.repominer.db.entities.SourceContainer.getProjectId()}}%end signature
}%end item
\item{ 
\index{hashCode()}
\hypertarget{it.unisa.sesa.repominer.db.entities.SourceContainer.hashCode()}{{\bf  hashCode}\\}
\texttt{public native int\ {\bf  hashCode}()
\label{it.unisa.sesa.repominer.db.entities.SourceContainer.hashCode()}}%end signature
}%end item
\item{ 
\index{setId(Integer)}
\hypertarget{it.unisa.sesa.repominer.db.entities.SourceContainer.setId(java.lang.Integer)}{{\bf  setId}\\}
\texttt{public void\ {\bf  setId}(\texttt{java.lang.Integer} {\bf  id})
\label{it.unisa.sesa.repominer.db.entities.SourceContainer.setId(java.lang.Integer)}}%end signature
}%end item
\item{ 
\index{setImportId(Integer)}
\hypertarget{it.unisa.sesa.repominer.db.entities.SourceContainer.setImportId(java.lang.Integer)}{{\bf  setImportId}\\}
\texttt{public void\ {\bf  setImportId}(\texttt{java.lang.Integer} {\bf  importId})
\label{it.unisa.sesa.repominer.db.entities.SourceContainer.setImportId(java.lang.Integer)}}%end signature
}%end item
\item{ 
\index{setName(String)}
\hypertarget{it.unisa.sesa.repominer.db.entities.SourceContainer.setName(java.lang.String)}{{\bf  setName}\\}
\texttt{public void\ {\bf  setName}(\texttt{java.lang.String} {\bf  name})
\label{it.unisa.sesa.repominer.db.entities.SourceContainer.setName(java.lang.String)}}%end signature
}%end item
\item{ 
\index{setProjectId(Integer)}
\hypertarget{it.unisa.sesa.repominer.db.entities.SourceContainer.setProjectId(java.lang.Integer)}{{\bf  setProjectId}\\}
\texttt{public void\ {\bf  setProjectId}(\texttt{java.lang.Integer} {\bf  projectId})
\label{it.unisa.sesa.repominer.db.entities.SourceContainer.setProjectId(java.lang.Integer)}}%end signature
}%end item
\item{ 
\index{toString()}
\hypertarget{it.unisa.sesa.repominer.db.entities.SourceContainer.toString()}{{\bf  toString}\\}
\texttt{public java.lang.String\ {\bf  toString}()
\label{it.unisa.sesa.repominer.db.entities.SourceContainer.toString()}}%end signature
}%end item
\end{itemize}
}
}
\subsection{\label{it.unisa.sesa.repominer.db.entities.Type}Class Type}{
\hypertarget{it.unisa.sesa.repominer.db.entities.Type}{}\vskip .1in 
\subsubsection{Declaration}{
\small public class Type
\\ {\bf  extends} java.lang.Object
\refdefined{java.lang.Object}}
\subsubsection{Constructor summary}{
\begin{verse}
\hyperlink{it.unisa.sesa.repominer.db.entities.Type()}{{\bf Type()}} \\
\hyperlink{it.unisa.sesa.repominer.db.entities.Type(java.lang.Integer, java.lang.Integer, java.lang.Integer, java.lang.String, java.lang.String, java.lang.Integer)}{{\bf Type(Integer, Integer, Integer, String, String, Integer)}} \\
\end{verse}
}
\subsubsection{Method summary}{
\begin{verse}
\hyperlink{it.unisa.sesa.repominer.db.entities.Type.equals(java.lang.Object)}{{\bf equals(Object)}} \\
\hyperlink{it.unisa.sesa.repominer.db.entities.Type.getHeaderFileLocation()}{{\bf getHeaderFileLocation()}} \\
\hyperlink{it.unisa.sesa.repominer.db.entities.Type.getId()}{{\bf getId()}} \\
\hyperlink{it.unisa.sesa.repominer.db.entities.Type.getImportId()}{{\bf getImportId()}} \\
\hyperlink{it.unisa.sesa.repominer.db.entities.Type.getLinesNumber()}{{\bf getLinesNumber()}} \\
\hyperlink{it.unisa.sesa.repominer.db.entities.Type.getSourceContainer()}{{\bf getSourceContainer()}} \\
\hyperlink{it.unisa.sesa.repominer.db.entities.Type.getSrcFileLocation()}{{\bf getSrcFileLocation()}} \\
\hyperlink{it.unisa.sesa.repominer.db.entities.Type.hashCode()}{{\bf hashCode()}} \\
\hyperlink{it.unisa.sesa.repominer.db.entities.Type.setHeaderFileLocation(java.lang.String)}{{\bf setHeaderFileLocation(String)}} \\
\hyperlink{it.unisa.sesa.repominer.db.entities.Type.setId(java.lang.Integer)}{{\bf setId(Integer)}} \\
\hyperlink{it.unisa.sesa.repominer.db.entities.Type.setImportId(java.lang.Integer)}{{\bf setImportId(Integer)}} \\
\hyperlink{it.unisa.sesa.repominer.db.entities.Type.setLinesNumber(java.lang.Integer)}{{\bf setLinesNumber(Integer)}} \\
\hyperlink{it.unisa.sesa.repominer.db.entities.Type.setSourceContainer(java.lang.Integer)}{{\bf setSourceContainer(Integer)}} \\
\hyperlink{it.unisa.sesa.repominer.db.entities.Type.setSrcFileLocation(java.lang.String)}{{\bf setSrcFileLocation(String)}} \\
\hyperlink{it.unisa.sesa.repominer.db.entities.Type.toString()}{{\bf toString()}} \\
\end{verse}
}
\subsubsection{Constructors}{
\vskip -2em
\begin{itemize}
\item{ 
\index{Type()}
\hypertarget{it.unisa.sesa.repominer.db.entities.Type()}{{\bf  Type}\\}
\texttt{public\ {\bf  Type}()
\label{it.unisa.sesa.repominer.db.entities.Type()}}%end signature
}%end item
\item{ 
\index{Type(Integer, Integer, Integer, String, String, Integer)}
\hypertarget{it.unisa.sesa.repominer.db.entities.Type(java.lang.Integer, java.lang.Integer, java.lang.Integer, java.lang.String, java.lang.String, java.lang.Integer)}{{\bf  Type}\\}
\texttt{public\ {\bf  Type}(\texttt{java.lang.Integer} {\bf  id},
\texttt{java.lang.Integer} {\bf  importId},
\texttt{java.lang.Integer} {\bf  linesNumber},
\texttt{java.lang.String} {\bf  srcFileLocation},
\texttt{java.lang.String} {\bf  headerFileLocation},
\texttt{java.lang.Integer} {\bf  sourceContainer})
\label{it.unisa.sesa.repominer.db.entities.Type(java.lang.Integer, java.lang.Integer, java.lang.Integer, java.lang.String, java.lang.String, java.lang.Integer)}}%end signature
}%end item
\end{itemize}
}
\subsubsection{Methods}{
\vskip -2em
\begin{itemize}
\item{ 
\index{equals(Object)}
\hypertarget{it.unisa.sesa.repominer.db.entities.Type.equals(java.lang.Object)}{{\bf  equals}\\}
\texttt{public boolean\ {\bf  equals}(\texttt{java.lang.Object} {\bf  arg0})
\label{it.unisa.sesa.repominer.db.entities.Type.equals(java.lang.Object)}}%end signature
}%end item
\item{ 
\index{getHeaderFileLocation()}
\hypertarget{it.unisa.sesa.repominer.db.entities.Type.getHeaderFileLocation()}{{\bf  getHeaderFileLocation}\\}
\texttt{public java.lang.String\ {\bf  getHeaderFileLocation}()
\label{it.unisa.sesa.repominer.db.entities.Type.getHeaderFileLocation()}}%end signature
}%end item
\item{ 
\index{getId()}
\hypertarget{it.unisa.sesa.repominer.db.entities.Type.getId()}{{\bf  getId}\\}
\texttt{public java.lang.Integer\ {\bf  getId}()
\label{it.unisa.sesa.repominer.db.entities.Type.getId()}}%end signature
}%end item
\item{ 
\index{getImportId()}
\hypertarget{it.unisa.sesa.repominer.db.entities.Type.getImportId()}{{\bf  getImportId}\\}
\texttt{public java.lang.Integer\ {\bf  getImportId}()
\label{it.unisa.sesa.repominer.db.entities.Type.getImportId()}}%end signature
}%end item
\item{ 
\index{getLinesNumber()}
\hypertarget{it.unisa.sesa.repominer.db.entities.Type.getLinesNumber()}{{\bf  getLinesNumber}\\}
\texttt{public java.lang.Integer\ {\bf  getLinesNumber}()
\label{it.unisa.sesa.repominer.db.entities.Type.getLinesNumber()}}%end signature
}%end item
\item{ 
\index{getSourceContainer()}
\hypertarget{it.unisa.sesa.repominer.db.entities.Type.getSourceContainer()}{{\bf  getSourceContainer}\\}
\texttt{public java.lang.Integer\ {\bf  getSourceContainer}()
\label{it.unisa.sesa.repominer.db.entities.Type.getSourceContainer()}}%end signature
}%end item
\item{ 
\index{getSrcFileLocation()}
\hypertarget{it.unisa.sesa.repominer.db.entities.Type.getSrcFileLocation()}{{\bf  getSrcFileLocation}\\}
\texttt{public java.lang.String\ {\bf  getSrcFileLocation}()
\label{it.unisa.sesa.repominer.db.entities.Type.getSrcFileLocation()}}%end signature
}%end item
\item{ 
\index{hashCode()}
\hypertarget{it.unisa.sesa.repominer.db.entities.Type.hashCode()}{{\bf  hashCode}\\}
\texttt{public native int\ {\bf  hashCode}()
\label{it.unisa.sesa.repominer.db.entities.Type.hashCode()}}%end signature
}%end item
\item{ 
\index{setHeaderFileLocation(String)}
\hypertarget{it.unisa.sesa.repominer.db.entities.Type.setHeaderFileLocation(java.lang.String)}{{\bf  setHeaderFileLocation}\\}
\texttt{public void\ {\bf  setHeaderFileLocation}(\texttt{java.lang.String} {\bf  headerFileLocation})
\label{it.unisa.sesa.repominer.db.entities.Type.setHeaderFileLocation(java.lang.String)}}%end signature
}%end item
\item{ 
\index{setId(Integer)}
\hypertarget{it.unisa.sesa.repominer.db.entities.Type.setId(java.lang.Integer)}{{\bf  setId}\\}
\texttt{public void\ {\bf  setId}(\texttt{java.lang.Integer} {\bf  id})
\label{it.unisa.sesa.repominer.db.entities.Type.setId(java.lang.Integer)}}%end signature
}%end item
\item{ 
\index{setImportId(Integer)}
\hypertarget{it.unisa.sesa.repominer.db.entities.Type.setImportId(java.lang.Integer)}{{\bf  setImportId}\\}
\texttt{public void\ {\bf  setImportId}(\texttt{java.lang.Integer} {\bf  importId})
\label{it.unisa.sesa.repominer.db.entities.Type.setImportId(java.lang.Integer)}}%end signature
}%end item
\item{ 
\index{setLinesNumber(Integer)}
\hypertarget{it.unisa.sesa.repominer.db.entities.Type.setLinesNumber(java.lang.Integer)}{{\bf  setLinesNumber}\\}
\texttt{public void\ {\bf  setLinesNumber}(\texttt{java.lang.Integer} {\bf  linesNumber})
\label{it.unisa.sesa.repominer.db.entities.Type.setLinesNumber(java.lang.Integer)}}%end signature
}%end item
\item{ 
\index{setSourceContainer(Integer)}
\hypertarget{it.unisa.sesa.repominer.db.entities.Type.setSourceContainer(java.lang.Integer)}{{\bf  setSourceContainer}\\}
\texttt{public void\ {\bf  setSourceContainer}(\texttt{java.lang.Integer} {\bf  sourceContainer})
\label{it.unisa.sesa.repominer.db.entities.Type.setSourceContainer(java.lang.Integer)}}%end signature
}%end item
\item{ 
\index{setSrcFileLocation(String)}
\hypertarget{it.unisa.sesa.repominer.db.entities.Type.setSrcFileLocation(java.lang.String)}{{\bf  setSrcFileLocation}\\}
\texttt{public void\ {\bf  setSrcFileLocation}(\texttt{java.lang.String} {\bf  srcFileLocation})
\label{it.unisa.sesa.repominer.db.entities.Type.setSrcFileLocation(java.lang.String)}}%end signature
}%end item
\item{ 
\index{toString()}
\hypertarget{it.unisa.sesa.repominer.db.entities.Type.toString()}{{\bf  toString}\\}
\texttt{public java.lang.String\ {\bf  toString}()
\label{it.unisa.sesa.repominer.db.entities.Type.toString()}}%end signature
}%end item
\end{itemize}
}
}
\subsection{\label{it.unisa.sesa.repominer.db.entities.TypeMetric}Class TypeMetric}{
\hypertarget{it.unisa.sesa.repominer.db.entities.TypeMetric}{}\vskip .1in 
\subsubsection{Declaration}{
\small public class TypeMetric
\\ {\bf  extends} java.lang.Object
\refdefined{java.lang.Object}}
\subsubsection{Constructor summary}{
\begin{verse}
\hyperlink{it.unisa.sesa.repominer.db.entities.TypeMetric()}{{\bf TypeMetric()}} \\
\hyperlink{it.unisa.sesa.repominer.db.entities.TypeMetric(java.lang.Integer, java.lang.Integer, java.lang.Double)}{{\bf TypeMetric(Integer, Integer, Double)}} \\
\end{verse}
}
\subsubsection{Method summary}{
\begin{verse}
\hyperlink{it.unisa.sesa.repominer.db.entities.TypeMetric.equals(java.lang.Object)}{{\bf equals(Object)}} \\
\hyperlink{it.unisa.sesa.repominer.db.entities.TypeMetric.getMetricId()}{{\bf getMetricId()}} \\
\hyperlink{it.unisa.sesa.repominer.db.entities.TypeMetric.getTypeId()}{{\bf getTypeId()}} \\
\hyperlink{it.unisa.sesa.repominer.db.entities.TypeMetric.getValue()}{{\bf getValue()}} \\
\hyperlink{it.unisa.sesa.repominer.db.entities.TypeMetric.hashCode()}{{\bf hashCode()}} \\
\hyperlink{it.unisa.sesa.repominer.db.entities.TypeMetric.setMetricId(java.lang.Integer)}{{\bf setMetricId(Integer)}} \\
\hyperlink{it.unisa.sesa.repominer.db.entities.TypeMetric.setTypeId(java.lang.Integer)}{{\bf setTypeId(Integer)}} \\
\hyperlink{it.unisa.sesa.repominer.db.entities.TypeMetric.setValue(java.lang.Double)}{{\bf setValue(Double)}} \\
\hyperlink{it.unisa.sesa.repominer.db.entities.TypeMetric.toString()}{{\bf toString()}} \\
\end{verse}
}
\subsubsection{Constructors}{
\vskip -2em
\begin{itemize}
\item{ 
\index{TypeMetric()}
\hypertarget{it.unisa.sesa.repominer.db.entities.TypeMetric()}{{\bf  TypeMetric}\\}
\texttt{public\ {\bf  TypeMetric}()
\label{it.unisa.sesa.repominer.db.entities.TypeMetric()}}%end signature
}%end item
\item{ 
\index{TypeMetric(Integer, Integer, Double)}
\hypertarget{it.unisa.sesa.repominer.db.entities.TypeMetric(java.lang.Integer, java.lang.Integer, java.lang.Double)}{{\bf  TypeMetric}\\}
\texttt{public\ {\bf  TypeMetric}(\texttt{java.lang.Integer} {\bf  typeId},
\texttt{java.lang.Integer} {\bf  metricId},
\texttt{java.lang.Double} {\bf  value})
\label{it.unisa.sesa.repominer.db.entities.TypeMetric(java.lang.Integer, java.lang.Integer, java.lang.Double)}}%end signature
}%end item
\end{itemize}
}
\subsubsection{Methods}{
\vskip -2em
\begin{itemize}
\item{ 
\index{equals(Object)}
\hypertarget{it.unisa.sesa.repominer.db.entities.TypeMetric.equals(java.lang.Object)}{{\bf  equals}\\}
\texttt{public boolean\ {\bf  equals}(\texttt{java.lang.Object} {\bf  arg0})
\label{it.unisa.sesa.repominer.db.entities.TypeMetric.equals(java.lang.Object)}}%end signature
}%end item
\item{ 
\index{getMetricId()}
\hypertarget{it.unisa.sesa.repominer.db.entities.TypeMetric.getMetricId()}{{\bf  getMetricId}\\}
\texttt{public java.lang.Integer\ {\bf  getMetricId}()
\label{it.unisa.sesa.repominer.db.entities.TypeMetric.getMetricId()}}%end signature
}%end item
\item{ 
\index{getTypeId()}
\hypertarget{it.unisa.sesa.repominer.db.entities.TypeMetric.getTypeId()}{{\bf  getTypeId}\\}
\texttt{public java.lang.Integer\ {\bf  getTypeId}()
\label{it.unisa.sesa.repominer.db.entities.TypeMetric.getTypeId()}}%end signature
}%end item
\item{ 
\index{getValue()}
\hypertarget{it.unisa.sesa.repominer.db.entities.TypeMetric.getValue()}{{\bf  getValue}\\}
\texttt{public java.lang.Double\ {\bf  getValue}()
\label{it.unisa.sesa.repominer.db.entities.TypeMetric.getValue()}}%end signature
}%end item
\item{ 
\index{hashCode()}
\hypertarget{it.unisa.sesa.repominer.db.entities.TypeMetric.hashCode()}{{\bf  hashCode}\\}
\texttt{public native int\ {\bf  hashCode}()
\label{it.unisa.sesa.repominer.db.entities.TypeMetric.hashCode()}}%end signature
}%end item
\item{ 
\index{setMetricId(Integer)}
\hypertarget{it.unisa.sesa.repominer.db.entities.TypeMetric.setMetricId(java.lang.Integer)}{{\bf  setMetricId}\\}
\texttt{public void\ {\bf  setMetricId}(\texttt{java.lang.Integer} {\bf  metricId})
\label{it.unisa.sesa.repominer.db.entities.TypeMetric.setMetricId(java.lang.Integer)}}%end signature
}%end item
\item{ 
\index{setTypeId(Integer)}
\hypertarget{it.unisa.sesa.repominer.db.entities.TypeMetric.setTypeId(java.lang.Integer)}{{\bf  setTypeId}\\}
\texttt{public void\ {\bf  setTypeId}(\texttt{java.lang.Integer} {\bf  typeId})
\label{it.unisa.sesa.repominer.db.entities.TypeMetric.setTypeId(java.lang.Integer)}}%end signature
}%end item
\item{ 
\index{setValue(Double)}
\hypertarget{it.unisa.sesa.repominer.db.entities.TypeMetric.setValue(java.lang.Double)}{{\bf  setValue}\\}
\texttt{public void\ {\bf  setValue}(\texttt{java.lang.Double} {\bf  value})
\label{it.unisa.sesa.repominer.db.entities.TypeMetric.setValue(java.lang.Double)}}%end signature
}%end item
\item{ 
\index{toString()}
\hypertarget{it.unisa.sesa.repominer.db.entities.TypeMetric.toString()}{{\bf  toString}\\}
\texttt{public java.lang.String\ {\bf  toString}()
\label{it.unisa.sesa.repominer.db.entities.TypeMetric.toString()}}%end signature
}%end item
\end{itemize}
}
}
}
\section{Package it.unisa.sesa.repominer.dbscan}{
\label{it.unisa.sesa.repominer.dbscan}\hypertarget{it.unisa.sesa.repominer.dbscan}{}
\hskip -.05in
\hbox to \hsize{\textit{ Package Contents\hfil Page}}
\vskip .13in
\hbox{{\bf  Classes}}
\entityintro{ChangePoint}{it.unisa.sesa.repominer.dbscan.ChangePoint}{An implementation of \hyperlink{it.unisa.sesa.repominer.db.entities.Change}{}{\small 
\refdefined{it.unisa.sesa.repominer.db.entities.Change}} for points with integer coordinates A single \hyperlink{it.unisa.sesa.repominer.dbscan.ChangePoint}{}{\small 
\refdefined{it.unisa.sesa.repominer.dbscan.ChangePoint}} has two instance variables, a \hyperlink{it.unisa.sesa.repominer.db.entities.Change}{}{\small 
\refdefined{it.unisa.sesa.repominer.db.entities.Change}} instance and an integer coordinate}
\entityintro{Cluster}{it.unisa.sesa.repominer.dbscan.Cluster}{Cluster holding a set of points}
\entityintro{DBSCAN}{it.unisa.sesa.repominer.dbscan.DBSCAN}{DBSCAN (density-based spatial clustering of applications with noise) algorithm

The DBSCAN algorithm forms clusters based on the idea of density connectivity, i.e. a point p is density connected to another point q, if there exists a chain of points p$_{i}$, with i=1...n and p$_{1}$=p and p$_{n}$=q, such that each pair p$_{i+1}$ is directly density-reachable from p$_{i}$.}
\vskip .1in
\vskip .1in
\subsection{\label{it.unisa.sesa.repominer.dbscan.ChangePoint}Class ChangePoint}{
\hypertarget{it.unisa.sesa.repominer.dbscan.ChangePoint}{}\vskip .1in 
An implementation of \hyperlink{it.unisa.sesa.repominer.db.entities.Change}{}{\small 
\refdefined{it.unisa.sesa.repominer.db.entities.Change}} for points with integer coordinates A single \hyperlink{it.unisa.sesa.repominer.dbscan.ChangePoint}{}{\small 
\refdefined{it.unisa.sesa.repominer.dbscan.ChangePoint}} has two instance variables, a \hyperlink{it.unisa.sesa.repominer.db.entities.Change}{}{\small 
\refdefined{it.unisa.sesa.repominer.db.entities.Change}} instance and an integer coordinate\vskip .1in 
\subsubsection{Declaration}{
\small public class ChangePoint
\\ {\bf  extends} java.lang.Object
\refdefined{java.lang.Object}}
\subsubsection{Constructor summary}{
\begin{verse}
\hyperlink{it.unisa.sesa.repominer.dbscan.ChangePoint(int, it.unisa.sesa.repominer.db.entities.Change)}{{\bf ChangePoint(int, Change)}} Build an instance of a point, wrapping a single \hyperlink{it.unisa.sesa.repominer.db.entities.Change}{}{\small 
\refdefined{it.unisa.sesa.repominer.db.entities.Change}} instance its coordinate\\
\end{verse}
}
\subsubsection{Method summary}{
\begin{verse}
\hyperlink{it.unisa.sesa.repominer.dbscan.ChangePoint.distanceFrom(it.unisa.sesa.repominer.dbscan.ChangePoint)}{{\bf distanceFrom(ChangePoint)}} Calculate distance between two ChangePoints\\
\hyperlink{it.unisa.sesa.repominer.dbscan.ChangePoint.equals(java.lang.Object)}{{\bf equals(Object)}} \\
\hyperlink{it.unisa.sesa.repominer.dbscan.ChangePoint.getChange()}{{\bf getChange()}} This method returns the \hyperlink{it.unisa.sesa.repominer.db.entities.Change}{}{\small 
\refdefined{it.unisa.sesa.repominer.db.entities.Change}} wrapped by this point\\
\hyperlink{it.unisa.sesa.repominer.dbscan.ChangePoint.getVisitedFlag()}{{\bf getVisitedFlag()}} This method return the flag value used to checking the state of this point; 0 means the point has not been visited; 1 if the point has been classified as noise; 2 if the point has been already put in a cluster\\
\hyperlink{it.unisa.sesa.repominer.dbscan.ChangePoint.getX()}{{\bf getX()}} This method returns the coordinate for this point\\
\hyperlink{it.unisa.sesa.repominer.dbscan.ChangePoint.hashCode()}{{\bf hashCode()}} \\
\hyperlink{it.unisa.sesa.repominer.dbscan.ChangePoint.isAlreadyInACluster()}{{\bf isAlreadyInACluster()}} \\
\hyperlink{it.unisa.sesa.repominer.dbscan.ChangePoint.isNoise()}{{\bf isNoise()}} This method check if this point has been classified as noise\\
\hyperlink{it.unisa.sesa.repominer.dbscan.ChangePoint.isNotVisited()}{{\bf isNotVisited()}} This method check if this point has been already visited by the algorithm\\
\hyperlink{it.unisa.sesa.repominer.dbscan.ChangePoint.setAlreadyInACluster()}{{\bf setAlreadyInACluster()}} Set the point as already in a cluster (flag value 2)\\
\hyperlink{it.unisa.sesa.repominer.dbscan.ChangePoint.setChange(it.unisa.sesa.repominer.db.entities.Change)}{{\bf setChange(Change)}} This method set the \hyperlink{it.unisa.sesa.repominer.db.entities.Change}{}{\small 
\refdefined{it.unisa.sesa.repominer.db.entities.Change}} specified by this point\\
\hyperlink{it.unisa.sesa.repominer.dbscan.ChangePoint.setNoise()}{{\bf setNoise()}} Set the point as noise (flag value 1)\\
\hyperlink{it.unisa.sesa.repominer.dbscan.ChangePoint.toString()}{{\bf toString()}} \\
\end{verse}
}
\subsubsection{Constructors}{
\vskip -2em
\begin{itemize}
\item{ 
\index{ChangePoint(int, Change)}
\hypertarget{it.unisa.sesa.repominer.dbscan.ChangePoint(int, it.unisa.sesa.repominer.db.entities.Change)}{{\bf  ChangePoint}\\}
\texttt{public\ {\bf  ChangePoint}(\texttt{int} {\bf  x},
\texttt{it.unisa.sesa.repominer.db.entities.Change} {\bf  pChange})
\label{it.unisa.sesa.repominer.dbscan.ChangePoint(int, it.unisa.sesa.repominer.db.entities.Change)}}%end signature
\begin{itemize}
\item{
{\bf  Description}

Build an instance of a point, wrapping a single \hyperlink{it.unisa.sesa.repominer.db.entities.Change}{}{\small 
\refdefined{it.unisa.sesa.repominer.db.entities.Change}} instance its coordinate
}
\item{
{\bf  Parameters}
  \begin{itemize}
   \item{
\texttt{x} -- }
   \item{
\texttt{pChange} -- }
  \end{itemize}
}%end item
\end{itemize}
}%end item
\end{itemize}
}
\subsubsection{Methods}{
\vskip -2em
\begin{itemize}
\item{ 
\index{distanceFrom(ChangePoint)}
\hypertarget{it.unisa.sesa.repominer.dbscan.ChangePoint.distanceFrom(it.unisa.sesa.repominer.dbscan.ChangePoint)}{{\bf  distanceFrom}\\}
\texttt{public double\ {\bf  distanceFrom}(\texttt{ChangePoint} {\bf  pChangePoint})
\label{it.unisa.sesa.repominer.dbscan.ChangePoint.distanceFrom(it.unisa.sesa.repominer.dbscan.ChangePoint)}}%end signature
\begin{itemize}
\item{
{\bf  Description}

Calculate distance between two ChangePoints
}
\item{
{\bf  Parameters}
  \begin{itemize}
   \item{
\texttt{pChangePoint} -- the point used to calculate distance from this point}
  \end{itemize}
}%end item
\item{{\bf  Returns} -- 
A double value that represent distance between two points 
}%end item
\end{itemize}
}%end item
\item{ 
\index{equals(Object)}
\hypertarget{it.unisa.sesa.repominer.dbscan.ChangePoint.equals(java.lang.Object)}{{\bf  equals}\\}
\texttt{public boolean\ {\bf  equals}(\texttt{java.lang.Object} {\bf  arg0})
\label{it.unisa.sesa.repominer.dbscan.ChangePoint.equals(java.lang.Object)}}%end signature
}%end item
\item{ 
\index{getChange()}
\hypertarget{it.unisa.sesa.repominer.dbscan.ChangePoint.getChange()}{{\bf  getChange}\\}
\texttt{public it.unisa.sesa.repominer.db.entities.Change\ {\bf  getChange}()
\label{it.unisa.sesa.repominer.dbscan.ChangePoint.getChange()}}%end signature
\begin{itemize}
\item{
{\bf  Description}

This method returns the \hyperlink{it.unisa.sesa.repominer.db.entities.Change}{}{\small 
\refdefined{it.unisa.sesa.repominer.db.entities.Change}} wrapped by this point
}
\item{{\bf  Returns} -- 
a \hyperlink{it.unisa.sesa.repominer.db.entities.Change}{}{\small 
\refdefined{it.unisa.sesa.repominer.db.entities.Change}} 
}%end item
\end{itemize}
}%end item
\item{ 
\index{getVisitedFlag()}
\hypertarget{it.unisa.sesa.repominer.dbscan.ChangePoint.getVisitedFlag()}{{\bf  getVisitedFlag}\\}
\texttt{public int\ {\bf  getVisitedFlag}()
\label{it.unisa.sesa.repominer.dbscan.ChangePoint.getVisitedFlag()}}%end signature
\begin{itemize}
\item{
{\bf  Description}

This method return the flag value used to checking the state of this point; 0 means the point has not been visited; 1 if the point has been classified as noise; 2 if the point has been already put in a cluster
}
\item{{\bf  Returns} -- 
the valued of the flag for this point 
}%end item
\end{itemize}
}%end item
\item{ 
\index{getX()}
\hypertarget{it.unisa.sesa.repominer.dbscan.ChangePoint.getX()}{{\bf  getX}\\}
\texttt{public int\ {\bf  getX}()
\label{it.unisa.sesa.repominer.dbscan.ChangePoint.getX()}}%end signature
\begin{itemize}
\item{
{\bf  Description}

This method returns the coordinate for this point
}
\item{{\bf  Returns} -- 
the coordinate for this point 
}%end item
\end{itemize}
}%end item
\item{ 
\index{hashCode()}
\hypertarget{it.unisa.sesa.repominer.dbscan.ChangePoint.hashCode()}{{\bf  hashCode}\\}
\texttt{public native int\ {\bf  hashCode}()
\label{it.unisa.sesa.repominer.dbscan.ChangePoint.hashCode()}}%end signature
}%end item
\item{ 
\index{isAlreadyInACluster()}
\hypertarget{it.unisa.sesa.repominer.dbscan.ChangePoint.isAlreadyInACluster()}{{\bf  isAlreadyInACluster}\\}
\texttt{public boolean\ {\bf  isAlreadyInACluster}()
\label{it.unisa.sesa.repominer.dbscan.ChangePoint.isAlreadyInACluster()}}%end signature
}%end item
\item{ 
\index{isNoise()}
\hypertarget{it.unisa.sesa.repominer.dbscan.ChangePoint.isNoise()}{{\bf  isNoise}\\}
\texttt{public boolean\ {\bf  isNoise}()
\label{it.unisa.sesa.repominer.dbscan.ChangePoint.isNoise()}}%end signature
\begin{itemize}
\item{
{\bf  Description}

This method check if this point has been classified as noise
}
\item{{\bf  Returns} -- 
true if the point is a noise point; false otherwise 
}%end item
\end{itemize}
}%end item
\item{ 
\index{isNotVisited()}
\hypertarget{it.unisa.sesa.repominer.dbscan.ChangePoint.isNotVisited()}{{\bf  isNotVisited}\\}
\texttt{public boolean\ {\bf  isNotVisited}()
\label{it.unisa.sesa.repominer.dbscan.ChangePoint.isNotVisited()}}%end signature
\begin{itemize}
\item{
{\bf  Description}

This method check if this point has been already visited by the algorithm
}
\item{{\bf  Returns} -- 
true if this point has been visited; false otherwise 
}%end item
\end{itemize}
}%end item
\item{ 
\index{setAlreadyInACluster()}
\hypertarget{it.unisa.sesa.repominer.dbscan.ChangePoint.setAlreadyInACluster()}{{\bf  setAlreadyInACluster}\\}
\texttt{public void\ {\bf  setAlreadyInACluster}()
\label{it.unisa.sesa.repominer.dbscan.ChangePoint.setAlreadyInACluster()}}%end signature
\begin{itemize}
\item{
{\bf  Description}

Set the point as already in a cluster (flag value 2)
}
\end{itemize}
}%end item
\item{ 
\index{setChange(Change)}
\hypertarget{it.unisa.sesa.repominer.dbscan.ChangePoint.setChange(it.unisa.sesa.repominer.db.entities.Change)}{{\bf  setChange}\\}
\texttt{public void\ {\bf  setChange}(\texttt{it.unisa.sesa.repominer.db.entities.Change} {\bf  change})
\label{it.unisa.sesa.repominer.dbscan.ChangePoint.setChange(it.unisa.sesa.repominer.db.entities.Change)}}%end signature
\begin{itemize}
\item{
{\bf  Description}

This method set the \hyperlink{it.unisa.sesa.repominer.db.entities.Change}{}{\small 
\refdefined{it.unisa.sesa.repominer.db.entities.Change}} specified by this point
}
\item{
{\bf  Parameters}
  \begin{itemize}
   \item{
\texttt{change} -- which the point refers}
  \end{itemize}
}%end item
\end{itemize}
}%end item
\item{ 
\index{setNoise()}
\hypertarget{it.unisa.sesa.repominer.dbscan.ChangePoint.setNoise()}{{\bf  setNoise}\\}
\texttt{public void\ {\bf  setNoise}()
\label{it.unisa.sesa.repominer.dbscan.ChangePoint.setNoise()}}%end signature
\begin{itemize}
\item{
{\bf  Description}

Set the point as noise (flag value 1)
}
\end{itemize}
}%end item
\item{ 
\index{toString()}
\hypertarget{it.unisa.sesa.repominer.dbscan.ChangePoint.toString()}{{\bf  toString}\\}
\texttt{public java.lang.String\ {\bf  toString}()
\label{it.unisa.sesa.repominer.dbscan.ChangePoint.toString()}}%end signature
}%end item
\end{itemize}
}
}
\subsection{\label{it.unisa.sesa.repominer.dbscan.Cluster}Class Cluster}{
\hypertarget{it.unisa.sesa.repominer.dbscan.Cluster}{}\vskip .1in 
Cluster holding a set of points\vskip .1in 
\subsubsection{Declaration}{
\small public class Cluster
\\ {\bf  extends} java.lang.Object
\refdefined{java.lang.Object}}
\subsubsection{Constructor summary}{
\begin{verse}
\hyperlink{it.unisa.sesa.repominer.dbscan.Cluster(it.unisa.sesa.repominer.dbscan.ChangePoint)}{{\bf Cluster(ChangePoint)}} Build a cluster centered at specified point\\
\end{verse}
}
\subsubsection{Method summary}{
\begin{verse}
\hyperlink{it.unisa.sesa.repominer.dbscan.Cluster.addPoint(it.unisa.sesa.repominer.dbscan.ChangePoint)}{{\bf addPoint(ChangePoint)}} Add a point to this cluster\\
\hyperlink{it.unisa.sesa.repominer.dbscan.Cluster.getCenter()}{{\bf getCenter()}} Get the point chosen to be the center of this cluster\\
\hyperlink{it.unisa.sesa.repominer.dbscan.Cluster.getPoints()}{{\bf getPoints()}} Get all points contained in the cluster\\
\end{verse}
}
\subsubsection{Constructors}{
\vskip -2em
\begin{itemize}
\item{ 
\index{Cluster(ChangePoint)}
\hypertarget{it.unisa.sesa.repominer.dbscan.Cluster(it.unisa.sesa.repominer.dbscan.ChangePoint)}{{\bf  Cluster}\\}
\texttt{public\ {\bf  Cluster}(\texttt{ChangePoint} {\bf  center})
\label{it.unisa.sesa.repominer.dbscan.Cluster(it.unisa.sesa.repominer.dbscan.ChangePoint)}}%end signature
\begin{itemize}
\item{
{\bf  Description}

Build a cluster centered at specified point
}
\item{
{\bf  Parameters}
  \begin{itemize}
   \item{
\texttt{center} -- the point which is to be the center of this cluster}
  \end{itemize}
}%end item
\end{itemize}
}%end item
\end{itemize}
}
\subsubsection{Methods}{
\vskip -2em
\begin{itemize}
\item{ 
\index{addPoint(ChangePoint)}
\hypertarget{it.unisa.sesa.repominer.dbscan.Cluster.addPoint(it.unisa.sesa.repominer.dbscan.ChangePoint)}{{\bf  addPoint}\\}
\texttt{public void\ {\bf  addPoint}(\texttt{ChangePoint} {\bf  pPoint})
\label{it.unisa.sesa.repominer.dbscan.Cluster.addPoint(it.unisa.sesa.repominer.dbscan.ChangePoint)}}%end signature
\begin{itemize}
\item{
{\bf  Description}

Add a point to this cluster
}
\item{
{\bf  Parameters}
  \begin{itemize}
   \item{
\texttt{pPoint} -- point to add}
  \end{itemize}
}%end item
\end{itemize}
}%end item
\item{ 
\index{getCenter()}
\hypertarget{it.unisa.sesa.repominer.dbscan.Cluster.getCenter()}{{\bf  getCenter}\\}
\texttt{public ChangePoint\ {\bf  getCenter}()
\label{it.unisa.sesa.repominer.dbscan.Cluster.getCenter()}}%end signature
\begin{itemize}
\item{
{\bf  Description}

Get the point chosen to be the center of this cluster
}
\item{{\bf  Returns} -- 
chosen cluster center 
}%end item
\end{itemize}
}%end item
\item{ 
\index{getPoints()}
\hypertarget{it.unisa.sesa.repominer.dbscan.Cluster.getPoints()}{{\bf  getPoints}\\}
\texttt{public java.util.List\ {\bf  getPoints}()
\label{it.unisa.sesa.repominer.dbscan.Cluster.getPoints()}}%end signature
\begin{itemize}
\item{
{\bf  Description}

Get all points contained in the cluster
}
\item{{\bf  Returns} -- 
points contained in the cluster 
}%end item
\end{itemize}
}%end item
\end{itemize}
}
}
\subsection{\label{it.unisa.sesa.repominer.dbscan.DBSCAN}Class DBSCAN}{
\hypertarget{it.unisa.sesa.repominer.dbscan.DBSCAN}{}\vskip .1in 
DBSCAN (density-based spatial clustering of applications with noise) algorithm

The DBSCAN algorithm forms clusters based on the idea of density connectivity, i.e. a point p is density connected to another point q, if there exists a chain of points p$_{i}$, with i=1...n and p$_{1}$=p and p$_{n}$=q, such that each pair p$_{i+1}$ is directly density-reachable from p$_{i}$. A point q is directly density-reachable from point p if it is in the $\epsilon$-neighborhood or this point.

Any point that is not density-reachable from a formed cluster is treated as noise, and will thus not be present in the results.

The algorithm requires two parameters:\begin{itemize}
\item{\vskip -.8ex eps: the distance that defines the}
\item{\vskip -.8ex epsilon-neighborhood of a point minPoints: the minimum number of density-connected points required to form a cluster}
\end{itemize}
\vskip .1in 
\subsubsection{Declaration}{
\small public class DBSCAN
\\ {\bf  extends} java.lang.Object
\refdefined{java.lang.Object}}
\subsubsection{Constructor summary}{
\begin{verse}
\hyperlink{it.unisa.sesa.repominer.dbscan.DBSCAN(double, int)}{{\bf DBSCAN(double, int)}} Creates a new instance of a DBSCAN\\
\end{verse}
}
\subsubsection{Method summary}{
\begin{verse}
\hyperlink{it.unisa.sesa.repominer.dbscan.DBSCAN.cluster(java.util.List)}{{\bf cluster(List)}} Perform DBSCAN cluster analysis\\
\hyperlink{it.unisa.sesa.repominer.dbscan.DBSCAN.getEps()}{{\bf getEps()}} Return the eps value (maximum radius of the neighborhood to be considered)\\
\hyperlink{it.unisa.sesa.repominer.dbscan.DBSCAN.getMinPoints()}{{\bf getMinPoints()}} Returns the minimum number of points needed for a cluster\\
\end{verse}
}
\subsubsection{Constructors}{
\vskip -2em
\begin{itemize}
\item{ 
\index{DBSCAN(double, int)}
\hypertarget{it.unisa.sesa.repominer.dbscan.DBSCAN(double, int)}{{\bf  DBSCAN}\\}
\texttt{public\ {\bf  DBSCAN}(\texttt{double} {\bf  eps},
\texttt{int} {\bf  minPoints})
\label{it.unisa.sesa.repominer.dbscan.DBSCAN(double, int)}}%end signature
\begin{itemize}
\item{
{\bf  Description}

Creates a new instance of a DBSCAN
}
\item{
{\bf  Parameters}
  \begin{itemize}
   \item{
\texttt{eps} -- maximum radius of the neighborhood to be considered}
   \item{
\texttt{minPoints} -- minimum number of points needed for a cluster}
  \end{itemize}
}%end item
\end{itemize}
}%end item
\end{itemize}
}
\subsubsection{Methods}{
\vskip -2em
\begin{itemize}
\item{ 
\index{cluster(List)}
\hypertarget{it.unisa.sesa.repominer.dbscan.DBSCAN.cluster(java.util.List)}{{\bf  cluster}\\}
\texttt{public java.util.List\ {\bf  cluster}(\texttt{java.util.List} {\bf  pPoints})
\label{it.unisa.sesa.repominer.dbscan.DBSCAN.cluster(java.util.List)}}%end signature
\begin{itemize}
\item{
{\bf  Description}

Perform DBSCAN cluster analysis
}
\item{
{\bf  Parameters}
  \begin{itemize}
   \item{
\texttt{pPoints} -- the points to cluster}
  \end{itemize}
}%end item
\item{{\bf  Returns} -- 
the list of Cluster 
}%end item
\end{itemize}
}%end item
\item{ 
\index{getEps()}
\hypertarget{it.unisa.sesa.repominer.dbscan.DBSCAN.getEps()}{{\bf  getEps}\\}
\texttt{public double\ {\bf  getEps}()
\label{it.unisa.sesa.repominer.dbscan.DBSCAN.getEps()}}%end signature
\begin{itemize}
\item{
{\bf  Description}

Return the eps value (maximum radius of the neighborhood to be considered)
}
\item{{\bf  Returns} -- 
maximum radius of the neighborhood 
}%end item
\end{itemize}
}%end item
\item{ 
\index{getMinPoints()}
\hypertarget{it.unisa.sesa.repominer.dbscan.DBSCAN.getMinPoints()}{{\bf  getMinPoints}\\}
\texttt{public int\ {\bf  getMinPoints}()
\label{it.unisa.sesa.repominer.dbscan.DBSCAN.getMinPoints()}}%end signature
\begin{itemize}
\item{
{\bf  Description}

Returns the minimum number of points needed for a cluster
}
\item{{\bf  Returns} -- 
minimum number of points needed for a cluster 
}%end item
\end{itemize}
}%end item
\end{itemize}
}
}
}
\section{Package it.unisa.sesa.repominer.metrics.exception}{
\label{it.unisa.sesa.repominer.metrics.exception}\hypertarget{it.unisa.sesa.repominer.metrics.exception}{}
\hskip -.05in
\hbox to \hsize{\textit{ Package Contents\hfil Page}}
\vskip .1in
\vskip .1in
\subsection{\label{it.unisa.sesa.repominer.metrics.exception.NoChangesException}Exception NoChangesException}{
\hypertarget{it.unisa.sesa.repominer.metrics.exception.NoChangesException}{}\vskip .1in 
This class represents an exception to be thrown when there are not changes stored into the database.\vskip .1in 
\subsubsection{Declaration}{
\small public class NoChangesException
\\ {\bf  extends} java.lang.Exception
\refdefined{java.lang.Exception}}
\subsubsection{Constructor summary}{
\begin{verse}
\hyperlink{it.unisa.sesa.repominer.metrics.exception.NoChangesException()}{{\bf NoChangesException()}} This method constructs an object of type NoChangesException by setting the message to a default string value "No changes found into db".\\
\end{verse}
}
\subsubsection{Constructors}{
\vskip -2em
\begin{itemize}
\item{ 
\index{NoChangesException()}
\hypertarget{it.unisa.sesa.repominer.metrics.exception.NoChangesException()}{{\bf  NoChangesException}\\}
\texttt{public\ {\bf  NoChangesException}()
\label{it.unisa.sesa.repominer.metrics.exception.NoChangesException()}}%end signature
\begin{itemize}
\item{
{\bf  Description}

This method constructs an object of type NoChangesException by setting the message to a default string value "No changes found into db".
}
\end{itemize}
}%end item
\end{itemize}
}
\subsubsection{Members inherited from class Throwable }{
\texttt{java.lang.Throwable} {\small 
\refdefined{java.lang.Throwable}}
{\small 

\vskip -2em
\begin{itemize}
\item{\vskip -1.5ex 
\texttt{public final synchronized void {\bf  addSuppressed}(\texttt{Throwable} {\bf  arg0})
}%end signature
}%end item
\item{\vskip -1.5ex 
\texttt{public synchronized Throwable {\bf  fillInStackTrace}()
}%end signature
}%end item
\item{\vskip -1.5ex 
\texttt{public synchronized Throwable {\bf  getCause}()
}%end signature
}%end item
\item{\vskip -1.5ex 
\texttt{public String {\bf  getLocalizedMessage}()
}%end signature
}%end item
\item{\vskip -1.5ex 
\texttt{public String {\bf  getMessage}()
}%end signature
}%end item
\item{\vskip -1.5ex 
\texttt{public StackTraceElement {\bf  getStackTrace}()
}%end signature
}%end item
\item{\vskip -1.5ex 
\texttt{public final synchronized Throwable {\bf  getSuppressed}()
}%end signature
}%end item
\item{\vskip -1.5ex 
\texttt{public synchronized Throwable {\bf  initCause}(\texttt{Throwable} {\bf  arg0})
}%end signature
}%end item
\item{\vskip -1.5ex 
\texttt{public void {\bf  printStackTrace}()
}%end signature
}%end item
\item{\vskip -1.5ex 
\texttt{public void {\bf  printStackTrace}(\texttt{java.io.PrintStream} {\bf  arg0})
}%end signature
}%end item
\item{\vskip -1.5ex 
\texttt{public void {\bf  printStackTrace}(\texttt{java.io.PrintWriter} {\bf  arg0})
}%end signature
}%end item
\item{\vskip -1.5ex 
\texttt{public void {\bf  setStackTrace}(\texttt{StackTraceElement\lbrack \rbrack } {\bf  arg0})
}%end signature
}%end item
\item{\vskip -1.5ex 
\texttt{public String {\bf  toString}()
}%end signature
}%end item
\end{itemize}
}
}
}
\section{Package it.unisa.sesa.repominer.metrics}{
\label{it.unisa.sesa.repominer.metrics}\hypertarget{it.unisa.sesa.repominer.metrics}{}
\hskip -.05in
\hbox to \hsize{\textit{ Package Contents\hfil Page}}
\vskip .13in
\hbox{{\bf  Classes}}
\entityintro{HistoryMetricsCalculator}{it.unisa.sesa.repominer.metrics.HistoryMetricsCalculator}{This class instantiates and calls methods of \hyperlink{it.unisa.sesa.repominer.metrics.ProjectMetrics}{}{\small 
\refdefined{it.unisa.sesa.repominer.metrics.ProjectMetrics}} and \hyperlink{it.unisa.sesa.repominer.metrics.PackageMetrics}{}{\small 
\refdefined{it.unisa.sesa.repominer.metrics.PackageMetrics}}.}
\entityintro{PackageMetrics}{it.unisa.sesa.repominer.metrics.PackageMetrics}{This class is responsible to calculate all the history metrics relative to a package.}
\entityintro{ProjectMetrics}{it.unisa.sesa.repominer.metrics.ProjectMetrics}{This class is responsible to calculate all the history metrics relative to a project.}
\vskip .1in
\vskip .1in
\subsection{\label{it.unisa.sesa.repominer.metrics.HistoryMetricsCalculator}Class HistoryMetricsCalculator}{
\hypertarget{it.unisa.sesa.repominer.metrics.HistoryMetricsCalculator}{}\vskip .1in 
This class instantiates and calls methods of \hyperlink{it.unisa.sesa.repominer.metrics.ProjectMetrics}{}{\small 
\refdefined{it.unisa.sesa.repominer.metrics.ProjectMetrics}} and \hyperlink{it.unisa.sesa.repominer.metrics.PackageMetrics}{}{\small 
\refdefined{it.unisa.sesa.repominer.metrics.PackageMetrics}}.\vskip .1in 
\subsubsection{Declaration}{
\small public class HistoryMetricsCalculator
\\ {\bf  extends} java.lang.Object
\refdefined{java.lang.Object}}
\subsubsection{Constructor summary}{
\begin{verse}
\hyperlink{it.unisa.sesa.repominer.metrics.HistoryMetricsCalculator()}{{\bf HistoryMetricsCalculator()}} \\
\end{verse}
}
\subsubsection{Method summary}{
\begin{verse}
\hyperlink{it.unisa.sesa.repominer.metrics.HistoryMetricsCalculator.calculateMetrics(it.unisa.sesa.repominer.db.entities.Project)}{{\bf calculateMetrics(Project)}} This method calculate project metrics for project passed as argument\\
\hyperlink{it.unisa.sesa.repominer.metrics.HistoryMetricsCalculator.calculateMetrics(it.unisa.sesa.repominer.db.entities.SourceContainer)}{{\bf calculateMetrics(SourceContainer)}} This method calculate package metrics for package passed as argument\\
\end{verse}
}
\subsubsection{Constructors}{
\vskip -2em
\begin{itemize}
\item{ 
\index{HistoryMetricsCalculator()}
\hypertarget{it.unisa.sesa.repominer.metrics.HistoryMetricsCalculator()}{{\bf  HistoryMetricsCalculator}\\}
\texttt{public\ {\bf  HistoryMetricsCalculator}()
\label{it.unisa.sesa.repominer.metrics.HistoryMetricsCalculator()}}%end signature
}%end item
\end{itemize}
}
\subsubsection{Methods}{
\vskip -2em
\begin{itemize}
\item{ 
\index{calculateMetrics(Project)}
\hypertarget{it.unisa.sesa.repominer.metrics.HistoryMetricsCalculator.calculateMetrics(it.unisa.sesa.repominer.db.entities.Project)}{{\bf  calculateMetrics}\\}
\texttt{public static void\ {\bf  calculateMetrics}(\texttt{it.unisa.sesa.repominer.db.entities.Project} {\bf  pProject})
\label{it.unisa.sesa.repominer.metrics.HistoryMetricsCalculator.calculateMetrics(it.unisa.sesa.repominer.db.entities.Project)}}%end signature
\begin{itemize}
\item{
{\bf  Description}

This method calculate project metrics for project passed as argument
}
\item{
{\bf  Parameters}
  \begin{itemize}
   \item{
\texttt{pProject} -- }
  \end{itemize}
}%end item
\end{itemize}
}%end item
\item{ 
\index{calculateMetrics(SourceContainer)}
\hypertarget{it.unisa.sesa.repominer.metrics.HistoryMetricsCalculator.calculateMetrics(it.unisa.sesa.repominer.db.entities.SourceContainer)}{{\bf  calculateMetrics}\\}
\texttt{public static void\ {\bf  calculateMetrics}(\texttt{it.unisa.sesa.repominer.db.entities.SourceContainer} {\bf  pSourceContainer})
\label{it.unisa.sesa.repominer.metrics.HistoryMetricsCalculator.calculateMetrics(it.unisa.sesa.repominer.db.entities.SourceContainer)}}%end signature
\begin{itemize}
\item{
{\bf  Description}

This method calculate package metrics for package passed as argument
}
\item{
{\bf  Parameters}
  \begin{itemize}
   \item{
\texttt{pSourceContainer} -- }
  \end{itemize}
}%end item
\end{itemize}
}%end item
\end{itemize}
}
}
\subsection{\label{it.unisa.sesa.repominer.metrics.PackageMetrics}Class PackageMetrics}{
\hypertarget{it.unisa.sesa.repominer.metrics.PackageMetrics}{}\vskip .1in 
This class is responsible to calculate all the history metrics relative to a package.\vskip .1in 
\subsubsection{Declaration}{
\small public class PackageMetrics
\\ {\bf  extends} java.lang.Object
\refdefined{java.lang.Object}}
\subsubsection{Constructor summary}{
\begin{verse}
\hyperlink{it.unisa.sesa.repominer.metrics.PackageMetrics()}{{\bf PackageMetrics()}} \\
\end{verse}
}
\subsubsection{Method summary}{
\begin{verse}
\hyperlink{it.unisa.sesa.repominer.metrics.PackageMetrics.getInsertionsAndDelitionsInfo(it.unisa.sesa.repominer.db.entities.SourceContainer)}{{\bf getInsertionsAndDelitionsInfo(SourceContainer)}} This method calculate the Lines metric.\\
\hyperlink{it.unisa.sesa.repominer.metrics.PackageMetrics.getMaxInsertionsOrDeletionsNumber(it.unisa.sesa.repominer.db.entities.SourceContainer)}{{\bf getMaxInsertionsOrDeletionsNumber(SourceContainer)}} This method calculate the second value of Lines metric, namely max number of insertion or deletion\\
\hyperlink{it.unisa.sesa.repominer.metrics.PackageMetrics.getMeanDimensionOfModifiedFiles(it.unisa.sesa.repominer.db.entities.SourceContainer)}{{\bf getMeanDimensionOfModifiedFiles(SourceContainer)}} This method calculate ChangeSetSize metric.\\
\hyperlink{it.unisa.sesa.repominer.metrics.PackageMetrics.getMeanInsertionsOrDeletionsNumber(it.unisa.sesa.repominer.db.entities.SourceContainer)}{{\bf getMeanInsertionsOrDeletionsNumber(SourceContainer)}} This method calculate the third value of Lines metric, namely mean number of insertion or deletion\\
\hyperlink{it.unisa.sesa.repominer.metrics.PackageMetrics.getMeanNumberOfChange(it.unisa.sesa.repominer.db.entities.SourceContainer)}{{\bf getMeanNumberOfChange(SourceContainer)}} This method calculate mean\_NCHANGE metric.\\
\hyperlink{it.unisa.sesa.repominer.metrics.PackageMetrics.getMeanNumberOfChangeForBugFix(it.unisa.sesa.repominer.db.entities.SourceContainer)}{{\bf getMeanNumberOfChangeForBugFix(SourceContainer)}} This method calculate mean\_NFIX metric.\\
\hyperlink{it.unisa.sesa.repominer.metrics.PackageMetrics.getMeanNumberOfChangeForRefactoring(it.unisa.sesa.repominer.db.entities.SourceContainer)}{{\bf getMeanNumberOfChangeForRefactoring(SourceContainer)}} This method calculate mean\_NREF metric.\\
\hyperlink{it.unisa.sesa.repominer.metrics.PackageMetrics.getNumberOfAuthor(it.unisa.sesa.repominer.db.entities.SourceContainer)}{{\bf getNumberOfAuthor(SourceContainer)}} This method calculate NR metric.\\
\hyperlink{it.unisa.sesa.repominer.metrics.PackageMetrics.getTotalInsertionsOrDeletionsNumber(it.unisa.sesa.repominer.db.entities.SourceContainer)}{{\bf getTotalInsertionsOrDeletionsNumber(SourceContainer)}} This method calculate the first value of Lines metric, namely sum of insertion or deletion\\
\end{verse}
}
\subsubsection{Constructors}{
\vskip -2em
\begin{itemize}
\item{ 
\index{PackageMetrics()}
\hypertarget{it.unisa.sesa.repominer.metrics.PackageMetrics()}{{\bf  PackageMetrics}\\}
\texttt{public\ {\bf  PackageMetrics}()
\label{it.unisa.sesa.repominer.metrics.PackageMetrics()}}%end signature
}%end item
\end{itemize}
}
\subsubsection{Methods}{
\vskip -2em
\begin{itemize}
\item{ 
\index{getInsertionsAndDelitionsInfo(SourceContainer)}
\hypertarget{it.unisa.sesa.repominer.metrics.PackageMetrics.getInsertionsAndDelitionsInfo(it.unisa.sesa.repominer.db.entities.SourceContainer)}{{\bf  getInsertionsAndDelitionsInfo}\\}
\texttt{public java.lang.Double\lbrack \rbrack \ {\bf  getInsertionsAndDelitionsInfo}(\texttt{it.unisa.sesa.repominer.db.entities.SourceContainer} {\bf  pSourceContainer})
\label{it.unisa.sesa.repominer.metrics.PackageMetrics.getInsertionsAndDelitionsInfo(it.unisa.sesa.repominer.db.entities.SourceContainer)}}%end signature
\begin{itemize}
\item{
{\bf  Description}

This method calculate the Lines metric. The Lines metric represent the total, mean an maximum number of insertion or deletion
}
\item{
{\bf  Parameters}
  \begin{itemize}
   \item{
\texttt{pSourceContainer} -- }
  \end{itemize}
}%end item
\item{{\bf  Returns} -- 
Lines metric value 
}%end item
\end{itemize}
}%end item
\item{ 
\index{getMaxInsertionsOrDeletionsNumber(SourceContainer)}
\hypertarget{it.unisa.sesa.repominer.metrics.PackageMetrics.getMaxInsertionsOrDeletionsNumber(it.unisa.sesa.repominer.db.entities.SourceContainer)}{{\bf  getMaxInsertionsOrDeletionsNumber}\\}
\texttt{public java.lang.Integer\ {\bf  getMaxInsertionsOrDeletionsNumber}(\texttt{it.unisa.sesa.repominer.db.entities.SourceContainer} {\bf  pSourceContainer})
\label{it.unisa.sesa.repominer.metrics.PackageMetrics.getMaxInsertionsOrDeletionsNumber(it.unisa.sesa.repominer.db.entities.SourceContainer)}}%end signature
\begin{itemize}
\item{
{\bf  Description}

This method calculate the second value of Lines metric, namely max number of insertion or deletion
}
\item{
{\bf  Parameters}
  \begin{itemize}
   \item{
\texttt{pSourceContainer} -- }
  \end{itemize}
}%end item
\item{{\bf  Returns} -- 
Second value for Lines metric 
}%end item
\end{itemize}
}%end item
\item{ 
\index{getMeanDimensionOfModifiedFiles(SourceContainer)}
\hypertarget{it.unisa.sesa.repominer.metrics.PackageMetrics.getMeanDimensionOfModifiedFiles(it.unisa.sesa.repominer.db.entities.SourceContainer)}{{\bf  getMeanDimensionOfModifiedFiles}\\}
\texttt{public double\ {\bf  getMeanDimensionOfModifiedFiles}(\texttt{it.unisa.sesa.repominer.db.entities.SourceContainer} {\bf  pSourceContainer})
\label{it.unisa.sesa.repominer.metrics.PackageMetrics.getMeanDimensionOfModifiedFiles(it.unisa.sesa.repominer.db.entities.SourceContainer)}}%end signature
\begin{itemize}
\item{
{\bf  Description}

This method calculate ChangeSetSize metric. The ChangeSetSize metric represent mean dimension of modified files in a package
}
\item{
{\bf  Parameters}
  \begin{itemize}
   \item{
\texttt{pSourceContainer} -- }
  \end{itemize}
}%end item
\item{{\bf  Returns} -- 
ChangeSetSize metric value 
}%end item
\end{itemize}
}%end item
\item{ 
\index{getMeanInsertionsOrDeletionsNumber(SourceContainer)}
\hypertarget{it.unisa.sesa.repominer.metrics.PackageMetrics.getMeanInsertionsOrDeletionsNumber(it.unisa.sesa.repominer.db.entities.SourceContainer)}{{\bf  getMeanInsertionsOrDeletionsNumber}\\}
\texttt{public java.lang.Double\ {\bf  getMeanInsertionsOrDeletionsNumber}(\texttt{it.unisa.sesa.repominer.db.entities.SourceContainer} {\bf  pSourceContainer})
\label{it.unisa.sesa.repominer.metrics.PackageMetrics.getMeanInsertionsOrDeletionsNumber(it.unisa.sesa.repominer.db.entities.SourceContainer)}}%end signature
\begin{itemize}
\item{
{\bf  Description}

This method calculate the third value of Lines metric, namely mean number of insertion or deletion
}
\item{
{\bf  Parameters}
  \begin{itemize}
   \item{
\texttt{pSourceContainer} -- }
  \end{itemize}
}%end item
\item{{\bf  Returns} -- 
Third value for Lines metric 
}%end item
\end{itemize}
}%end item
\item{ 
\index{getMeanNumberOfChange(SourceContainer)}
\hypertarget{it.unisa.sesa.repominer.metrics.PackageMetrics.getMeanNumberOfChange(it.unisa.sesa.repominer.db.entities.SourceContainer)}{{\bf  getMeanNumberOfChange}\\}
\texttt{public double\ {\bf  getMeanNumberOfChange}(\texttt{it.unisa.sesa.repominer.db.entities.SourceContainer} {\bf  pSourceContainer})
\label{it.unisa.sesa.repominer.metrics.PackageMetrics.getMeanNumberOfChange(it.unisa.sesa.repominer.db.entities.SourceContainer)}}%end signature
\begin{itemize}
\item{
{\bf  Description}

This method calculate mean\_NCHANGE metric. The mean\_NCHANGE metric represent mean number of file changes in a package
}
\item{
{\bf  Parameters}
  \begin{itemize}
   \item{
\texttt{pSourceContainer} -- }
  \end{itemize}
}%end item
\item{{\bf  Returns} -- 
mean\_NCHANGE metric value 
}%end item
\end{itemize}
}%end item
\item{ 
\index{getMeanNumberOfChangeForBugFix(SourceContainer)}
\hypertarget{it.unisa.sesa.repominer.metrics.PackageMetrics.getMeanNumberOfChangeForBugFix(it.unisa.sesa.repominer.db.entities.SourceContainer)}{{\bf  getMeanNumberOfChangeForBugFix}\\}
\texttt{public double\ {\bf  getMeanNumberOfChangeForBugFix}(\texttt{it.unisa.sesa.repominer.db.entities.SourceContainer} {\bf  pSourceContainer})
\label{it.unisa.sesa.repominer.metrics.PackageMetrics.getMeanNumberOfChangeForBugFix(it.unisa.sesa.repominer.db.entities.SourceContainer)}}%end signature
\begin{itemize}
\item{
{\bf  Description}

This method calculate mean\_NFIX metric. The mean\_NFIX metric represent mean number of file changes in a package caused by bug fixes
}
\item{
{\bf  Parameters}
  \begin{itemize}
   \item{
\texttt{pSourceContainer} -- }
  \end{itemize}
}%end item
\item{{\bf  Returns} -- 
mean\_NFIX metric value 
}%end item
\end{itemize}
}%end item
\item{ 
\index{getMeanNumberOfChangeForRefactoring(SourceContainer)}
\hypertarget{it.unisa.sesa.repominer.metrics.PackageMetrics.getMeanNumberOfChangeForRefactoring(it.unisa.sesa.repominer.db.entities.SourceContainer)}{{\bf  getMeanNumberOfChangeForRefactoring}\\}
\texttt{public double\ {\bf  getMeanNumberOfChangeForRefactoring}(\texttt{it.unisa.sesa.repominer.db.entities.SourceContainer} {\bf  pSourceContainer})
\label{it.unisa.sesa.repominer.metrics.PackageMetrics.getMeanNumberOfChangeForRefactoring(it.unisa.sesa.repominer.db.entities.SourceContainer)}}%end signature
\begin{itemize}
\item{
{\bf  Description}

This method calculate mean\_NREF metric. The mean\_NREF metric represent mean number of file changes in a package caused by refactoring operation
}
\item{
{\bf  Parameters}
  \begin{itemize}
   \item{
\texttt{pSourceContainer} -- }
  \end{itemize}
}%end item
\item{{\bf  Returns} -- 
mean\_NREF metric value 
}%end item
\end{itemize}
}%end item
\item{ 
\index{getNumberOfAuthor(SourceContainer)}
\hypertarget{it.unisa.sesa.repominer.metrics.PackageMetrics.getNumberOfAuthor(it.unisa.sesa.repominer.db.entities.SourceContainer)}{{\bf  getNumberOfAuthor}\\}
\texttt{public int\ {\bf  getNumberOfAuthor}(\texttt{it.unisa.sesa.repominer.db.entities.SourceContainer} {\bf  pSourceContainer})
\label{it.unisa.sesa.repominer.metrics.PackageMetrics.getNumberOfAuthor(it.unisa.sesa.repominer.db.entities.SourceContainer)}}%end signature
\begin{itemize}
\item{
{\bf  Description}

This method calculate NR metric. The NR metric represent system number of revision
}
\item{
{\bf  Parameters}
  \begin{itemize}
   \item{
\texttt{pSourceContainer} -- }
  \end{itemize}
}%end item
\item{{\bf  Returns} -- 
NR metric value 
}%end item
\end{itemize}
}%end item
\item{ 
\index{getTotalInsertionsOrDeletionsNumber(SourceContainer)}
\hypertarget{it.unisa.sesa.repominer.metrics.PackageMetrics.getTotalInsertionsOrDeletionsNumber(it.unisa.sesa.repominer.db.entities.SourceContainer)}{{\bf  getTotalInsertionsOrDeletionsNumber}\\}
\texttt{public java.lang.Integer\ {\bf  getTotalInsertionsOrDeletionsNumber}(\texttt{it.unisa.sesa.repominer.db.entities.SourceContainer} {\bf  pSourceContainer})
\label{it.unisa.sesa.repominer.metrics.PackageMetrics.getTotalInsertionsOrDeletionsNumber(it.unisa.sesa.repominer.db.entities.SourceContainer)}}%end signature
\begin{itemize}
\item{
{\bf  Description}

This method calculate the first value of Lines metric, namely sum of insertion or deletion
}
\item{
{\bf  Parameters}
  \begin{itemize}
   \item{
\texttt{pSourceContainer} -- }
  \end{itemize}
}%end item
\item{{\bf  Returns} -- 
First value for Lines metric 
}%end item
\end{itemize}
}%end item
\end{itemize}
}
}
\subsection{\label{it.unisa.sesa.repominer.metrics.ProjectMetrics}Class ProjectMetrics}{
\hypertarget{it.unisa.sesa.repominer.metrics.ProjectMetrics}{}\vskip .1in 
This class is responsible to calculate all the history metrics relative to a project.\vskip .1in 
\subsubsection{Declaration}{
\small public class ProjectMetrics
\\ {\bf  extends} java.lang.Object
\refdefined{java.lang.Object}}
\subsubsection{Constructor summary}{
\begin{verse}
\hyperlink{it.unisa.sesa.repominer.metrics.ProjectMetrics()}{{\bf ProjectMetrics()}} \\
\end{verse}
}
\subsubsection{Method summary}{
\begin{verse}
\hyperlink{it.unisa.sesa.repominer.metrics.ProjectMetrics.getBCCMMetric(it.unisa.sesa.repominer.db.entities.Project, java.util.Date, java.util.Date)}{{\bf getBCCMMetric(Project, Date, Date)}} This method calculate the value of BCC Metric; it calculates this value only considering changes occurred in time based period between the start and the end date specified in preference panel\\
\hyperlink{it.unisa.sesa.repominer.metrics.ProjectMetrics.getBCCPeriodBased(it.unisa.sesa.repominer.db.entities.Project, int, java.lang.String)}{{\bf getBCCPeriodBased(Project, int, String)}} This method calculate a set of values of BCC Metric; it calculates this values broken the history of changes into equal length periods based on calendar time from the start of the project.\\
\hyperlink{it.unisa.sesa.repominer.metrics.ProjectMetrics.getECCBurstBased(it.unisa.sesa.repominer.db.entities.Project, int, int, boolean)}{{\bf getECCBurstBased(Project, int, int, boolean)}} This method break history into burst based periods and then calculates the value of ECC Model on this periods\\
\hyperlink{it.unisa.sesa.repominer.metrics.ProjectMetrics.getECCModificationBased(it.unisa.sesa.repominer.db.entities.Project, int, boolean)}{{\bf getECCModificationBased(Project, int, boolean)}} This method calculate a set of values of ECC Model using modification limit period based; it calculates this values broken the history of changes into period based on number of modifications specified in preferences panel.\\
\hyperlink{it.unisa.sesa.repominer.metrics.ProjectMetrics.getECCPeriodBased(it.unisa.sesa.repominer.db.entities.Project, int, java.lang.String, java.lang.Boolean)}{{\bf getECCPeriodBased(Project, int, String, Boolean)}} This method calculate a set of values of ECC Model using time based periods method; it calculates this values broken the history of changes into equal length periods based on calendar time from the start of the project.\\
\hyperlink{it.unisa.sesa.repominer.metrics.ProjectMetrics.getNumberOfRevision(it.unisa.sesa.repominer.db.entities.Project)}{{\bf getNumberOfRevision(Project)}} Calculate the NR metric, that represents the system number of revision\\
\end{verse}
}
\subsubsection{Constructors}{
\vskip -2em
\begin{itemize}
\item{ 
\index{ProjectMetrics()}
\hypertarget{it.unisa.sesa.repominer.metrics.ProjectMetrics()}{{\bf  ProjectMetrics}\\}
\texttt{public\ {\bf  ProjectMetrics}()
\label{it.unisa.sesa.repominer.metrics.ProjectMetrics()}}%end signature
}%end item
\end{itemize}
}
\subsubsection{Methods}{
\vskip -2em
\begin{itemize}
\item{ 
\index{getBCCMMetric(Project, Date, Date)}
\hypertarget{it.unisa.sesa.repominer.metrics.ProjectMetrics.getBCCMMetric(it.unisa.sesa.repominer.db.entities.Project, java.util.Date, java.util.Date)}{{\bf  getBCCMMetric}\\}
\texttt{public it.unisa.sesa.repominer.db.entities.ProjectMetric\ {\bf  getBCCMMetric}(\texttt{it.unisa.sesa.repominer.db.entities.Project} {\bf  pProject},
\texttt{java.util.Date} {\bf  pPeriodStart},
\texttt{java.util.Date} {\bf  pPeriodEnd}) throws it.unisa.sesa.repominer.metrics.exception.NoChangesException
\label{it.unisa.sesa.repominer.metrics.ProjectMetrics.getBCCMMetric(it.unisa.sesa.repominer.db.entities.Project, java.util.Date, java.util.Date)}}%end signature
\begin{itemize}
\item{
{\bf  Description}

This method calculate the value of BCC Metric; it calculates this value only considering changes occurred in time based period between the start and the end date specified in preference panel
}
\item{
{\bf  Parameters}
  \begin{itemize}
   \item{
\texttt{pProject} -- project analyzed}
   \item{
\texttt{pPeriodStart} -- start date for analysis}
   \item{
\texttt{pPeriodEnd} -- end date for analysis}
  \end{itemize}
}%end item
\item{{\bf  Returns} -- 
BCC Metric 
}%end item
\end{itemize}
}%end item
\item{ 
\index{getBCCPeriodBased(Project, int, String)}
\hypertarget{it.unisa.sesa.repominer.metrics.ProjectMetrics.getBCCPeriodBased(it.unisa.sesa.repominer.db.entities.Project, int, java.lang.String)}{{\bf  getBCCPeriodBased}\\}
\texttt{public java.util.List\ {\bf  getBCCPeriodBased}(\texttt{it.unisa.sesa.repominer.db.entities.Project} {\bf  pProject},
\texttt{int} {\bf  periodLength},
\texttt{java.lang.String} {\bf  periodType})
\label{it.unisa.sesa.repominer.metrics.ProjectMetrics.getBCCPeriodBased(it.unisa.sesa.repominer.db.entities.Project, int, java.lang.String)}}%end signature
\begin{itemize}
\item{
{\bf  Description}

This method calculate a set of values of BCC Metric; it calculates this values broken the history of changes into equal length periods based on calendar time from the start of the project. The length of a single interval time is specified in preference panel
}
\item{
{\bf  Parameters}
  \begin{itemize}
   \item{
\texttt{pProject} -- project analyzed}
   \item{
\texttt{periodLength} -- length of period for analysis}
   \item{
\texttt{periodType} -- weeks, months or year for period of analysis}
  \end{itemize}
}%end item
\item{{\bf  Returns} -- 
some BCC Model values time period based 
}%end item
\end{itemize}
}%end item
\item{ 
\index{getECCBurstBased(Project, int, int, boolean)}
\hypertarget{it.unisa.sesa.repominer.metrics.ProjectMetrics.getECCBurstBased(it.unisa.sesa.repominer.db.entities.Project, int, int, boolean)}{{\bf  getECCBurstBased}\\}
\texttt{public java.util.List\ {\bf  getECCBurstBased}(\texttt{it.unisa.sesa.repominer.db.entities.Project} {\bf  pProject},
\texttt{int} {\bf  pEps},
\texttt{int} {\bf  pMinPoints},
\texttt{boolean} {\bf  pIsStatic}) throws it.unisa.sesa.repominer.metrics.exception.NoChangesException
\label{it.unisa.sesa.repominer.metrics.ProjectMetrics.getECCBurstBased(it.unisa.sesa.repominer.db.entities.Project, int, int, boolean)}}%end signature
\begin{itemize}
\item{
{\bf  Description}

This method break history into burst based periods and then calculates the value of ECC Model on this periods
}
\item{
{\bf  Parameters}
  \begin{itemize}
   \item{
\texttt{pProject} -- project analyzed}
   \item{
\texttt{pEps} -- eps parameter for dbscan analysis}
   \item{
\texttt{pMinPoints} -- minPoint parameter for dbscan analysis}
   \item{
\texttt{pIsStatic} -- if true calculates ECCM with Normalized Static Entropy; if false calculates ECCM with Adaptive Sizing Entropy}
  \end{itemize}
}%end item
\item{{\bf  Returns} -- 
a list of ECC Model values 
}%end item
\end{itemize}
}%end item
\item{ 
\index{getECCModificationBased(Project, int, boolean)}
\hypertarget{it.unisa.sesa.repominer.metrics.ProjectMetrics.getECCModificationBased(it.unisa.sesa.repominer.db.entities.Project, int, boolean)}{{\bf  getECCModificationBased}\\}
\texttt{public java.util.List\ {\bf  getECCModificationBased}(\texttt{it.unisa.sesa.repominer.db.entities.Project} {\bf  pProject},
\texttt{int} {\bf  pLimit},
\texttt{boolean} {\bf  pIsStatic}) throws it.unisa.sesa.repominer.metrics.exception.NoChangesException
\label{it.unisa.sesa.repominer.metrics.ProjectMetrics.getECCModificationBased(it.unisa.sesa.repominer.db.entities.Project, int, boolean)}}%end signature
\begin{itemize}
\item{
{\bf  Description}

This method calculate a set of values of ECC Model using modification limit period based; it calculates this values broken the history of changes into period based on number of modifications specified in preferences panel. To prevent a period where little development may have occurred from spanning a long time, we impose a limit of 3 months on a period even if the modification limit was no reached
}
\item{
{\bf  Parameters}
  \begin{itemize}
   \item{
\texttt{pProject} -- project analyzed}
   \item{
\texttt{pLimit} -- limit of modifications for breaking periods}
   \item{
\texttt{pIsStatic} -- if true calculated with Normalized Static Entropy; if false calculated with our Adaptive Sizing Entropy}
  \end{itemize}
}%end item
\item{{\bf  Returns} -- 
some ECC Model value calculated with modification limit period method 
}%end item
\end{itemize}
}%end item
\item{ 
\index{getECCPeriodBased(Project, int, String, Boolean)}
\hypertarget{it.unisa.sesa.repominer.metrics.ProjectMetrics.getECCPeriodBased(it.unisa.sesa.repominer.db.entities.Project, int, java.lang.String, java.lang.Boolean)}{{\bf  getECCPeriodBased}\\}
\texttt{public java.util.List\ {\bf  getECCPeriodBased}(\texttt{it.unisa.sesa.repominer.db.entities.Project} {\bf  pProject},
\texttt{int} {\bf  pPeriod},
\texttt{java.lang.String} {\bf  periodType},
\texttt{java.lang.Boolean} {\bf  pIsStatic})
\label{it.unisa.sesa.repominer.metrics.ProjectMetrics.getECCPeriodBased(it.unisa.sesa.repominer.db.entities.Project, int, java.lang.String, java.lang.Boolean)}}%end signature
\begin{itemize}
\item{
{\bf  Description}

This method calculate a set of values of ECC Model using time based periods method; it calculates this values broken the history of changes into equal length periods based on calendar time from the start of the project. The length of a single interval time is specified in preference panel
}
\item{
{\bf  Parameters}
  \begin{itemize}
   \item{
\texttt{pProject} -- project analyzed}
   \item{
\texttt{pPeriod} -- length of period for analysis}
   \item{
\texttt{periodType} -- weeks, months or years for period of analysis}
   \item{
\texttt{pIsStatic} -- if true calculates ECCM with Normalized Static Entropy; if false calculates ECCM with Adaptive Sizing Entropy}
  \end{itemize}
}%end item
\item{{\bf  Returns} -- 
some ECC Model values calculated with time based periods method 
}%end item
\end{itemize}
}%end item
\item{ 
\index{getNumberOfRevision(Project)}
\hypertarget{it.unisa.sesa.repominer.metrics.ProjectMetrics.getNumberOfRevision(it.unisa.sesa.repominer.db.entities.Project)}{{\bf  getNumberOfRevision}\\}
\texttt{public it.unisa.sesa.repominer.db.entities.ProjectMetric\ {\bf  getNumberOfRevision}(\texttt{it.unisa.sesa.repominer.db.entities.Project} {\bf  pProject})
\label{it.unisa.sesa.repominer.metrics.ProjectMetrics.getNumberOfRevision(it.unisa.sesa.repominer.db.entities.Project)}}%end signature
\begin{itemize}
\item{
{\bf  Description}

Calculate the NR metric, that represents the system number of revision
}
\item{
{\bf  Parameters}
  \begin{itemize}
   \item{
\texttt{pProject} -- project analyzed}
  \end{itemize}
}%end item
\item{{\bf  Returns} -- 
the NR metric value 
}%end item
\end{itemize}
}%end item
\end{itemize}
}
}
}
\section{Package it.unisa.sesa.repominer.preferences.exceptions}{
\label{it.unisa.sesa.repominer.preferences.exceptions}\hypertarget{it.unisa.sesa.repominer.preferences.exceptions}{}
\hskip -.05in
\hbox to \hsize{\textit{ Package Contents\hfil Page}}
\vskip .1in
\vskip .1in
\subsection{\label{it.unisa.sesa.repominer.preferences.exceptions.IntegerPreferenceException}Exception IntegerPreferenceException}{
\hypertarget{it.unisa.sesa.repominer.preferences.exceptions.IntegerPreferenceException}{}\vskip .1in 
This exception is thrown when a preference is not an int value\vskip .1in 
\subsubsection{Declaration}{
\small public class IntegerPreferenceException
\\ {\bf  extends} java.lang.Exception
\refdefined{java.lang.Exception}}
\subsubsection{Constructor summary}{
\begin{verse}
\hyperlink{it.unisa.sesa.repominer.preferences.exceptions.IntegerPreferenceException(java.lang.String, int)}{{\bf IntegerPreferenceException(String, int)}} This constructor creates an IntegerPreferenceException initializing some instance variables.\\
\end{verse}
}
\subsubsection{Method summary}{
\begin{verse}
\hyperlink{it.unisa.sesa.repominer.preferences.exceptions.IntegerPreferenceException.getPreferenceName()}{{\bf getPreferenceName()}} This method returns the name of the preference that caused this exception to be thrown.\\
\hyperlink{it.unisa.sesa.repominer.preferences.exceptions.IntegerPreferenceException.getPreferenceValue()}{{\bf getPreferenceValue()}} This method returns the value of the preference that caused this exception to be thrown.\\
\end{verse}
}
\subsubsection{Constructors}{
\vskip -2em
\begin{itemize}
\item{ 
\index{IntegerPreferenceException(String, int)}
\hypertarget{it.unisa.sesa.repominer.preferences.exceptions.IntegerPreferenceException(java.lang.String, int)}{{\bf  IntegerPreferenceException}\\}
\texttt{public\ {\bf  IntegerPreferenceException}(\texttt{java.lang.String} {\bf  pPreferenceName},
\texttt{int} {\bf  pPreferenceValue})
\label{it.unisa.sesa.repominer.preferences.exceptions.IntegerPreferenceException(java.lang.String, int)}}%end signature
\begin{itemize}
\item{
{\bf  Description}

This constructor creates an IntegerPreferenceException initializing some instance variables.
}
\item{
{\bf  Parameters}
  \begin{itemize}
   \item{
\texttt{pPreferenceName} -- The name of the preference that caused this exception to be thrown.}
   \item{
\texttt{pPreferenceValue} -- The value of the preference that caused this exception to be thrown.}
  \end{itemize}
}%end item
\end{itemize}
}%end item
\end{itemize}
}
\subsubsection{Methods}{
\vskip -2em
\begin{itemize}
\item{ 
\index{getPreferenceName()}
\hypertarget{it.unisa.sesa.repominer.preferences.exceptions.IntegerPreferenceException.getPreferenceName()}{{\bf  getPreferenceName}\\}
\texttt{public java.lang.String\ {\bf  getPreferenceName}()
\label{it.unisa.sesa.repominer.preferences.exceptions.IntegerPreferenceException.getPreferenceName()}}%end signature
\begin{itemize}
\item{
{\bf  Description}

This method returns the name of the preference that caused this exception to be thrown.
}
\item{{\bf  Returns} -- 
A String object representing the name of the preference that caused this exception to be thrown. 
}%end item
\end{itemize}
}%end item
\item{ 
\index{getPreferenceValue()}
\hypertarget{it.unisa.sesa.repominer.preferences.exceptions.IntegerPreferenceException.getPreferenceValue()}{{\bf  getPreferenceValue}\\}
\texttt{public int\ {\bf  getPreferenceValue}()
\label{it.unisa.sesa.repominer.preferences.exceptions.IntegerPreferenceException.getPreferenceValue()}}%end signature
\begin{itemize}
\item{
{\bf  Description}

This method returns the value of the preference that caused this exception to be thrown.
}
\item{{\bf  Returns} -- 
An int value representing the value of the preference that caused this exception to be thrown. 
}%end item
\end{itemize}
}%end item
\end{itemize}
}
\subsubsection{Members inherited from class Throwable }{
\texttt{java.lang.Throwable} {\small 
\refdefined{java.lang.Throwable}}
{\small 

\vskip -2em
\begin{itemize}
\item{\vskip -1.5ex 
\texttt{public final synchronized void {\bf  addSuppressed}(\texttt{Throwable} {\bf  arg0})
}%end signature
}%end item
\item{\vskip -1.5ex 
\texttt{public synchronized Throwable {\bf  fillInStackTrace}()
}%end signature
}%end item
\item{\vskip -1.5ex 
\texttt{public synchronized Throwable {\bf  getCause}()
}%end signature
}%end item
\item{\vskip -1.5ex 
\texttt{public String {\bf  getLocalizedMessage}()
}%end signature
}%end item
\item{\vskip -1.5ex 
\texttt{public String {\bf  getMessage}()
}%end signature
}%end item
\item{\vskip -1.5ex 
\texttt{public StackTraceElement {\bf  getStackTrace}()
}%end signature
}%end item
\item{\vskip -1.5ex 
\texttt{public final synchronized Throwable {\bf  getSuppressed}()
}%end signature
}%end item
\item{\vskip -1.5ex 
\texttt{public synchronized Throwable {\bf  initCause}(\texttt{Throwable} {\bf  arg0})
}%end signature
}%end item
\item{\vskip -1.5ex 
\texttt{public void {\bf  printStackTrace}()
}%end signature
}%end item
\item{\vskip -1.5ex 
\texttt{public void {\bf  printStackTrace}(\texttt{java.io.PrintStream} {\bf  arg0})
}%end signature
}%end item
\item{\vskip -1.5ex 
\texttt{public void {\bf  printStackTrace}(\texttt{java.io.PrintWriter} {\bf  arg0})
}%end signature
}%end item
\item{\vskip -1.5ex 
\texttt{public void {\bf  setStackTrace}(\texttt{StackTraceElement\lbrack \rbrack } {\bf  arg0})
}%end signature
}%end item
\item{\vskip -1.5ex 
\texttt{public String {\bf  toString}()
}%end signature
}%end item
\end{itemize}
}
}
\subsection{\label{it.unisa.sesa.repominer.preferences.exceptions.PeriodLengthTooLong}Exception PeriodLengthTooLong}{
\hypertarget{it.unisa.sesa.repominer.preferences.exceptions.PeriodLengthTooLong}{}\vskip .1in 
\subsubsection{Declaration}{
\small public class PeriodLengthTooLong
\\ {\bf  extends} java.lang.Exception
\refdefined{java.lang.Exception}}
\subsubsection{Constructor summary}{
\begin{verse}
\hyperlink{it.unisa.sesa.repominer.preferences.exceptions.PeriodLengthTooLong(int, int)}{{\bf PeriodLengthTooLong(int, int)}} This constructor builds an object of type PeriodLengthTooLong by setting the message of the exception to a string representation indicating the limit exceeded and the value that exceeded that limit.\\
\end{verse}
}
\subsubsection{Constructors}{
\vskip -2em
\begin{itemize}
\item{ 
\index{PeriodLengthTooLong(int, int)}
\hypertarget{it.unisa.sesa.repominer.preferences.exceptions.PeriodLengthTooLong(int, int)}{{\bf  PeriodLengthTooLong}\\}
\texttt{public\ {\bf  PeriodLengthTooLong}(\texttt{int} {\bf  pLimit},
\texttt{int} {\bf  pLength})
\label{it.unisa.sesa.repominer.preferences.exceptions.PeriodLengthTooLong(int, int)}}%end signature
\begin{itemize}
\item{
{\bf  Description}

This constructor builds an object of type PeriodLengthTooLong by setting the message of the exception to a string representation indicating the limit exceeded and the value that exceeded that limit.
}
\item{
{\bf  Parameters}
  \begin{itemize}
   \item{
\texttt{pLimit} -- }
   \item{
\texttt{pLength} -- }
  \end{itemize}
}%end item
\end{itemize}
}%end item
\end{itemize}
}
\subsubsection{Members inherited from class Throwable }{
\texttt{java.lang.Throwable} {\small 
\refdefined{java.lang.Throwable}}
{\small 

\vskip -2em
\begin{itemize}
\item{\vskip -1.5ex 
\texttt{public final synchronized void {\bf  addSuppressed}(\texttt{Throwable} {\bf  arg0})
}%end signature
}%end item
\item{\vskip -1.5ex 
\texttt{public synchronized Throwable {\bf  fillInStackTrace}()
}%end signature
}%end item
\item{\vskip -1.5ex 
\texttt{public synchronized Throwable {\bf  getCause}()
}%end signature
}%end item
\item{\vskip -1.5ex 
\texttt{public String {\bf  getLocalizedMessage}()
}%end signature
}%end item
\item{\vskip -1.5ex 
\texttt{public String {\bf  getMessage}()
}%end signature
}%end item
\item{\vskip -1.5ex 
\texttt{public StackTraceElement {\bf  getStackTrace}()
}%end signature
}%end item
\item{\vskip -1.5ex 
\texttt{public final synchronized Throwable {\bf  getSuppressed}()
}%end signature
}%end item
\item{\vskip -1.5ex 
\texttt{public synchronized Throwable {\bf  initCause}(\texttt{Throwable} {\bf  arg0})
}%end signature
}%end item
\item{\vskip -1.5ex 
\texttt{public void {\bf  printStackTrace}()
}%end signature
}%end item
\item{\vskip -1.5ex 
\texttt{public void {\bf  printStackTrace}(\texttt{java.io.PrintStream} {\bf  arg0})
}%end signature
}%end item
\item{\vskip -1.5ex 
\texttt{public void {\bf  printStackTrace}(\texttt{java.io.PrintWriter} {\bf  arg0})
}%end signature
}%end item
\item{\vskip -1.5ex 
\texttt{public void {\bf  setStackTrace}(\texttt{StackTraceElement\lbrack \rbrack } {\bf  arg0})
}%end signature
}%end item
\item{\vskip -1.5ex 
\texttt{public String {\bf  toString}()
}%end signature
}%end item
\end{itemize}
}
}
}
\section{Package it.unisa.sesa.repominer.preferences}{
\label{it.unisa.sesa.repominer.preferences}\hypertarget{it.unisa.sesa.repominer.preferences}{}
\hskip -.05in
\hbox to \hsize{\textit{ Package Contents\hfil Page}}
\vskip .13in
\hbox{{\bf  Classes}}
\entityintro{PreferenceConstants}{it.unisa.sesa.repominer.preferences.PreferenceConstants}{Constant definitions for plug-in preferences}
\entityintro{PreferenceInitializer}{it.unisa.sesa.repominer.preferences.PreferenceInitializer}{Class used to initialize default preference values.}
\entityintro{PreferencePage}{it.unisa.sesa.repominer.preferences.PreferencePage}{This class represents a preference page that is contributed to the Preferences dialog.}
\entityintro{Preferences}{it.unisa.sesa.repominer.preferences.Preferences}{}
\vskip .1in
\vskip .1in
\subsection{\label{it.unisa.sesa.repominer.preferences.PreferenceConstants}Class PreferenceConstants}{
\hypertarget{it.unisa.sesa.repominer.preferences.PreferenceConstants}{}\vskip .1in 
Constant definitions for plug-in preferences\vskip .1in 
\subsubsection{Declaration}{
\small public class PreferenceConstants
\\ {\bf  extends} java.lang.Object
\refdefined{java.lang.Object}}
\subsubsection{Field summary}{
\begin{verse}
\hyperlink{it.unisa.sesa.repominer.preferences.PreferenceConstants.BURST_EPS}{{\bf BURST\_EPS}} \\
\hyperlink{it.unisa.sesa.repominer.preferences.PreferenceConstants.BURST_MINPOINTS}{{\bf BURST\_MINPOINTS}} \\
\hyperlink{it.unisa.sesa.repominer.preferences.PreferenceConstants.ECCM_BURST_VALUE}{{\bf ECCM\_BURST\_VALUE}} \\
\hyperlink{it.unisa.sesa.repominer.preferences.PreferenceConstants.ECCM_MODALITY}{{\bf ECCM\_MODALITY}} \\
\hyperlink{it.unisa.sesa.repominer.preferences.PreferenceConstants.ECCM_MODIFICATION_LIMIT}{{\bf ECCM\_MODIFICATION\_LIMIT}} \\
\hyperlink{it.unisa.sesa.repominer.preferences.PreferenceConstants.ECCM_MODIFICATION_VALUE}{{\bf ECCM\_MODIFICATION\_VALUE}} \\
\hyperlink{it.unisa.sesa.repominer.preferences.PreferenceConstants.ECCM_TIME_VALUE}{{\bf ECCM\_TIME\_VALUE}} \\
\hyperlink{it.unisa.sesa.repominer.preferences.PreferenceConstants.P_DBHOST}{{\bf P\_DBHOST}} \\
\hyperlink{it.unisa.sesa.repominer.preferences.PreferenceConstants.P_DBNAME}{{\bf P\_DBNAME}} \\
\hyperlink{it.unisa.sesa.repominer.preferences.PreferenceConstants.P_DBPASS}{{\bf P\_DBPASS}} \\
\hyperlink{it.unisa.sesa.repominer.preferences.PreferenceConstants.P_DBPORT}{{\bf P\_DBPORT}} \\
\hyperlink{it.unisa.sesa.repominer.preferences.PreferenceConstants.P_DBUSER}{{\bf P\_DBUSER}} \\
\hyperlink{it.unisa.sesa.repominer.preferences.PreferenceConstants.PERIOD_END}{{\bf PERIOD\_END}} \\
\hyperlink{it.unisa.sesa.repominer.preferences.PreferenceConstants.PERIOD_LENGTH}{{\bf PERIOD\_LENGTH}} \\
\hyperlink{it.unisa.sesa.repominer.preferences.PreferenceConstants.PERIOD_START}{{\bf PERIOD\_START}} \\
\hyperlink{it.unisa.sesa.repominer.preferences.PreferenceConstants.PERIOD_TYPE}{{\bf PERIOD\_TYPE}} \\
\hyperlink{it.unisa.sesa.repominer.preferences.PreferenceConstants.PERIOD_TYPE_MONTH}{{\bf PERIOD\_TYPE\_MONTH}} \\
\hyperlink{it.unisa.sesa.repominer.preferences.PreferenceConstants.PERIOD_TYPE_WEEK}{{\bf PERIOD\_TYPE\_WEEK}} \\
\hyperlink{it.unisa.sesa.repominer.preferences.PreferenceConstants.PERIOD_TYPE_YEAR}{{\bf PERIOD\_TYPE\_YEAR}} \\
\end{verse}
}
\subsubsection{Constructor summary}{
\begin{verse}
\hyperlink{it.unisa.sesa.repominer.preferences.PreferenceConstants()}{{\bf PreferenceConstants()}} \\
\end{verse}
}
\subsubsection{Fields}{
\begin{itemize}
\item{
\index{P\_DBPORT}
\label{it.unisa.sesa.repominer.preferences.PreferenceConstants.P_DBPORT}\hypertarget{it.unisa.sesa.repominer.preferences.PreferenceConstants.P_DBPORT}{public static final java.lang.String {\bf  P\_DBPORT}}
}
\item{
\index{P\_DBHOST}
\label{it.unisa.sesa.repominer.preferences.PreferenceConstants.P_DBHOST}\hypertarget{it.unisa.sesa.repominer.preferences.PreferenceConstants.P_DBHOST}{public static final java.lang.String {\bf  P\_DBHOST}}
}
\item{
\index{P\_DBNAME}
\label{it.unisa.sesa.repominer.preferences.PreferenceConstants.P_DBNAME}\hypertarget{it.unisa.sesa.repominer.preferences.PreferenceConstants.P_DBNAME}{public static final java.lang.String {\bf  P\_DBNAME}}
}
\item{
\index{P\_DBUSER}
\label{it.unisa.sesa.repominer.preferences.PreferenceConstants.P_DBUSER}\hypertarget{it.unisa.sesa.repominer.preferences.PreferenceConstants.P_DBUSER}{public static final java.lang.String {\bf  P\_DBUSER}}
}
\item{
\index{P\_DBPASS}
\label{it.unisa.sesa.repominer.preferences.PreferenceConstants.P_DBPASS}\hypertarget{it.unisa.sesa.repominer.preferences.PreferenceConstants.P_DBPASS}{public static final java.lang.String {\bf  P\_DBPASS}}
}
\item{
\index{PERIOD\_START}
\label{it.unisa.sesa.repominer.preferences.PreferenceConstants.PERIOD_START}\hypertarget{it.unisa.sesa.repominer.preferences.PreferenceConstants.PERIOD_START}{public static final java.lang.String {\bf  PERIOD\_START}}
}
\item{
\index{PERIOD\_END}
\label{it.unisa.sesa.repominer.preferences.PreferenceConstants.PERIOD_END}\hypertarget{it.unisa.sesa.repominer.preferences.PreferenceConstants.PERIOD_END}{public static final java.lang.String {\bf  PERIOD\_END}}
}
\item{
\index{PERIOD\_LENGTH}
\label{it.unisa.sesa.repominer.preferences.PreferenceConstants.PERIOD_LENGTH}\hypertarget{it.unisa.sesa.repominer.preferences.PreferenceConstants.PERIOD_LENGTH}{public static final java.lang.String {\bf  PERIOD\_LENGTH}}
}
\item{
\index{PERIOD\_TYPE}
\label{it.unisa.sesa.repominer.preferences.PreferenceConstants.PERIOD_TYPE}\hypertarget{it.unisa.sesa.repominer.preferences.PreferenceConstants.PERIOD_TYPE}{public static final java.lang.String {\bf  PERIOD\_TYPE}}
}
\item{
\index{PERIOD\_TYPE\_WEEK}
\label{it.unisa.sesa.repominer.preferences.PreferenceConstants.PERIOD_TYPE_WEEK}\hypertarget{it.unisa.sesa.repominer.preferences.PreferenceConstants.PERIOD_TYPE_WEEK}{public static final java.lang.String {\bf  PERIOD\_TYPE\_WEEK}}
}
\item{
\index{PERIOD\_TYPE\_MONTH}
\label{it.unisa.sesa.repominer.preferences.PreferenceConstants.PERIOD_TYPE_MONTH}\hypertarget{it.unisa.sesa.repominer.preferences.PreferenceConstants.PERIOD_TYPE_MONTH}{public static final java.lang.String {\bf  PERIOD\_TYPE\_MONTH}}
}
\item{
\index{PERIOD\_TYPE\_YEAR}
\label{it.unisa.sesa.repominer.preferences.PreferenceConstants.PERIOD_TYPE_YEAR}\hypertarget{it.unisa.sesa.repominer.preferences.PreferenceConstants.PERIOD_TYPE_YEAR}{public static final java.lang.String {\bf  PERIOD\_TYPE\_YEAR}}
}
\item{
\index{ECCM\_MODALITY}
\label{it.unisa.sesa.repominer.preferences.PreferenceConstants.ECCM_MODALITY}\hypertarget{it.unisa.sesa.repominer.preferences.PreferenceConstants.ECCM_MODALITY}{public static final java.lang.String {\bf  ECCM\_MODALITY}}
}
\item{
\index{ECCM\_MODIFICATION\_LIMIT}
\label{it.unisa.sesa.repominer.preferences.PreferenceConstants.ECCM_MODIFICATION_LIMIT}\hypertarget{it.unisa.sesa.repominer.preferences.PreferenceConstants.ECCM_MODIFICATION_LIMIT}{public static final java.lang.String {\bf  ECCM\_MODIFICATION\_LIMIT}}
}
\item{
\index{ECCM\_MODIFICATION\_VALUE}
\label{it.unisa.sesa.repominer.preferences.PreferenceConstants.ECCM_MODIFICATION_VALUE}\hypertarget{it.unisa.sesa.repominer.preferences.PreferenceConstants.ECCM_MODIFICATION_VALUE}{public static final java.lang.String {\bf  ECCM\_MODIFICATION\_VALUE}}
}
\item{
\index{ECCM\_TIME\_VALUE}
\label{it.unisa.sesa.repominer.preferences.PreferenceConstants.ECCM_TIME_VALUE}\hypertarget{it.unisa.sesa.repominer.preferences.PreferenceConstants.ECCM_TIME_VALUE}{public static final java.lang.String {\bf  ECCM\_TIME\_VALUE}}
}
\item{
\index{ECCM\_BURST\_VALUE}
\label{it.unisa.sesa.repominer.preferences.PreferenceConstants.ECCM_BURST_VALUE}\hypertarget{it.unisa.sesa.repominer.preferences.PreferenceConstants.ECCM_BURST_VALUE}{public static final java.lang.String {\bf  ECCM\_BURST\_VALUE}}
}
\item{
\index{BURST\_EPS}
\label{it.unisa.sesa.repominer.preferences.PreferenceConstants.BURST_EPS}\hypertarget{it.unisa.sesa.repominer.preferences.PreferenceConstants.BURST_EPS}{public static final java.lang.String {\bf  BURST\_EPS}}
}
\item{
\index{BURST\_MINPOINTS}
\label{it.unisa.sesa.repominer.preferences.PreferenceConstants.BURST_MINPOINTS}\hypertarget{it.unisa.sesa.repominer.preferences.PreferenceConstants.BURST_MINPOINTS}{public static final java.lang.String {\bf  BURST\_MINPOINTS}}
}
\end{itemize}
}
\subsubsection{Constructors}{
\vskip -2em
\begin{itemize}
\item{ 
\index{PreferenceConstants()}
\hypertarget{it.unisa.sesa.repominer.preferences.PreferenceConstants()}{{\bf  PreferenceConstants}\\}
\texttt{public\ {\bf  PreferenceConstants}()
\label{it.unisa.sesa.repominer.preferences.PreferenceConstants()}}%end signature
}%end item
\end{itemize}
}
}
\subsection{\label{it.unisa.sesa.repominer.preferences.PreferenceInitializer}Class PreferenceInitializer}{
\hypertarget{it.unisa.sesa.repominer.preferences.PreferenceInitializer}{}\vskip .1in 
Class used to initialize default preference values.\vskip .1in 
\subsubsection{Declaration}{
\small public class PreferenceInitializer
\\ {\bf  extends} AbstractPreferenceInitializer
\refdefined{.AbstractPreferenceInitializer}}
\subsubsection{Constructor summary}{
\begin{verse}
\hyperlink{it.unisa.sesa.repominer.preferences.PreferenceInitializer()}{{\bf PreferenceInitializer()}} \\
\end{verse}
}
\subsubsection{Method summary}{
\begin{verse}
\hyperlink{it.unisa.sesa.repominer.preferences.PreferenceInitializer.initializeDefaultPreferences()}{{\bf initializeDefaultPreferences()}} \\
\end{verse}
}
\subsubsection{Constructors}{
\vskip -2em
\begin{itemize}
\item{ 
\index{PreferenceInitializer()}
\hypertarget{it.unisa.sesa.repominer.preferences.PreferenceInitializer()}{{\bf  PreferenceInitializer}\\}
\texttt{public\ {\bf  PreferenceInitializer}()
\label{it.unisa.sesa.repominer.preferences.PreferenceInitializer()}}%end signature
}%end item
\end{itemize}
}
\subsubsection{Methods}{
\vskip -2em
\begin{itemize}
\item{ 
\index{initializeDefaultPreferences()}
\hypertarget{it.unisa.sesa.repominer.preferences.PreferenceInitializer.initializeDefaultPreferences()}{{\bf  initializeDefaultPreferences}\\}
\texttt{public void\ {\bf  initializeDefaultPreferences}()
\label{it.unisa.sesa.repominer.preferences.PreferenceInitializer.initializeDefaultPreferences()}}%end signature
}%end item
\end{itemize}
}
}
\subsection{\label{it.unisa.sesa.repominer.preferences.PreferencePage}Class PreferencePage}{
\hypertarget{it.unisa.sesa.repominer.preferences.PreferencePage}{}\vskip .1in 
This class represents a preference page that is contributed to the Preferences dialog. By subclassing FieldEditorPreferencePage, we can use the field support built into JFace that allows us to create a page that is small and knows how to save, restore and apply itself.

This page is used to modify preferences only. They are stored in the preference store that belongs to the main plug-in class. That way, preferences can be accessed directly via the preference store.\vskip .1in 
\subsubsection{Declaration}{
\small public class PreferencePage
\\ {\bf  extends} FieldEditorPreferencePage
\refdefined{.FieldEditorPreferencePage}}
\subsubsection{Constructor summary}{
\begin{verse}
\hyperlink{it.unisa.sesa.repominer.preferences.PreferencePage()}{{\bf PreferencePage()}} \\
\end{verse}
}
\subsubsection{Method summary}{
\begin{verse}
\hyperlink{it.unisa.sesa.repominer.preferences.PreferencePage.createFieldEditors()}{{\bf createFieldEditors()}} Creates the field editors.\\
\hyperlink{it.unisa.sesa.repominer.preferences.PreferencePage.init(IWorkbench)}{{\bf init(IWorkbench)}} \\
\end{verse}
}
\subsubsection{Constructors}{
\vskip -2em
\begin{itemize}
\item{ 
\index{PreferencePage()}
\hypertarget{it.unisa.sesa.repominer.preferences.PreferencePage()}{{\bf  PreferencePage}\\}
\texttt{public\ {\bf  PreferencePage}()
\label{it.unisa.sesa.repominer.preferences.PreferencePage()}}%end signature
}%end item
\end{itemize}
}
\subsubsection{Methods}{
\vskip -2em
\begin{itemize}
\item{ 
\index{createFieldEditors()}
\hypertarget{it.unisa.sesa.repominer.preferences.PreferencePage.createFieldEditors()}{{\bf  createFieldEditors}\\}
\texttt{public void\ {\bf  createFieldEditors}()
\label{it.unisa.sesa.repominer.preferences.PreferencePage.createFieldEditors()}}%end signature
\begin{itemize}
\item{
{\bf  Description}

Creates the field editors. Field editors are abstractions of the common GUI blocks needed to manipulate various types of preferences. Each field editor knows how to save and restore itself.
}
\end{itemize}
}%end item
\item{ 
\index{init(IWorkbench)}
\hypertarget{it.unisa.sesa.repominer.preferences.PreferencePage.init(IWorkbench)}{{\bf  init}\\}
\texttt{public void\ {\bf  init}(\texttt{IWorkbench} {\bf  workbench})
\label{it.unisa.sesa.repominer.preferences.PreferencePage.init(IWorkbench)}}%end signature
}%end item
\end{itemize}
}
}
\subsection{\label{it.unisa.sesa.repominer.preferences.Preferences}Class Preferences}{
\hypertarget{it.unisa.sesa.repominer.preferences.Preferences}{}\vskip .1in 
\subsubsection{Declaration}{
\small public class Preferences
\\ {\bf  extends} java.lang.Object
\refdefined{java.lang.Object}}
\subsubsection{Constructor summary}{
\begin{verse}
\hyperlink{it.unisa.sesa.repominer.preferences.Preferences()}{{\bf Preferences()}} \\
\end{verse}
}
\subsubsection{Method summary}{
\begin{verse}
\hyperlink{it.unisa.sesa.repominer.preferences.Preferences.getDatabaseHost()}{{\bf getDatabaseHost()}} This method fetches the value of database host from the Eclipse preference store.\\
\hyperlink{it.unisa.sesa.repominer.preferences.Preferences.getDatabaseName()}{{\bf getDatabaseName()}} This method fetches the value of database name from the Eclipse preference store.\\
\hyperlink{it.unisa.sesa.repominer.preferences.Preferences.getDatabasePassword()}{{\bf getDatabasePassword()}} This method fetches the value of database user's password from the Eclipse preference store.\\
\hyperlink{it.unisa.sesa.repominer.preferences.Preferences.getDatabasePort()}{{\bf getDatabasePort()}} This method fetches the value of database port from the Eclipse preference store.\\
\hyperlink{it.unisa.sesa.repominer.preferences.Preferences.getDatabaseUser()}{{\bf getDatabaseUser()}} This method fetches the value of database username from the Eclipse preference store.\\
\hyperlink{it.unisa.sesa.repominer.preferences.Preferences.getECCMModality()}{{\bf getECCMModality()}} This method fetches the value of the Extended Code Change Model metric modality from the Eclipse preference store.\\
\hyperlink{it.unisa.sesa.repominer.preferences.Preferences.getECCMModificationLimit()}{{\bf getECCMModificationLimit()}} This method fetches the value of the modification limit value for the Extended Code Change Model metric.\\
\hyperlink{it.unisa.sesa.repominer.preferences.Preferences.getEpsParameter()}{{\bf getEpsParameter()}} This method fetches the value of radius parameter for DBSCAN algorithm from Eclipse preference store.\\
\hyperlink{it.unisa.sesa.repominer.preferences.Preferences.getMinPointsParameter()}{{\bf getMinPointsParameter()}} This method fetches the value of the minimum size per cluster relative to the DBSCAN algorithm from Eclipse preference store.\\
\hyperlink{it.unisa.sesa.repominer.preferences.Preferences.getPeriodEndingDate()}{{\bf getPeriodEndingDate()}} This method fetches the value for the period ending date for the Basic Code Change Model metric.\\
\hyperlink{it.unisa.sesa.repominer.preferences.Preferences.getPeriodLength()}{{\bf getPeriodLength()}} This method fetches the value of the period length from the Eclipse preference store.\\
\hyperlink{it.unisa.sesa.repominer.preferences.Preferences.getPeriodStartingDate()}{{\bf getPeriodStartingDate()}} This method fetches the value for the period starting date for the Basic Code Change Model metric.\\
\hyperlink{it.unisa.sesa.repominer.preferences.Preferences.getPeriodType()}{{\bf getPeriodType()}} This method fetches the value of the period type from Eclipse preference store.\\
\end{verse}
}
\subsubsection{Constructors}{
\vskip -2em
\begin{itemize}
\item{ 
\index{Preferences()}
\hypertarget{it.unisa.sesa.repominer.preferences.Preferences()}{{\bf  Preferences}\\}
\texttt{public\ {\bf  Preferences}()
\label{it.unisa.sesa.repominer.preferences.Preferences()}}%end signature
}%end item
\end{itemize}
}
\subsubsection{Methods}{
\vskip -2em
\begin{itemize}
\item{ 
\index{getDatabaseHost()}
\hypertarget{it.unisa.sesa.repominer.preferences.Preferences.getDatabaseHost()}{{\bf  getDatabaseHost}\\}
\texttt{public static java.lang.String\ {\bf  getDatabaseHost}()
\label{it.unisa.sesa.repominer.preferences.Preferences.getDatabaseHost()}}%end signature
\begin{itemize}
\item{
{\bf  Description}

This method fetches the value of database host from the Eclipse preference store.
}
\item{{\bf  Returns} -- 
A String object representing the database host to which the application connects. 
}%end item
\end{itemize}
}%end item
\item{ 
\index{getDatabaseName()}
\hypertarget{it.unisa.sesa.repominer.preferences.Preferences.getDatabaseName()}{{\bf  getDatabaseName}\\}
\texttt{public static java.lang.String\ {\bf  getDatabaseName}()
\label{it.unisa.sesa.repominer.preferences.Preferences.getDatabaseName()}}%end signature
\begin{itemize}
\item{
{\bf  Description}

This method fetches the value of database name from the Eclipse preference store.
}
\item{{\bf  Returns} -- 
A String object representing the database name to which the application connects. 
}%end item
\end{itemize}
}%end item
\item{ 
\index{getDatabasePassword()}
\hypertarget{it.unisa.sesa.repominer.preferences.Preferences.getDatabasePassword()}{{\bf  getDatabasePassword}\\}
\texttt{public static java.lang.String\ {\bf  getDatabasePassword}()
\label{it.unisa.sesa.repominer.preferences.Preferences.getDatabasePassword()}}%end signature
\begin{itemize}
\item{
{\bf  Description}

This method fetches the value of database user's password from the Eclipse preference store.
}
\item{{\bf  Returns} -- 
A String object representing the database users's password with which the application logs into DBMS. 
}%end item
\end{itemize}
}%end item
\item{ 
\index{getDatabasePort()}
\hypertarget{it.unisa.sesa.repominer.preferences.Preferences.getDatabasePort()}{{\bf  getDatabasePort}\\}
\texttt{public static int\ {\bf  getDatabasePort}()
\label{it.unisa.sesa.repominer.preferences.Preferences.getDatabasePort()}}%end signature
\begin{itemize}
\item{
{\bf  Description}

This method fetches the value of database port from the Eclipse preference store.
}
\item{{\bf  Returns} -- 
A int value representing the port on which the DBMS is listening to. 
}%end item
\end{itemize}
}%end item
\item{ 
\index{getDatabaseUser()}
\hypertarget{it.unisa.sesa.repominer.preferences.Preferences.getDatabaseUser()}{{\bf  getDatabaseUser}\\}
\texttt{public static java.lang.String\ {\bf  getDatabaseUser}()
\label{it.unisa.sesa.repominer.preferences.Preferences.getDatabaseUser()}}%end signature
\begin{itemize}
\item{
{\bf  Description}

This method fetches the value of database username from the Eclipse preference store.
}
\item{{\bf  Returns} -- 
A String object representing the database username with which the application logs into DBMS. 
}%end item
\end{itemize}
}%end item
\item{ 
\index{getECCMModality()}
\hypertarget{it.unisa.sesa.repominer.preferences.Preferences.getECCMModality()}{{\bf  getECCMModality}\\}
\texttt{public static java.lang.String\ {\bf  getECCMModality}()
\label{it.unisa.sesa.repominer.preferences.Preferences.getECCMModality()}}%end signature
\begin{itemize}
\item{
{\bf  Description}

This method fetches the value of the Extended Code Change Model metric modality from the Eclipse preference store.
}
\item{{\bf  Returns} -- 
A String object representing the Extended Code Change Model metric modality. 
}%end item
\end{itemize}
}%end item
\item{ 
\index{getECCMModificationLimit()}
\hypertarget{it.unisa.sesa.repominer.preferences.Preferences.getECCMModificationLimit()}{{\bf  getECCMModificationLimit}\\}
\texttt{public static int\ {\bf  getECCMModificationLimit}() throws it.unisa.sesa.repominer.preferences.exceptions.IntegerPreferenceException
\label{it.unisa.sesa.repominer.preferences.Preferences.getECCMModificationLimit()}}%end signature
\begin{itemize}
\item{
{\bf  Description}

This method fetches the value of the modification limit value for the Extended Code Change Model metric. Value is fetched from the Eclipse preference store.
}
\item{{\bf  Returns} -- 
An int value representing value for the modification limit. 
}%end item
\item{{\bf  Throws}
  \begin{itemize}
   \item{\vskip -.6ex \texttt{it.unisa.sesa.repominer.preferences.exceptions.IntegerPreferenceException} -- This exception is thrown when the value is not a positive int number.}
  \end{itemize}
}%end item
\end{itemize}
}%end item
\item{ 
\index{getEpsParameter()}
\hypertarget{it.unisa.sesa.repominer.preferences.Preferences.getEpsParameter()}{{\bf  getEpsParameter}\\}
\texttt{public static int\ {\bf  getEpsParameter}() throws it.unisa.sesa.repominer.preferences.exceptions.IntegerPreferenceException
\label{it.unisa.sesa.repominer.preferences.Preferences.getEpsParameter()}}%end signature
\begin{itemize}
\item{
{\bf  Description}

This method fetches the value of radius parameter for DBSCAN algorithm from Eclipse preference store.
}
\item{{\bf  Returns} -- 
A int value representing the radius of the DBSCAN algorithm. 
}%end item
\item{{\bf  Throws}
  \begin{itemize}
   \item{\vskip -.6ex \texttt{it.unisa.sesa.repominer.preferences.exceptions.IntegerPreferenceException} -- Thrown when the value fetched is not a positive number.}
  \end{itemize}
}%end item
\end{itemize}
}%end item
\item{ 
\index{getMinPointsParameter()}
\hypertarget{it.unisa.sesa.repominer.preferences.Preferences.getMinPointsParameter()}{{\bf  getMinPointsParameter}\\}
\texttt{public static int\ {\bf  getMinPointsParameter}() throws it.unisa.sesa.repominer.preferences.exceptions.IntegerPreferenceException
\label{it.unisa.sesa.repominer.preferences.Preferences.getMinPointsParameter()}}%end signature
\begin{itemize}
\item{
{\bf  Description}

This method fetches the value of the minimum size per cluster relative to the DBSCAN algorithm from Eclipse preference store.
}
\item{{\bf  Returns} -- 
A int value representing the minimum size per cluster relative to the DBSCAN algorithm. 
}%end item
\item{{\bf  Throws}
  \begin{itemize}
   \item{\vskip -.6ex \texttt{it.unisa.sesa.repominer.preferences.exceptions.IntegerPreferenceException} -- Thrown when the value fetched is not a positive number.}
  \end{itemize}
}%end item
\end{itemize}
}%end item
\item{ 
\index{getPeriodEndingDate()}
\hypertarget{it.unisa.sesa.repominer.preferences.Preferences.getPeriodEndingDate()}{{\bf  getPeriodEndingDate}\\}
\texttt{public static java.util.Date\ {\bf  getPeriodEndingDate}() throws java.text.ParseException
\label{it.unisa.sesa.repominer.preferences.Preferences.getPeriodEndingDate()}}%end signature
\begin{itemize}
\item{
{\bf  Description}

This method fetches the value for the period ending date for the Basic Code Change Model metric. Value is fetched from the Eclipse preference store.
}
\item{{\bf  Returns} -- 
A Date object representing the ending date for the interval in which to calculate Basic Code Change Model metric. 
}%end item
\item{{\bf  Throws}
  \begin{itemize}
   \item{\vskip -.6ex \texttt{java.text.ParseException} -- This exception is thrown when the string fetched from Eclipse preference store is not parsable as a valid Date.}
  \end{itemize}
}%end item
\end{itemize}
}%end item
\item{ 
\index{getPeriodLength()}
\hypertarget{it.unisa.sesa.repominer.preferences.Preferences.getPeriodLength()}{{\bf  getPeriodLength}\\}
\texttt{public static int\ {\bf  getPeriodLength}() throws it.unisa.sesa.repominer.preferences.exceptions.IntegerPreferenceException, it.unisa.sesa.repominer.preferences.exceptions.PeriodLengthTooLong
\label{it.unisa.sesa.repominer.preferences.Preferences.getPeriodLength()}}%end signature
\begin{itemize}
\item{
{\bf  Description}

This method fetches the value of the period length from the Eclipse preference store.
}
\item{{\bf  Returns} -- 
An int value representing the value for the period length. 
}%end item
\item{{\bf  Throws}
  \begin{itemize}
   \item{\vskip -.6ex \texttt{it.unisa.sesa.repominer.preferences.exceptions.IntegerPreferenceException} -- Thrown when the value fetched is not positive number.}
   \item{\vskip -.6ex \texttt{it.unisa.sesa.repominer.preferences.exceptions.PeriodLengthTooLong} -- Thrown when the value exceeds a limit (4000 when considering weeks, 1000 when months, 80 years).}
  \end{itemize}
}%end item
\end{itemize}
}%end item
\item{ 
\index{getPeriodStartingDate()}
\hypertarget{it.unisa.sesa.repominer.preferences.Preferences.getPeriodStartingDate()}{{\bf  getPeriodStartingDate}\\}
\texttt{public static java.util.Date\ {\bf  getPeriodStartingDate}() throws java.text.ParseException
\label{it.unisa.sesa.repominer.preferences.Preferences.getPeriodStartingDate()}}%end signature
\begin{itemize}
\item{
{\bf  Description}

This method fetches the value for the period starting date for the Basic Code Change Model metric. Value is fetched from the Eclipse preference store.
}
\item{{\bf  Returns} -- 
A Date object representing the starting date for the interval in which to calculate Basic Code Change Model metric. 
}%end item
\item{{\bf  Throws}
  \begin{itemize}
   \item{\vskip -.6ex \texttt{java.text.ParseException} -- This exception is thrown when the string fetched from Eclipse preference store is not parsable as a valid Date.}
  \end{itemize}
}%end item
\end{itemize}
}%end item
\item{ 
\index{getPeriodType()}
\hypertarget{it.unisa.sesa.repominer.preferences.Preferences.getPeriodType()}{{\bf  getPeriodType}\\}
\texttt{public static java.lang.String\ {\bf  getPeriodType}()
\label{it.unisa.sesa.repominer.preferences.Preferences.getPeriodType()}}%end signature
\begin{itemize}
\item{
{\bf  Description}

This method fetches the value of the period type from Eclipse preference store.
}
\item{{\bf  Returns} -- 
A String object representing the period type relative to the Extended Code Change Model metric. 
}%end item
\end{itemize}
}%end item
\end{itemize}
}
}
}
\section{Package it.unisa.sesa.repominer.util.snippets}{
\label{it.unisa.sesa.repominer.util.snippets}\hypertarget{it.unisa.sesa.repominer.util.snippets}{}
\hskip -.05in
\hbox to \hsize{\textit{ Package Contents\hfil Page}}
\vskip .13in
\hbox{{\bf  Classes}}
\entityintro{Test}{it.unisa.sesa.repominer.util.snippets.Test}{}
\vskip .1in
\vskip .1in
\subsection{\label{it.unisa.sesa.repominer.util.snippets.Test}Class Test}{
\hypertarget{it.unisa.sesa.repominer.util.snippets.Test}{}\vskip .1in 
\subsubsection{Declaration}{
\small public class Test
\\ {\bf  extends} java.lang.Object
\refdefined{java.lang.Object}}
\subsubsection{Constructor summary}{
\begin{verse}
\hyperlink{it.unisa.sesa.repominer.util.snippets.Test()}{{\bf Test()}} \\
\end{verse}
}
\subsubsection{Method summary}{
\begin{verse}
\hyperlink{it.unisa.sesa.repominer.util.snippets.Test.main(java.lang.String[])}{{\bf main(String\lbrack \rbrack )}} \\
\end{verse}
}
\subsubsection{Constructors}{
\vskip -2em
\begin{itemize}
\item{ 
\index{Test()}
\hypertarget{it.unisa.sesa.repominer.util.snippets.Test()}{{\bf  Test}\\}
\texttt{public\ {\bf  Test}()
\label{it.unisa.sesa.repominer.util.snippets.Test()}}%end signature
}%end item
\end{itemize}
}
\subsubsection{Methods}{
\vskip -2em
\begin{itemize}
\item{ 
\index{main(String\lbrack \rbrack )}
\hypertarget{it.unisa.sesa.repominer.util.snippets.Test.main(java.lang.String[])}{{\bf  main}\\}
\texttt{public static void\ {\bf  main}(\texttt{java.lang.String\lbrack \rbrack } {\bf  args})
\label{it.unisa.sesa.repominer.util.snippets.Test.main(java.lang.String[])}}%end signature
}%end item
\end{itemize}
}
}
}
\section{Package it.unisa.sesa.repominer.util}{
\label{it.unisa.sesa.repominer.util}\hypertarget{it.unisa.sesa.repominer.util}{}
\hskip -.05in
\hbox to \hsize{\textit{ Package Contents\hfil Page}}
\vskip .13in
\hbox{{\bf  Classes}}
\entityintro{Utils}{it.unisa.sesa.repominer.util.Utils}{}
\vskip .1in
\vskip .1in
\subsection{\label{it.unisa.sesa.repominer.util.Utils}Class Utils}{
\hypertarget{it.unisa.sesa.repominer.util.Utils}{}\vskip .1in 
\subsubsection{Declaration}{
\small public class Utils
\\ {\bf  extends} java.lang.Object
\refdefined{java.lang.Object}}
\subsubsection{Constructor summary}{
\begin{verse}
\hyperlink{it.unisa.sesa.repominer.util.Utils()}{{\bf Utils()}} \\
\end{verse}
}
\subsubsection{Method summary}{
\begin{verse}
\hyperlink{it.unisa.sesa.repominer.util.Utils.dateToCalendar(java.util.Date)}{{\bf dateToCalendar(Date)}} This method convert an instance of Date into an instance of Calendar\\
\hyperlink{it.unisa.sesa.repominer.util.Utils.daysBetween(java.util.Calendar, java.util.Calendar)}{{\bf daysBetween(Calendar, Calendar)}} This method calculates difference in days between two dates\\
\hyperlink{it.unisa.sesa.repominer.util.Utils.log2(double)}{{\bf log2(double)}} This function calculate the base-2 logarithm of a float number\\
\hyperlink{it.unisa.sesa.repominer.util.Utils.msgIsBugFixing(java.lang.String)}{{\bf msgIsBugFixing(String)}} Checks whether or not a message is likely to be a bug-fixing one.\\
\hyperlink{it.unisa.sesa.repominer.util.Utils.msgIsRefactoring(java.lang.String)}{{\bf msgIsRefactoring(String)}} Checks whether or not a message is likely to be a refactoring one.\\
\hyperlink{it.unisa.sesa.repominer.util.Utils.stringToDate(java.lang.String)}{{\bf stringToDate(String)}} This method convert a String into a Date\\
\end{verse}
}
\subsubsection{Constructors}{
\vskip -2em
\begin{itemize}
\item{ 
\index{Utils()}
\hypertarget{it.unisa.sesa.repominer.util.Utils()}{{\bf  Utils}\\}
\texttt{public\ {\bf  Utils}()
\label{it.unisa.sesa.repominer.util.Utils()}}%end signature
}%end item
\end{itemize}
}
\subsubsection{Methods}{
\vskip -2em
\begin{itemize}
\item{ 
\index{dateToCalendar(Date)}
\hypertarget{it.unisa.sesa.repominer.util.Utils.dateToCalendar(java.util.Date)}{{\bf  dateToCalendar}\\}
\texttt{public static java.util.Calendar\ {\bf  dateToCalendar}(\texttt{java.util.Date} {\bf  pDate})
\label{it.unisa.sesa.repominer.util.Utils.dateToCalendar(java.util.Date)}}%end signature
\begin{itemize}
\item{
{\bf  Description}

This method convert an instance of Date into an instance of Calendar
}
\item{
{\bf  Parameters}
  \begin{itemize}
   \item{
\texttt{pDate} -- }
  \end{itemize}
}%end item
\item{{\bf  Returns} -- 
A Calendar object 
}%end item
\end{itemize}
}%end item
\item{ 
\index{daysBetween(Calendar, Calendar)}
\hypertarget{it.unisa.sesa.repominer.util.Utils.daysBetween(java.util.Calendar, java.util.Calendar)}{{\bf  daysBetween}\\}
\texttt{public static int\ {\bf  daysBetween}(\texttt{java.util.Calendar} {\bf  startDate},
\texttt{java.util.Calendar} {\bf  endDate})
\label{it.unisa.sesa.repominer.util.Utils.daysBetween(java.util.Calendar, java.util.Calendar)}}%end signature
\begin{itemize}
\item{
{\bf  Description}

This method calculates difference in days between two dates
}
\item{
{\bf  Parameters}
  \begin{itemize}
   \item{
\texttt{startDate} -- start date}
   \item{
\texttt{endDate} -- end date}
  \end{itemize}
}%end item
\item{{\bf  Returns} -- 
days from start date to end date 
}%end item
\end{itemize}
}%end item
\item{ 
\index{log2(double)}
\hypertarget{it.unisa.sesa.repominer.util.Utils.log2(double)}{{\bf  log2}\\}
\texttt{public static double\ {\bf  log2}(\texttt{double} {\bf  pValue})
\label{it.unisa.sesa.repominer.util.Utils.log2(double)}}%end signature
\begin{itemize}
\item{
{\bf  Description}

This function calculate the base-2 logarithm of a float number
}
\item{
{\bf  Parameters}
  \begin{itemize}
   \item{
\texttt{pValue} -- }
  \end{itemize}
}%end item
\item{{\bf  Returns} -- 
Base-2 logarithm value 
}%end item
\end{itemize}
}%end item
\item{ 
\index{msgIsBugFixing(String)}
\hypertarget{it.unisa.sesa.repominer.util.Utils.msgIsBugFixing(java.lang.String)}{{\bf  msgIsBugFixing}\\}
\texttt{public static boolean\ {\bf  msgIsBugFixing}(\texttt{java.lang.String} {\bf  msg})
\label{it.unisa.sesa.repominer.util.Utils.msgIsBugFixing(java.lang.String)}}%end signature
\begin{itemize}
\item{
{\bf  Description}

Checks whether or not a message is likely to be a bug-fixing one.
}
\item{
{\bf  Parameters}
  \begin{itemize}
   \item{
\texttt{msg} -- The message to check.}
  \end{itemize}
}%end item
\item{{\bf  Returns} -- 
true if message is bug-fixing, false otherwise. 
}%end item
\end{itemize}
}%end item
\item{ 
\index{msgIsRefactoring(String)}
\hypertarget{it.unisa.sesa.repominer.util.Utils.msgIsRefactoring(java.lang.String)}{{\bf  msgIsRefactoring}\\}
\texttt{public static boolean\ {\bf  msgIsRefactoring}(\texttt{java.lang.String} {\bf  msg})
\label{it.unisa.sesa.repominer.util.Utils.msgIsRefactoring(java.lang.String)}}%end signature
\begin{itemize}
\item{
{\bf  Description}

Checks whether or not a message is likely to be a refactoring one.
}
\item{
{\bf  Parameters}
  \begin{itemize}
   \item{
\texttt{msg} -- The message to check.}
  \end{itemize}
}%end item
\item{{\bf  Returns} -- 
true if message is refactoring, false otherwise. 
}%end item
\end{itemize}
}%end item
\item{ 
\index{stringToDate(String)}
\hypertarget{it.unisa.sesa.repominer.util.Utils.stringToDate(java.lang.String)}{{\bf  stringToDate}\\}
\texttt{public static java.util.Date\ {\bf  stringToDate}(\texttt{java.lang.String} {\bf  pString}) throws java.text.ParseException
\label{it.unisa.sesa.repominer.util.Utils.stringToDate(java.lang.String)}}%end signature
\begin{itemize}
\item{
{\bf  Description}

This method convert a String into a Date
}
\item{
{\bf  Parameters}
  \begin{itemize}
   \item{
\texttt{pString} -- }
  \end{itemize}
}%end item
\item{{\bf  Returns} -- 
A Date Object 
}%end item
\end{itemize}
}%end item
\end{itemize}
}
}
}

}
	\chapter{Glossario}
% prova longtable
\begin{center}
\begin{longtable}{lp{.6\columnwidth}}
\hline 
\rowcolor[gray]{.80}
\textbf{Termine} & \textbf{Descrizione} \tabularnewline
\hline 
SIS & 
Starting Impact Set: artefatti che, dopo una analisi delle specifiche della change request, si presume di dover intaccare andando a implementare il cambiamento richiesto
\tabularnewline
\hline
CIS & 
Candidate Impact Set: ulteriore livello di analisi sugli artefatti individuati dal SIS
\tabularnewline
\hline
Recall & 
Metrica per la valutazione del processo di impact analysis che misura la percentuale degli impatti reali inclusi in CIS
\tabularnewline
\hline
Precision & 
Metrica per la valutazione del processo di impact analysis che misura la percentuale degli impatti candidati che sono impatti reali
\tabularnewline
\hline
DIS & 
Discovered Impact Set: sottostima dell'impatto; si tratta degli artefatti che fanno parte dell'AID ma che non sono stati inclusi nel CIS al momento dell'analisi
\tabularnewline
\hline
AID & 
Actual Impact Set: insieme degli artefatti realmente modificati
\tabularnewline
\hline
FPIS & 
False Positive Impact Set: sovrastima dell'impatto; si tratta degli artefatti che sono stati inclusi nel CIS ma che non sono stati impattati
\tabularnewline
\hline
\end{longtable}
\end{center}





\end{document}
